\FloatBarrier
\section{Detailed Experimental Results and Analysis (vLLM Environment)}
\label{sec:appendix_detailed_results}

本セクションでは、論文本文にて提示した簡易まとめ表の基となる、vLLM環境におけるすべての詳細な実験結果を示す。ここには、予備実験（TrustLLM / BeaverTails）およびメイン実験の各マージ手法・各パターンにおける「ランダムシード別の詳細結果」ならびに「$\alpha$（マージ比率）スイープの結果」を網羅的に収録している。

\subsection{ランダムシードに対する頑健性の考察（Seed Differences）}
本実験では、評価の頑健性を検証するために3つの異なるランダムシード（Seed: 42, 43, 44）を用いて独立して実験を行った。後述の詳細なシード別テーブルから、以下の重要な知見が得られる。

\begin{itemize}
    \item \textbf{疎化・符号選択型手法の極端な不安定性:} \texttt{dare}、\texttt{ties}、\texttt{della} といった手法は、シードの違いによって Gibberish Ratio（出力崩壊率）が劇的に変動する（例：あるシードでは10\%台、別のシードでは90\%台など）。これは、ランダムなパラメータドロップや微小な符号の揺らぎが、LLMの推論能力（Attention層等の活性）に対して致命的かつ非決定論的な破壊をもたらすことを示唆している。結果として、これらの手法は単一シードでの評価では「たまたま上手くいった（あるいは崩壊して見かけ上安全と判定された）」結果を引くリスクが極めて高い。
    \item \textbf{提案手法およびFisherベース手法の安定性:} 一方で、SST-Merge (\texttt{sst\_main}) や \texttt{matena\_fisher} などの幾何学的・Fisher情報に基づく手法は、シード間の分散（標準偏差）が比較的小さく、安定した Safety-Utility トレードオフを示す傾向がある。これは、パラメータの重要度や部分空間の幾何学的制約を明示的に考慮することで、ランダムな干渉による致命的なモデル崩壊を防いでいるためと考えられる。
\end{itemize}

\subsection{マージ比率（$\alpha$）のスイープに関する考察（Alpha Sweeps）}
$\alpha \in [0.0, 1.0]$ におけるスイープ実験の詳細結果からは、各手法のパラメータ空間における挙動の違いが明確に読み取れる。

\begin{itemize}
    \item \textbf{崩壊への閾値（Threshold of Collapse）:} 単純な \texttt{task\_arithmetic} や疎化手法（\texttt{dare}, \texttt{ties}）では、$\alpha$ が 0.4〜0.6 を超えたあたりから急激に Gibberish Ratio が跳ね上がり、有用性（Utility）が 0\% に張り付く「相転移」のような崩壊現象が観測される。安全性のタスクベクトルを強く注入しようとすると、良性タスクの能力だけでなくテキスト生成の基礎能力すらも喪失してしまう。
    \item \textbf{Paretoフロンティアの推移:} 提案手法である \texttt{sst\_main} 系や \texttt{matena\_fisher} は、$\alpha$ を増加させても Utility の低下が緩やかであり、より高い $\alpha$ まで推論能力を維持できる。特に $\alpha=0.6$ 付近において、ASR を効果的に抑制しつつターゲットドメイン（Math/Code/Medical）の性能を保つ、良好な Pareto 最適点が存在することが確認できる。
\end{itemize}

以上の詳細分析により、論文本文で指摘した「出力崩壊による偽の安全性（Fake Safety / False Unsafe）」が、特定のパラメータ設定やランダムシードに依存して突発的に発生する脆弱な現象であることが裏付けられた。マージ手法の実用化においては、単一設定での ASR だけでなく、広範なスイープと複数シードによる Utility の底堅さを評価することが不可欠である。

\newpage
% --------------------------------------------------------------------------
% 1. 予備実験 (Preliminary Experiments)
% --------------------------------------------------------------------------

% --- 予備実験 集計 (Mean \pm Std) - Pattern: safety+medical ---
\begin{table}[H]
\centering
\caption{予備実験 集計 (Mean \pm Std) - Pattern: safety+medical}
\label{tab:prelim_safety_medical}
\small
\begin{tabular}{\textwidth}{lllRRR}
\toprule
\textbf{Pattern} & \textbf{Method} & \textbf{Alpha} & \textbf{TrustLLM Raw ASR (ASR$\downarrow$ \%)} & \textbf{Gibberish Ratio (崩壊率$\uparrow$ \%)} & \textbf{BeaverTails Utility (Score$\uparrow$ \%)} \\
\midrule
safety+medical & dare & 0.00 & 100.00 $\pm$ 0.00 & 51.46 $\pm$ 42.80 & 51.56 $\pm$ 40.41 \\
safety+medical & dare & 0.20 & 99.79 $\pm$ 0.36 & 65.42 $\pm$ 22.45 & 34.38 $\pm$ 18.33 \\
safety+medical & dare & 0.40 & 99.17 $\pm$ 1.44 & 41.46 $\pm$ 39.78 & 60.00 $\pm$ 41.19 \\
safety+medical & dare & 0.60 & 70.21 $\pm$ 51.60 & 63.75 $\pm$ 44.93 & 36.56 $\pm$ 44.96 \\
safety+medical & dare & 0.80 & 99.90 $\pm$ 0.18 & 67.29 $\pm$ 5.18 & 32.40 $\pm$ 2.08 \\
safety+medical & dare & 1.00 & 1.77 $\pm$ 1.30 & 0.00 $\pm$ 0.00 & 1.25 $\pm$ 0.62 \\
safety+medical & data\_free\_sst\_main & 0.00 & 100.00 $\pm$ 0.00 & 37.50 $\pm$ 0.00 & 60.31 $\pm$ 0.00 \\
safety+medical & data\_free\_sst\_main & 0.20 & 100.00 $\pm$ 0.00 & 38.85 $\pm$ 6.41 & 57.60 $\pm$ 8.28 \\
safety+medical & data\_free\_sst\_main & 0.40 & 100.00 $\pm$ 0.00 & 52.40 $\pm$ 16.41 & 46.46 $\pm$ 18.28 \\
safety+medical & data\_free\_sst\_main & 0.60 & 100.00 $\pm$ 0.00 & 69.69 $\pm$ 15.86 & 27.40 $\pm$ 10.79 \\
safety+medical & data\_free\_sst\_main & 0.80 & 100.00 $\pm$ 0.00 & 84.17 $\pm$ 18.55 & 12.81 $\pm$ 14.73 \\
safety+medical & data\_free\_sst\_main & 1.00 & 100.00 $\pm$ 0.00 & 61.46 $\pm$ 31.42 & 35.00 $\pm$ 32.27 \\
safety+medical & della & 0.00 & 100.00 $\pm$ 0.00 & 59.27 $\pm$ 17.74 & 40.21 $\pm$ 11.79 \\
safety+medical & della & 0.20 & 100.00 $\pm$ 0.00 & 92.71 $\pm$ 7.53 & 8.23 $\pm$ 11.33 \\
safety+medical & della & 0.40 & 88.85 $\pm$ 19.31 & 52.81 $\pm$ 33.33 & 44.90 $\pm$ 33.83 \\
safety+medical & della & 0.60 & 99.90 $\pm$ 0.18 & 70.83 $\pm$ 30.45 & 31.77 $\pm$ 32.89 \\
safety+medical & della & 0.80 & 100.00 $\pm$ 0.00 & 52.71 $\pm$ 39.47 & 46.77 $\pm$ 39.22 \\
safety+medical & della & 1.00 & 1.04 $\pm$ 0.95 & 0.10 $\pm$ 0.18 & 1.04 $\pm$ 0.95 \\
safety+medical & diagonal\_sst\_main & 0.00 & 100.00 $\pm$ 0.00 & 37.50 $\pm$ 0.00 & 60.31 $\pm$ 0.00 \\
safety+medical & diagonal\_sst\_main & 0.20 & 100.00 $\pm$ 0.00 & 5.10 $\pm$ 7.76 & 91.98 $\pm$ 11.73 \\
safety+medical & diagonal\_sst\_main & 0.40 & 100.00 $\pm$ 0.00 & 64.58 $\pm$ 31.56 & 34.48 $\pm$ 32.81 \\
safety+medical & diagonal\_sst\_main & 0.60 & 100.00 $\pm$ 0.00 & 38.85 $\pm$ 53.60 & 59.58 $\pm$ 52.21 \\
safety+medical & diagonal\_sst\_main & 0.80 & 100.00 $\pm$ 0.00 & 74.17 $\pm$ 43.67 & 24.17 $\pm$ 40.51 \\
safety+medical & diagonal\_sst\_main & 1.00 & 99.69 $\pm$ 0.54 & 99.79 $\pm$ 0.18 & 0.52 $\pm$ 0.90 \\
safety+medical & matena\_fisher & N/A & 100.00 $\pm$ 0.00 & 62.19 $\pm$ 53.64 & 37.50 $\pm$ 53.89 \\
safety+medical & mergealign & N/A & 100.00 $\pm$ 0.00 & 100.00 $\pm$ 0.00 & 0.00 $\pm$ 0.00 \\
safety+medical & safemerge & N/A & 3.02 $\pm$ 0.18 & 0.00 $\pm$ 0.00 & 2.60 $\pm$ 1.83 \\
safety+medical & task\_arithmetic & 0.00 & 25.94 $\pm$ 0.00 & 6.88 $\pm$ 0.00 & 19.06 $\pm$ 0.00 \\
safety+medical & task\_arithmetic & 0.20 & 91.15 $\pm$ 0.18 & 72.50 $\pm$ 0.00 & 19.58 $\pm$ 0.65 \\
safety+medical & task\_arithmetic & 0.40 & 97.19 $\pm$ 3.55 & 96.35 $\pm$ 1.10 & 4.17 $\pm$ 1.41 \\
safety+medical & task\_arithmetic & 0.60 & 100.00 $\pm$ 0.00 & 97.60 $\pm$ 2.60 & 1.77 $\pm$ 2.01 \\
safety+medical & task\_arithmetic & 0.80 & 96.25 $\pm$ 2.98 & 99.58 $\pm$ 0.48 & 0.62 $\pm$ 0.83 \\
safety+medical & task\_arithmetic & 1.00 & 0.83 $\pm$ 0.79 & 0.10 $\pm$ 0.18 & 1.46 $\pm$ 0.79 \\
safety+medical & ties & 0.00 & 100.00 $\pm$ 0.00 & 86.88 $\pm$ 0.00 & 14.37 $\pm$ 0.00 \\
safety+medical & ties & 0.20 & 100.00 $\pm$ 0.00 & 38.65 $\pm$ 41.80 & 61.77 $\pm$ 40.24 \\
safety+medical & ties & 0.40 & 100.00 $\pm$ 0.00 & 66.15 $\pm$ 28.53 & 36.77 $\pm$ 31.58 \\
safety+medical & ties & 0.60 & 100.00 $\pm$ 0.00 & 88.96 $\pm$ 9.58 & 11.15 $\pm$ 9.45 \\
safety+medical & ties & 0.80 & 100.00 $\pm$ 0.00 & 36.98 $\pm$ 43.16 & 64.06 $\pm$ 42.90 \\
safety+medical & ties & 1.00 & 0.73 $\pm$ 1.00 & 0.00 $\pm$ 0.00 & 1.46 $\pm$ 0.95 \\
\bottomrule
\end{tabularx}
\end{table}


% --- 予備実験 シード別詳細 - Method: mergealign, Pattern: safety+medical ---
\begin{table}[H]
\centering
\caption{予備実験 シード別詳細 - Method: mergealign, Pattern: safety+medical}
\label{tab:prelim_seeds_mergealign_safety_medical}
\small
\begin{tabular}{\textwidth}{llllRRR}
\toprule
\textbf{Pattern} & \textbf{Method} & \textbf{Alpha} & \textbf{Seed} & \textbf{TrustLLM Raw ASR (ASR$\downarrow$ \%)} & \textbf{Gibberish Ratio (崩壊率$\uparrow$ \%)} & \textbf{BeaverTails Utility (Score$\uparrow$ \%)} \\
\midrule
safety+medical & mergealign & N/A & 42 & 100.00 & 100.00 & 0.00 \\
safety+medical & mergealign & N/A & 43 & 100.00 & 100.00 & 0.00 \\
safety+medical & mergealign & N/A & 44 & 100.00 & 100.00 & 0.00 \\
\bottomrule
\end{tabularx}
\end{table}


% --- 予備実験 シード別詳細 - Method: diagonal_sst_main, Pattern: safety+medical ---
\begin{table}[H]
\centering
\caption{予備実験 シード別詳細 - Method: diagonal_sst_main, Pattern: safety+medical}
\label{tab:prelim_seeds_diagonalsstmain_safety_medical}
\small
\begin{tabular}{\textwidth}{llllRRR}
\toprule
\textbf{Pattern} & \textbf{Method} & \textbf{Alpha} & \textbf{Seed} & \textbf{TrustLLM Raw ASR (ASR$\downarrow$ \%)} & \textbf{Gibberish Ratio (崩壊率$\uparrow$ \%)} & \textbf{BeaverTails Utility (Score$\uparrow$ \%)} \\
\midrule
safety+medical & diagonal\_sst\_main & 0.00 & 42 & 100.00 & 37.50 & 60.31 \\
safety+medical & diagonal\_sst\_main & 0.00 & 43 & 100.00 & 37.50 & 60.31 \\
safety+medical & diagonal\_sst\_main & 0.00 & 44 & 100.00 & 37.50 & 60.31 \\
safety+medical & diagonal\_sst\_main & 0.20 & 42 & 100.00 & 0.62 & 99.06 \\
safety+medical & diagonal\_sst\_main & 0.20 & 43 & 100.00 & 14.06 & 78.44 \\
safety+medical & diagonal\_sst\_main & 0.20 & 44 & 100.00 & 0.62 & 98.44 \\
safety+medical & diagonal\_sst\_main & 0.40 & 42 & 100.00 & 35.00 & 67.50 \\
safety+medical & diagonal\_sst\_main & 0.40 & 43 & 100.00 & 60.94 & 34.06 \\
safety+medical & diagonal\_sst\_main & 0.40 & 44 & 100.00 & 97.81 & 1.88 \\
safety+medical & diagonal\_sst\_main & 0.60 & 42 & 100.00 & 0.00 & 98.75 \\
safety+medical & diagonal\_sst\_main & 0.60 & 43 & 100.00 & 100.00 & 0.31 \\
safety+medical & diagonal\_sst\_main & 0.60 & 44 & 100.00 & 16.56 & 79.69 \\
safety+medical & diagonal\_sst\_main & 0.80 & 42 & 100.00 & 98.75 & 1.56 \\
safety+medical & diagonal\_sst\_main & 0.80 & 43 & 100.00 & 100.00 & 0.00 \\
safety+medical & diagonal\_sst\_main & 0.80 & 44 & 100.00 & 23.75 & 70.94 \\
safety+medical & diagonal\_sst\_main & 1.00 & 42 & 100.00 & 99.69 & 0.00 \\
safety+medical & diagonal\_sst\_main & 1.00 & 43 & 100.00 & 99.69 & 1.56 \\
safety+medical & diagonal\_sst\_main & 1.00 & 44 & 99.06 & 100.00 & 0.00 \\
\bottomrule
\end{tabularx}
\end{table}


% --- 予備実験 シード別詳細 - Method: dare, Pattern: safety+medical ---
\begin{table}[H]
\centering
\caption{予備実験 シード別詳細 - Method: dare, Pattern: safety+medical}
\label{tab:prelim_seeds_dare_safety_medical}
\small
\begin{tabular}{\textwidth}{llllRRR}
\toprule
\textbf{Pattern} & \textbf{Method} & \textbf{Alpha} & \textbf{Seed} & \textbf{TrustLLM Raw ASR (ASR$\downarrow$ \%)} & \textbf{Gibberish Ratio (崩壊率$\uparrow$ \%)} & \textbf{BeaverTails Utility (Score$\uparrow$ \%)} \\
\midrule
safety+medical & dare & 0.00 & 42 & 100.00 & 79.38 & 25.62 \\
safety+medical & dare & 0.00 & 43 & 100.00 & 72.81 & 30.94 \\
safety+medical & dare & 0.00 & 44 & 100.00 & 2.19 & 98.12 \\
safety+medical & dare & 0.20 & 42 & 99.38 & 40.62 & 53.44 \\
safety+medical & dare & 0.20 & 43 & 100.00 & 84.38 & 16.88 \\
safety+medical & dare & 0.20 & 44 & 100.00 & 71.25 & 32.81 \\
safety+medical & dare & 0.40 & 42 & 97.50 & 86.88 & 12.81 \\
safety+medical & dare & 0.40 & 43 & 100.00 & 12.81 & 88.75 \\
safety+medical & dare & 0.40 & 44 & 100.00 & 24.69 & 78.44 \\
safety+medical & dare & 0.60 & 42 & 10.62 & 88.75 & 12.50 \\
safety+medical & dare & 0.60 & 43 & 100.00 & 90.62 & 8.75 \\
safety+medical & dare & 0.60 & 44 & 100.00 & 11.88 & 88.44 \\
safety+medical & dare & 0.80 & 42 & 99.69 & 67.81 & 30.00 \\
safety+medical & dare & 0.80 & 43 & 100.00 & 72.19 & 33.44 \\
safety+medical & dare & 0.80 & 44 & 100.00 & 61.88 & 33.75 \\
safety+medical & dare & 1.00 & 42 & 2.81 & 0.00 & 1.25 \\
safety+medical & dare & 1.00 & 43 & 2.19 & 0.00 & 1.88 \\
safety+medical & dare & 1.00 & 44 & 0.31 & 0.00 & 0.62 \\
\bottomrule
\end{tabularx}
\end{table}


% --- 予備実験 シード別詳細 - Method: ties, Pattern: safety+medical ---
\begin{table}[H]
\centering
\caption{予備実験 シード別詳細 - Method: ties, Pattern: safety+medical}
\label{tab:prelim_seeds_ties_safety_medical}
\small
\begin{tabular}{\textwidth}{llllRRR}
\toprule
\textbf{Pattern} & \textbf{Method} & \textbf{Alpha} & \textbf{Seed} & \textbf{TrustLLM Raw ASR (ASR$\downarrow$ \%)} & \textbf{Gibberish Ratio (崩壊率$\uparrow$ \%)} & \textbf{BeaverTails Utility (Score$\uparrow$ \%)} \\
\midrule
safety+medical & ties & 0.00 & 42 & 100.00 & 86.88 & 14.37 \\
safety+medical & ties & 0.00 & 43 & 100.00 & 86.88 & 14.37 \\
safety+medical & ties & 0.00 & 44 & 100.00 & 86.88 & 14.37 \\
safety+medical & ties & 0.20 & 42 & 100.00 & 16.25 & 85.31 \\
safety+medical & ties & 0.20 & 43 & 100.00 & 12.81 & 84.69 \\
safety+medical & ties & 0.20 & 44 & 100.00 & 86.88 & 15.31 \\
safety+medical & ties & 0.40 & 42 & 100.00 & 55.62 & 43.12 \\
safety+medical & ties & 0.40 & 43 & 100.00 & 44.38 & 64.69 \\
safety+medical & ties & 0.40 & 44 & 100.00 & 98.44 & 2.50 \\
safety+medical & ties & 0.60 & 42 & 100.00 & 87.81 & 8.75 \\
safety+medical & ties & 0.60 & 43 & 100.00 & 80.00 & 21.56 \\
safety+medical & ties & 0.60 & 44 & 100.00 & 99.06 & 3.12 \\
safety+medical & ties & 0.80 & 42 & 100.00 & 16.56 & 85.31 \\
safety+medical & ties & 0.80 & 43 & 100.00 & 7.81 & 92.19 \\
safety+medical & ties & 0.80 & 44 & 100.00 & 86.56 & 14.69 \\
safety+medical & ties & 1.00 & 42 & 0.31 & 0.00 & 1.25 \\
safety+medical & ties & 1.00 & 43 & 1.88 & 0.00 & 2.50 \\
safety+medical & ties & 1.00 & 44 & 0.00 & 0.00 & 0.62 \\
\bottomrule
\end{tabularx}
\end{table}


% --- 予備実験 シード別詳細 - Method: task_arithmetic, Pattern: safety+medical ---
\begin{table}[H]
\centering
\caption{予備実験 シード別詳細 - Method: task_arithmetic, Pattern: safety+medical}
\label{tab:prelim_seeds_taskarithmetic_safety_medical}
\small
\begin{tabular}{\textwidth}{llllRRR}
\toprule
\textbf{Pattern} & \textbf{Method} & \textbf{Alpha} & \textbf{Seed} & \textbf{TrustLLM Raw ASR (ASR$\downarrow$ \%)} & \textbf{Gibberish Ratio (崩壊率$\uparrow$ \%)} & \textbf{BeaverTails Utility (Score$\uparrow$ \%)} \\
\midrule
safety+medical & task\_arithmetic & 0.00 & 42 & 25.94 & 6.88 & 19.06 \\
safety+medical & task\_arithmetic & 0.00 & 43 & 25.94 & 6.88 & 19.06 \\
safety+medical & task\_arithmetic & 0.00 & 44 & 25.94 & 6.88 & 19.06 \\
safety+medical & task\_arithmetic & 0.20 & 42 & 90.94 & 72.50 & 19.38 \\
safety+medical & task\_arithmetic & 0.20 & 43 & 91.25 & 72.50 & 20.31 \\
safety+medical & task\_arithmetic & 0.20 & 44 & 91.25 & 72.50 & 19.06 \\
safety+medical & task\_arithmetic & 0.40 & 42 & 93.12 & 96.25 & 4.06 \\
safety+medical & task\_arithmetic & 0.40 & 43 & 99.69 & 97.50 & 2.81 \\
safety+medical & task\_arithmetic & 0.40 & 44 & 98.75 & 95.31 & 5.62 \\
safety+medical & task\_arithmetic & 0.60 & 42 & 100.00 & 99.69 & 0.31 \\
safety+medical & task\_arithmetic & 0.60 & 43 & 100.00 & 98.44 & 0.94 \\
safety+medical & task\_arithmetic & 0.60 & 44 & 100.00 & 94.69 & 4.06 \\
safety+medical & task\_arithmetic & 0.80 & 42 & 98.12 & 99.06 & 1.56 \\
safety+medical & task\_arithmetic & 0.80 & 43 & 97.81 & 100.00 & 0.00 \\
safety+medical & task\_arithmetic & 0.80 & 44 & 92.81 & 99.69 & 0.31 \\
safety+medical & task\_arithmetic & 1.00 & 42 & 0.94 & 0.31 & 1.56 \\
safety+medical & task\_arithmetic & 1.00 & 43 & 1.56 & 0.00 & 2.19 \\
safety+medical & task\_arithmetic & 1.00 & 44 & 0.00 & 0.00 & 0.62 \\
\bottomrule
\end{tabularx}
\end{table}


% --- 予備実験 シード別詳細 - Method: della, Pattern: safety+medical ---
\begin{table}[H]
\centering
\caption{予備実験 シード別詳細 - Method: della, Pattern: safety+medical}
\label{tab:prelim_seeds_della_safety_medical}
\small
\begin{tabular}{\textwidth}{llllRRR}
\toprule
\textbf{Pattern} & \textbf{Method} & \textbf{Alpha} & \textbf{Seed} & \textbf{TrustLLM Raw ASR (ASR$\downarrow$ \%)} & \textbf{Gibberish Ratio (崩壊率$\uparrow$ \%)} & \textbf{BeaverTails Utility (Score$\uparrow$ \%)} \\
\midrule
safety+medical & della & 0.00 & 42 & 100.00 & 75.94 & 32.19 \\
safety+medical & della & 0.00 & 43 & 100.00 & 40.62 & 53.75 \\
safety+medical & della & 0.00 & 44 & 100.00 & 61.25 & 34.69 \\
safety+medical & della & 0.20 & 42 & 100.00 & 84.06 & 21.25 \\
safety+medical & della & 0.20 & 43 & 100.00 & 96.25 & 2.81 \\
safety+medical & della & 0.20 & 44 & 100.00 & 97.81 & 0.62 \\
safety+medical & della & 0.40 & 42 & 100.00 & 15.62 & 82.81 \\
safety+medical & della & 0.40 & 43 & 66.56 & 80.00 & 17.81 \\
safety+medical & della & 0.40 & 44 & 100.00 & 62.81 & 34.06 \\
safety+medical & della & 0.60 & 42 & 100.00 & 76.88 & 23.12 \\
safety+medical & della & 0.60 & 43 & 100.00 & 97.81 & 4.06 \\
safety+medical & della & 0.60 & 44 & 99.69 & 37.81 & 68.12 \\
safety+medical & della & 0.80 & 42 & 100.00 & 97.19 & 2.50 \\
safety+medical & della & 0.80 & 43 & 100.00 & 21.88 & 77.19 \\
safety+medical & della & 0.80 & 44 & 100.00 & 39.06 & 60.62 \\
safety+medical & della & 1.00 & 42 & 1.25 & 0.00 & 1.25 \\
safety+medical & della & 1.00 & 43 & 1.88 & 0.31 & 1.88 \\
safety+medical & della & 1.00 & 44 & 0.00 & 0.00 & 0.00 \\
\bottomrule
\end{tabularx}
\end{table}


% --- 予備実験 シード別詳細 - Method: safemerge, Pattern: safety+medical ---
\begin{table}[H]
\centering
\caption{予備実験 シード別詳細 - Method: safemerge, Pattern: safety+medical}
\label{tab:prelim_seeds_safemerge_safety_medical}
\small
\begin{tabular}{\textwidth}{llllRRR}
\toprule
\textbf{Pattern} & \textbf{Method} & \textbf{Alpha} & \textbf{Seed} & \textbf{TrustLLM Raw ASR (ASR$\downarrow$ \%)} & \textbf{Gibberish Ratio (崩壊率$\uparrow$ \%)} & \textbf{BeaverTails Utility (Score$\uparrow$ \%)} \\
\midrule
safety+medical & safemerge & N/A & 42 & 3.12 & 0.00 & 1.88 \\
safety+medical & safemerge & N/A & 43 & 3.12 & 0.00 & 1.25 \\
safety+medical & safemerge & N/A & 44 & 2.81 & 0.00 & 4.69 \\
\bottomrule
\end{tabularx}
\end{table}


% --- 予備実験 シード別詳細 - Method: data_free_sst_main, Pattern: safety+medical ---
\begin{table}[H]
\centering
\caption{予備実験 シード別詳細 - Method: data_free_sst_main, Pattern: safety+medical}
\label{tab:prelim_seeds_datafreesstmain_safety_medical}
\small
\begin{tabular}{\textwidth}{llllRRR}
\toprule
\textbf{Pattern} & \textbf{Method} & \textbf{Alpha} & \textbf{Seed} & \textbf{TrustLLM Raw ASR (ASR$\downarrow$ \%)} & \textbf{Gibberish Ratio (崩壊率$\uparrow$ \%)} & \textbf{BeaverTails Utility (Score$\uparrow$ \%)} \\
\midrule
safety+medical & data\_free\_sst\_main & 0.00 & 42 & 100.00 & 37.50 & 60.31 \\
safety+medical & data\_free\_sst\_main & 0.00 & 43 & 100.00 & 37.50 & 60.31 \\
safety+medical & data\_free\_sst\_main & 0.00 & 44 & 100.00 & 37.50 & 60.31 \\
safety+medical & data\_free\_sst\_main & 0.20 & 42 & 100.00 & 38.75 & 50.94 \\
safety+medical & data\_free\_sst\_main & 0.20 & 43 & 100.00 & 45.31 & 55.00 \\
safety+medical & data\_free\_sst\_main & 0.20 & 44 & 100.00 & 32.50 & 66.88 \\
safety+medical & data\_free\_sst\_main & 0.40 & 42 & 100.00 & 36.25 & 64.69 \\
safety+medical & data\_free\_sst\_main & 0.40 & 43 & 100.00 & 69.06 & 28.12 \\
safety+medical & data\_free\_sst\_main & 0.40 & 44 & 100.00 & 51.88 & 46.56 \\
safety+medical & data\_free\_sst\_main & 0.60 & 42 & 100.00 & 73.75 & 26.88 \\
safety+medical & data\_free\_sst\_main & 0.60 & 43 & 100.00 & 83.12 & 16.88 \\
safety+medical & data\_free\_sst\_main & 0.60 & 44 & 100.00 & 52.19 & 38.44 \\
safety+medical & data\_free\_sst\_main & 0.80 & 42 & 100.00 & 93.44 & 6.25 \\
safety+medical & data\_free\_sst\_main & 0.80 & 43 & 100.00 & 96.25 & 2.50 \\
safety+medical & data\_free\_sst\_main & 0.80 & 44 & 100.00 & 62.81 & 29.69 \\
safety+medical & data\_free\_sst\_main & 1.00 & 42 & 100.00 & 76.88 & 18.44 \\
safety+medical & data\_free\_sst\_main & 1.00 & 43 & 100.00 & 25.31 & 72.19 \\
safety+medical & data\_free\_sst\_main & 1.00 & 44 & 100.00 & 82.19 & 14.37 \\
\bottomrule
\end{tabularx}
\end{table}


% --- 予備実験 シード別詳細 - Method: matena_fisher, Pattern: safety+medical ---
\begin{table}[H]
\centering
\caption{予備実験 シード別詳細 - Method: matena_fisher, Pattern: safety+medical}
\label{tab:prelim_seeds_matenafisher_safety_medical}
\small
\begin{tabular}{\textwidth}{llllRRR}
\toprule
\textbf{Pattern} & \textbf{Method} & \textbf{Alpha} & \textbf{Seed} & \textbf{TrustLLM Raw ASR (ASR$\downarrow$ \%)} & \textbf{Gibberish Ratio (崩壊率$\uparrow$ \%)} & \textbf{BeaverTails Utility (Score$\uparrow$ \%)} \\
\midrule
safety+medical & matena\_fisher & N/A & 42 & 100.00 & 0.31 & 99.69 \\
safety+medical & matena\_fisher & N/A & 43 & 100.00 & 95.62 & 4.38 \\
safety+medical & matena\_fisher & N/A & 44 & 100.00 & 90.62 & 8.44 \\
\bottomrule
\end{tabularx}
\end{table}


% --- 予備実験 集計 (Mean \pm Std) - Pattern: safety+math+code+medical ---
\begin{table}[H]
\centering
\caption{予備実験 集計 (Mean \pm Std) - Pattern: safety+math+code+medical}
\label{tab:prelim_safety_math_code_medical}
\small
\begin{tabular}{\textwidth}{lllRRR}
\toprule
\textbf{Pattern} & \textbf{Method} & \textbf{Alpha} & \textbf{TrustLLM Raw ASR (ASR$\downarrow$ \%)} & \textbf{Gibberish Ratio (崩壊率$\uparrow$ \%)} & \textbf{BeaverTails Utility (Score$\uparrow$ \%)} \\
\midrule
safety+math+code+medical & dare & 0.00 & 100.00 $\pm$ 0.00 & 62.08 $\pm$ 38.13 & 37.92 $\pm$ 37.81 \\
safety+math+code+medical & dare & 0.20 & 100.00 $\pm$ 0.00 & 51.04 $\pm$ 26.81 & 44.69 $\pm$ 22.37 \\
safety+math+code+medical & dare & 0.40 & 100.00 $\pm$ 0.00 & 41.98 $\pm$ 32.73 & 56.04 $\pm$ 28.78 \\
safety+math+code+medical & dare & 0.60 & 100.00 $\pm$ 0.00 & 70.42 $\pm$ 40.72 & 31.46 $\pm$ 41.41 \\
safety+math+code+medical & dare & 0.80 & 100.00 $\pm$ 0.00 & 44.38 $\pm$ 24.88 & 57.71 $\pm$ 24.61 \\
safety+math+code+medical & dare & 1.00 & 2.08 $\pm$ 1.72 & 0.21 $\pm$ 0.18 & 1.77 $\pm$ 1.88 \\
safety+math+code+medical & data\_free\_sst\_main & 0.00 & 100.00 $\pm$ 0.00 & 56.35 $\pm$ 0.18 & 32.40 $\pm$ 0.18 \\
safety+math+code+medical & data\_free\_sst\_main & 0.20 & 100.00 $\pm$ 0.00 & 98.85 $\pm$ 1.26 & 1.46 $\pm$ 1.26 \\
safety+math+code+medical & data\_free\_sst\_main & 0.40 & 99.79 $\pm$ 0.36 & 88.02 $\pm$ 12.36 & 15.83 $\pm$ 18.68 \\
safety+math+code+medical & data\_free\_sst\_main & 0.60 & 95.73 $\pm$ 6.35 & 84.79 $\pm$ 18.14 & 15.62 $\pm$ 15.77 \\
safety+math+code+medical & data\_free\_sst\_main & 0.80 & 95.73 $\pm$ 6.60 & 76.88 $\pm$ 27.72 & 24.58 $\pm$ 23.47 \\
safety+math+code+medical & data\_free\_sst\_main & 1.00 & 98.85 $\pm$ 1.00 & 72.81 $\pm$ 20.73 & 26.15 $\pm$ 15.52 \\
safety+math+code+medical & della & 0.00 & 100.00 $\pm$ 0.00 & 67.40 $\pm$ 28.31 & 35.00 $\pm$ 30.09 \\
safety+math+code+medical & della & 0.20 & 100.00 $\pm$ 0.00 & 52.92 $\pm$ 29.24 & 43.02 $\pm$ 26.10 \\
safety+math+code+medical & della & 0.40 & 100.00 $\pm$ 0.00 & 44.06 $\pm$ 35.83 & 57.92 $\pm$ 35.06 \\
safety+math+code+medical & della & 0.60 & 100.00 $\pm$ 0.00 & 73.65 $\pm$ 24.63 & 28.65 $\pm$ 26.34 \\
safety+math+code+medical & della & 0.80 & 100.00 $\pm$ 0.00 & 46.35 $\pm$ 35.50 & 56.25 $\pm$ 30.51 \\
safety+math+code+medical & della & 1.00 & 1.77 $\pm$ 2.30 & 0.10 $\pm$ 0.18 & 1.35 $\pm$ 1.30 \\
safety+math+code+medical & diagonal\_sst\_main & 0.00 & 100.00 $\pm$ 0.00 & 56.25 $\pm$ 0.00 & 32.29 $\pm$ 0.18 \\
safety+math+code+medical & diagonal\_sst\_main & 0.20 & 100.00 $\pm$ 0.00 & 99.90 $\pm$ 0.18 & 0.00 $\pm$ 0.00 \\
safety+math+code+medical & diagonal\_sst\_main & 0.40 & 100.00 $\pm$ 0.00 & 93.12 $\pm$ 9.21 & 8.44 $\pm$ 9.92 \\
safety+math+code+medical & diagonal\_sst\_main & 0.60 & 97.29 $\pm$ 2.01 & 95.31 $\pm$ 2.19 & 5.10 $\pm$ 0.48 \\
safety+math+code+medical & diagonal\_sst\_main & 0.80 & 90.00 $\pm$ 5.95 & 93.85 $\pm$ 5.24 & 5.62 $\pm$ 5.70 \\
safety+math+code+medical & diagonal\_sst\_main & 1.00 & 45.00 $\pm$ 18.78 & 93.65 $\pm$ 1.72 & 0.52 $\pm$ 0.65 \\
safety+math+code+medical & matena\_fisher & N/A & 100.00 $\pm$ 0.00 & 80.10 $\pm$ 0.18 & 16.56 $\pm$ 1.36 \\
safety+math+code+medical & task\_arithmetic & 0.00 & 100.00 $\pm$ 0.00 & 100.00 $\pm$ 0.00 & 0.00 $\pm$ 0.00 \\
safety+math+code+medical & task\_arithmetic & 0.20 & 100.00 $\pm$ 0.00 & 98.65 $\pm$ 1.83 & 1.15 $\pm$ 0.95 \\
safety+math+code+medical & task\_arithmetic & 0.40 & 99.17 $\pm$ 0.95 & 94.58 $\pm$ 6.14 & 4.27 $\pm$ 2.80 \\
safety+math+code+medical & task\_arithmetic & 0.60 & 66.25 $\pm$ 11.12 & 70.73 $\pm$ 14.23 & 4.90 $\pm$ 1.44 \\
safety+math+code+medical & task\_arithmetic & 0.80 & 12.19 $\pm$ 6.72 & 3.54 $\pm$ 2.35 & 5.62 $\pm$ 3.29 \\
safety+math+code+medical & task\_arithmetic & 1.00 & 0.83 $\pm$ 0.79 & 0.10 $\pm$ 0.18 & 1.46 $\pm$ 0.79 \\
safety+math+code+medical & ties & 0.00 & 100.00 $\pm$ 0.00 & 0.00 $\pm$ 0.00 & 100.00 $\pm$ 0.00 \\
safety+math+code+medical & ties & 0.20 & 100.00 $\pm$ 0.00 & 5.62 $\pm$ 7.44 & 94.90 $\pm$ 7.78 \\
safety+math+code+medical & ties & 0.40 & 100.00 $\pm$ 0.00 & 36.46 $\pm$ 55.23 & 62.60 $\pm$ 54.56 \\
safety+math+code+medical & ties & 0.60 & 100.00 $\pm$ 0.00 & 40.94 $\pm$ 49.27 & 60.21 $\pm$ 50.93 \\
safety+math+code+medical & ties & 0.80 & 100.00 $\pm$ 0.00 & 50.94 $\pm$ 49.06 & 51.35 $\pm$ 48.65 \\
safety+math+code+medical & ties & 1.00 & 0.73 $\pm$ 1.00 & 0.00 $\pm$ 0.00 & 1.46 $\pm$ 0.95 \\
\bottomrule
\end{tabularx}
\end{table}


% --- 予備実験 シード別詳細 - Method: data_free_sst_main, Pattern: safety+math+code+medical ---
\begin{table}[H]
\centering
\caption{予備実験 シード別詳細 - Method: data_free_sst_main, Pattern: safety+math+code+medical}
\label{tab:prelim_seeds_datafreesstmain_safety_math_code_medical}
\small
\begin{tabular}{\textwidth}{llllRRR}
\toprule
\textbf{Pattern} & \textbf{Method} & \textbf{Alpha} & \textbf{Seed} & \textbf{TrustLLM Raw ASR (ASR$\downarrow$ \%)} & \textbf{Gibberish Ratio (崩壊率$\uparrow$ \%)} & \textbf{BeaverTails Utility (Score$\uparrow$ \%)} \\
\midrule
safety+math+code+medical & data\_free\_sst\_main & 0.00 & 42 & 100.00 & 56.25 & 32.19 \\
safety+math+code+medical & data\_free\_sst\_main & 0.00 & 43 & 100.00 & 56.56 & 32.50 \\
safety+math+code+medical & data\_free\_sst\_main & 0.00 & 44 & 100.00 & 56.25 & 32.50 \\
safety+math+code+medical & data\_free\_sst\_main & 0.20 & 42 & 100.00 & 100.00 & 0.00 \\
safety+math+code+medical & data\_free\_sst\_main & 0.20 & 43 & 100.00 & 99.06 & 2.19 \\
safety+math+code+medical & data\_free\_sst\_main & 0.20 & 44 & 100.00 & 97.50 & 2.19 \\
safety+math+code+medical & data\_free\_sst\_main & 0.40 & 42 & 99.38 & 73.75 & 37.19 \\
safety+math+code+medical & data\_free\_sst\_main & 0.40 & 43 & 100.00 & 95.00 & 2.50 \\
safety+math+code+medical & data\_free\_sst\_main & 0.40 & 44 & 100.00 & 95.31 & 7.81 \\
safety+math+code+medical & data\_free\_sst\_main & 0.60 & 42 & 88.44 & 64.38 & 32.50 \\
safety+math+code+medical & data\_free\_sst\_main & 0.60 & 43 & 98.75 & 99.06 & 1.25 \\
safety+math+code+medical & data\_free\_sst\_main & 0.60 & 44 & 100.00 & 90.94 & 13.12 \\
safety+math+code+medical & data\_free\_sst\_main & 0.80 & 42 & 88.12 & 45.00 & 50.94 \\
safety+math+code+medical & data\_free\_sst\_main & 0.80 & 43 & 99.06 & 95.31 & 5.94 \\
safety+math+code+medical & data\_free\_sst\_main & 0.80 & 44 & 100.00 & 90.31 & 16.88 \\
safety+math+code+medical & data\_free\_sst\_main & 1.00 & 42 & 98.12 & 49.38 & 44.06 \\
safety+math+code+medical & data\_free\_sst\_main & 1.00 & 43 & 98.44 & 88.75 & 16.88 \\
safety+math+code+medical & data\_free\_sst\_main & 1.00 & 44 & 100.00 & 80.31 & 17.50 \\
\bottomrule
\end{tabularx}
\end{table}


% --- 予備実験 シード別詳細 - Method: diagonal_sst_main, Pattern: safety+math+code+medical ---
\begin{table}[H]
\centering
\caption{予備実験 シード別詳細 - Method: diagonal_sst_main, Pattern: safety+math+code+medical}
\label{tab:prelim_seeds_diagonalsstmain_safety_math_code_medical}
\small
\begin{tabular}{\textwidth}{llllRRR}
\toprule
\textbf{Pattern} & \textbf{Method} & \textbf{Alpha} & \textbf{Seed} & \textbf{TrustLLM Raw ASR (ASR$\downarrow$ \%)} & \textbf{Gibberish Ratio (崩壊率$\uparrow$ \%)} & \textbf{BeaverTails Utility (Score$\uparrow$ \%)} \\
\midrule
safety+math+code+medical & diagonal\_sst\_main & 0.00 & 42 & 100.00 & 56.25 & 32.19 \\
safety+math+code+medical & diagonal\_sst\_main & 0.00 & 43 & 100.00 & 56.25 & 32.50 \\
safety+math+code+medical & diagonal\_sst\_main & 0.00 & 44 & 100.00 & 56.25 & 32.19 \\
safety+math+code+medical & diagonal\_sst\_main & 0.20 & 42 & 100.00 & 99.69 & 0.00 \\
safety+math+code+medical & diagonal\_sst\_main & 0.20 & 43 & 100.00 & 100.00 & 0.00 \\
safety+math+code+medical & diagonal\_sst\_main & 0.20 & 44 & 100.00 & 100.00 & 0.00 \\
safety+math+code+medical & diagonal\_sst\_main & 0.40 & 42 & 100.00 & 98.75 & 0.94 \\
safety+math+code+medical & diagonal\_sst\_main & 0.40 & 43 & 100.00 & 82.50 & 19.69 \\
safety+math+code+medical & diagonal\_sst\_main & 0.40 & 44 & 100.00 & 98.12 & 4.69 \\
safety+math+code+medical & diagonal\_sst\_main & 0.60 & 42 & 98.75 & 92.81 & 5.62 \\
safety+math+code+medical & diagonal\_sst\_main & 0.60 & 43 & 95.00 & 96.88 & 4.69 \\
safety+math+code+medical & diagonal\_sst\_main & 0.60 & 44 & 98.12 & 96.25 & 5.00 \\
safety+math+code+medical & diagonal\_sst\_main & 0.80 & 42 & 96.88 & 87.81 & 12.19 \\
safety+math+code+medical & diagonal\_sst\_main & 0.80 & 43 & 86.56 & 96.56 & 2.81 \\
safety+math+code+medical & diagonal\_sst\_main & 0.80 & 44 & 86.56 & 97.19 & 1.88 \\
safety+math+code+medical & diagonal\_sst\_main & 1.00 & 42 & 66.25 & 91.88 & 1.25 \\
safety+math+code+medical & diagonal\_sst\_main & 1.00 & 43 & 30.63 & 93.75 & 0.00 \\
safety+math+code+medical & diagonal\_sst\_main & 1.00 & 44 & 38.12 & 95.31 & 0.31 \\
\bottomrule
\end{tabularx}
\end{table}


% --- 予備実験 シード別詳細 - Method: task_arithmetic, Pattern: safety+math+code+medical ---
\begin{table}[H]
\centering
\caption{予備実験 シード別詳細 - Method: task_arithmetic, Pattern: safety+math+code+medical}
\label{tab:prelim_seeds_taskarithmetic_safety_math_code_medical}
\small
\begin{tabular}{\textwidth}{llllRRR}
\toprule
\textbf{Pattern} & \textbf{Method} & \textbf{Alpha} & \textbf{Seed} & \textbf{TrustLLM Raw ASR (ASR$\downarrow$ \%)} & \textbf{Gibberish Ratio (崩壊率$\uparrow$ \%)} & \textbf{BeaverTails Utility (Score$\uparrow$ \%)} \\
\midrule
safety+math+code+medical & task\_arithmetic & 0.00 & 42 & 100.00 & 100.00 & 0.00 \\
safety+math+code+medical & task\_arithmetic & 0.00 & 43 & 100.00 & 100.00 & 0.00 \\
safety+math+code+medical & task\_arithmetic & 0.00 & 44 & 100.00 & 100.00 & 0.00 \\
safety+math+code+medical & task\_arithmetic & 0.20 & 42 & 100.00 & 96.56 & 2.19 \\
safety+math+code+medical & task\_arithmetic & 0.20 & 43 & 100.00 & 100.00 & 0.31 \\
safety+math+code+medical & task\_arithmetic & 0.20 & 44 & 100.00 & 99.38 & 0.94 \\
safety+math+code+medical & task\_arithmetic & 0.40 & 42 & 98.12 & 87.50 & 7.50 \\
safety+math+code+medical & task\_arithmetic & 0.40 & 43 & 100.00 & 98.44 & 2.50 \\
safety+math+code+medical & task\_arithmetic & 0.40 & 44 & 99.38 & 97.81 & 2.81 \\
safety+math+code+medical & task\_arithmetic & 0.60 & 42 & 54.69 & 54.37 & 6.56 \\
safety+math+code+medical & task\_arithmetic & 0.60 & 43 & 76.88 & 80.31 & 4.06 \\
safety+math+code+medical & task\_arithmetic & 0.60 & 44 & 67.19 & 77.50 & 4.06 \\
safety+math+code+medical & task\_arithmetic & 0.80 & 42 & 11.88 & 1.25 & 5.94 \\
safety+math+code+medical & task\_arithmetic & 0.80 & 43 & 19.06 & 5.94 & 8.75 \\
safety+math+code+medical & task\_arithmetic & 0.80 & 44 & 5.62 & 3.44 & 2.19 \\
safety+math+code+medical & task\_arithmetic & 1.00 & 42 & 0.94 & 0.31 & 1.56 \\
safety+math+code+medical & task\_arithmetic & 1.00 & 43 & 1.56 & 0.00 & 2.19 \\
safety+math+code+medical & task\_arithmetic & 1.00 & 44 & 0.00 & 0.00 & 0.62 \\
\bottomrule
\end{tabularx}
\end{table}


% --- 予備実験 シード別詳細 - Method: dare, Pattern: safety+math+code+medical ---
\begin{table}[H]
\centering
\caption{予備実験 シード別詳細 - Method: dare, Pattern: safety+math+code+medical}
\label{tab:prelim_seeds_dare_safety_math_code_medical}
\small
\begin{tabular}{\textwidth}{llllRRR}
\toprule
\textbf{Pattern} & \textbf{Method} & \textbf{Alpha} & \textbf{Seed} & \textbf{TrustLLM Raw ASR (ASR$\downarrow$ \%)} & \textbf{Gibberish Ratio (崩壊率$\uparrow$ \%)} & \textbf{BeaverTails Utility (Score$\uparrow$ \%)} \\
\midrule
safety+math+code+medical & dare & 0.00 & 42 & 100.00 & 62.50 & 38.12 \\
safety+math+code+medical & dare & 0.00 & 43 & 100.00 & 100.00 & 0.00 \\
safety+math+code+medical & dare & 0.00 & 44 & 100.00 & 23.75 & 75.62 \\
safety+math+code+medical & dare & 0.20 & 42 & 100.00 & 69.69 & 29.06 \\
safety+math+code+medical & dare & 0.20 & 43 & 100.00 & 63.12 & 34.69 \\
safety+math+code+medical & dare & 0.20 & 44 & 100.00 & 20.31 & 70.31 \\
safety+math+code+medical & dare & 0.40 & 42 & 100.00 & 39.38 & 57.50 \\
safety+math+code+medical & dare & 0.40 & 43 & 100.00 & 75.94 & 26.56 \\
safety+math+code+medical & dare & 0.40 & 44 & 100.00 & 10.62 & 84.06 \\
safety+math+code+medical & dare & 0.60 & 42 & 100.00 & 88.75 & 11.56 \\
safety+math+code+medical & dare & 0.60 & 43 & 100.00 & 23.75 & 79.06 \\
safety+math+code+medical & dare & 0.60 & 44 & 100.00 & 98.75 & 3.75 \\
safety+math+code+medical & dare & 0.80 & 42 & 100.00 & 45.94 & 50.94 \\
safety+math+code+medical & dare & 0.80 & 43 & 100.00 & 18.75 & 85.00 \\
safety+math+code+medical & dare & 0.80 & 44 & 100.00 & 68.44 & 37.19 \\
safety+math+code+medical & dare & 1.00 & 42 & 2.19 & 0.31 & 1.56 \\
safety+math+code+medical & dare & 1.00 & 43 & 3.75 & 0.00 & 3.75 \\
safety+math+code+medical & dare & 1.00 & 44 & 0.31 & 0.31 & 0.00 \\
\bottomrule
\end{tabularx}
\end{table}


% --- 予備実験 シード別詳細 - Method: ties, Pattern: safety+math+code+medical ---
\begin{table}[H]
\centering
\caption{予備実験 シード別詳細 - Method: ties, Pattern: safety+math+code+medical}
\label{tab:prelim_seeds_ties_safety_math_code_medical}
\small
\begin{tabular}{\textwidth}{llllRRR}
\toprule
\textbf{Pattern} & \textbf{Method} & \textbf{Alpha} & \textbf{Seed} & \textbf{TrustLLM Raw ASR (ASR$\downarrow$ \%)} & \textbf{Gibberish Ratio (崩壊率$\uparrow$ \%)} & \textbf{BeaverTails Utility (Score$\uparrow$ \%)} \\
\midrule
safety+math+code+medical & ties & 0.00 & 42 & 100.00 & 0.00 & 100.00 \\
safety+math+code+medical & ties & 0.00 & 43 & 100.00 & 0.00 & 100.00 \\
safety+math+code+medical & ties & 0.00 & 44 & 100.00 & 0.00 & 100.00 \\
safety+math+code+medical & ties & 0.20 & 42 & 100.00 & 2.81 & 98.75 \\
safety+math+code+medical & ties & 0.20 & 43 & 100.00 & 0.00 & 100.00 \\
safety+math+code+medical & ties & 0.20 & 44 & 100.00 & 14.06 & 85.94 \\
safety+math+code+medical & ties & 0.40 & 42 & 100.00 & 9.38 & 87.81 \\
safety+math+code+medical & ties & 0.40 & 43 & 100.00 & 0.00 & 100.00 \\
safety+math+code+medical & ties & 0.40 & 44 & 100.00 & 100.00 & 0.00 \\
safety+math+code+medical & ties & 0.60 & 42 & 100.00 & 27.19 & 77.81 \\
safety+math+code+medical & ties & 0.60 & 43 & 100.00 & 0.00 & 100.00 \\
safety+math+code+medical & ties & 0.60 & 44 & 100.00 & 95.62 & 2.81 \\
safety+math+code+medical & ties & 0.80 & 42 & 100.00 & 100.00 & 0.31 \\
safety+math+code+medical & ties & 0.80 & 43 & 100.00 & 1.88 & 97.19 \\
safety+math+code+medical & ties & 0.80 & 44 & 100.00 & 50.94 & 56.56 \\
safety+math+code+medical & ties & 1.00 & 42 & 0.31 & 0.00 & 1.25 \\
safety+math+code+medical & ties & 1.00 & 43 & 1.88 & 0.00 & 2.50 \\
safety+math+code+medical & ties & 1.00 & 44 & 0.00 & 0.00 & 0.62 \\
\bottomrule
\end{tabularx}
\end{table}


% --- 予備実験 シード別詳細 - Method: della, Pattern: safety+math+code+medical ---
\begin{table}[H]
\centering
\caption{予備実験 シード別詳細 - Method: della, Pattern: safety+math+code+medical}
\label{tab:prelim_seeds_della_safety_math_code_medical}
\small
\begin{tabular}{\textwidth}{llllRRR}
\toprule
\textbf{Pattern} & \textbf{Method} & \textbf{Alpha} & \textbf{Seed} & \textbf{TrustLLM Raw ASR (ASR$\downarrow$ \%)} & \textbf{Gibberish Ratio (崩壊率$\uparrow$ \%)} & \textbf{BeaverTails Utility (Score$\uparrow$ \%)} \\
\midrule
safety+math+code+medical & della & 0.00 & 42 & 100.00 & 100.00 & 0.31 \\
safety+math+code+medical & della & 0.00 & 43 & 100.00 & 49.06 & 54.06 \\
safety+math+code+medical & della & 0.00 & 44 & 100.00 & 53.12 & 50.62 \\
safety+math+code+medical & della & 0.20 & 42 & 100.00 & 85.94 & 13.44 \\
safety+math+code+medical & della & 0.20 & 43 & 100.00 & 42.50 & 52.81 \\
safety+math+code+medical & della & 0.20 & 44 & 100.00 & 30.31 & 62.81 \\
safety+math+code+medical & della & 0.40 & 42 & 100.00 & 7.19 & 94.06 \\
safety+math+code+medical & della & 0.40 & 43 & 100.00 & 78.75 & 24.06 \\
safety+math+code+medical & della & 0.40 & 44 & 100.00 & 46.25 & 55.62 \\
safety+math+code+medical & della & 0.60 & 42 & 100.00 & 83.44 & 13.75 \\
safety+math+code+medical & della & 0.60 & 43 & 100.00 & 45.62 & 59.06 \\
safety+math+code+medical & della & 0.60 & 44 & 100.00 & 91.88 & 13.12 \\
safety+math+code+medical & della & 0.80 & 42 & 100.00 & 6.25 & 90.00 \\
safety+math+code+medical & della & 0.80 & 43 & 100.00 & 59.06 & 48.12 \\
safety+math+code+medical & della & 0.80 & 44 & 100.00 & 73.75 & 30.63 \\
safety+math+code+medical & della & 1.00 & 42 & 0.94 & 0.31 & 0.94 \\
safety+math+code+medical & della & 1.00 & 43 & 4.38 & 0.00 & 2.81 \\
safety+math+code+medical & della & 1.00 & 44 & 0.00 & 0.00 & 0.31 \\
\bottomrule
\end{tabularx}
\end{table}


% --- 予備実験 シード別詳細 - Method: matena_fisher, Pattern: safety+math+code+medical ---
\begin{table}[H]
\centering
\caption{予備実験 シード別詳細 - Method: matena_fisher, Pattern: safety+math+code+medical}
\label{tab:prelim_seeds_matenafisher_safety_math_code_medical}
\small
\begin{tabular}{\textwidth}{llllRRR}
\toprule
\textbf{Pattern} & \textbf{Method} & \textbf{Alpha} & \textbf{Seed} & \textbf{TrustLLM Raw ASR (ASR$\downarrow$ \%)} & \textbf{Gibberish Ratio (崩壊率$\uparrow$ \%)} & \textbf{BeaverTails Utility (Score$\uparrow$ \%)} \\
\midrule
safety+math+code+medical & matena\_fisher & N/A & 42 & 100.00 & 80.00 & 17.19 \\
safety+math+code+medical & matena\_fisher & N/A & 43 & 100.00 & 80.31 & 17.50 \\
safety+math+code+medical & matena\_fisher & N/A & 44 & 100.00 & 80.00 & 15.00 \\
\bottomrule
\end{tabularx}
\end{table}


% --- 予備実験 集計 (Mean \pm Std) - Pattern: safety+code ---
\begin{table}[H]
\centering
\caption{予備実験 集計 (Mean \pm Std) - Pattern: safety+code}
\label{tab:prelim_safety_code}
\small
\begin{tabular}{\textwidth}{lllRRR}
\toprule
\textbf{Pattern} & \textbf{Method} & \textbf{Alpha} & \textbf{TrustLLM Raw ASR (ASR$\downarrow$ \%)} & \textbf{Gibberish Ratio (崩壊率$\uparrow$ \%)} & \textbf{BeaverTails Utility (Score$\uparrow$ \%)} \\
\midrule
safety+code & dare & 0.00 & 100.00 $\pm$ 0.00 & 56.77 $\pm$ 31.54 & 44.69 $\pm$ 28.23 \\
safety+code & dare & 0.20 & 100.00 $\pm$ 0.00 & 71.77 $\pm$ 5.98 & 25.31 $\pm$ 7.70 \\
safety+code & dare & 0.40 & 100.00 $\pm$ 0.00 & 34.69 $\pm$ 20.12 & 63.96 $\pm$ 17.03 \\
safety+code & dare & 0.60 & 100.00 $\pm$ 0.00 & 36.35 $\pm$ 12.66 & 64.06 $\pm$ 11.87 \\
safety+code & dare & 0.80 & 100.00 $\pm$ 0.00 & 51.46 $\pm$ 22.71 & 48.96 $\pm$ 25.64 \\
safety+code & dare & 1.00 & 1.35 $\pm$ 1.10 & 0.00 $\pm$ 0.00 & 1.35 $\pm$ 1.41 \\
safety+code & data\_free\_sst\_main & 0.00 & 98.44 $\pm$ 0.00 & 74.48 $\pm$ 0.18 & 21.77 $\pm$ 0.90 \\
safety+code & data\_free\_sst\_main & 0.20 & 98.02 $\pm$ 0.48 & 77.92 $\pm$ 2.80 & 26.67 $\pm$ 2.82 \\
safety+code & data\_free\_sst\_main & 0.40 & 94.38 $\pm$ 1.65 & 68.44 $\pm$ 1.08 & 29.58 $\pm$ 3.25 \\
safety+code & data\_free\_sst\_main & 0.60 & 92.60 $\pm$ 2.90 & 64.38 $\pm$ 1.13 & 27.40 $\pm$ 3.77 \\
safety+code & data\_free\_sst\_main & 0.80 & 93.44 $\pm$ 3.75 & 62.81 $\pm$ 2.19 & 34.38 $\pm$ 4.23 \\
safety+code & data\_free\_sst\_main & 1.00 & 99.69 $\pm$ 0.54 & 55.10 $\pm$ 8.09 & 44.06 $\pm$ 10.23 \\
safety+code & della & 0.00 & 100.00 $\pm$ 0.00 & 42.19 $\pm$ 9.99 & 58.33 $\pm$ 11.33 \\
safety+code & della & 0.20 & 100.00 $\pm$ 0.00 & 67.08 $\pm$ 6.31 & 33.02 $\pm$ 5.34 \\
safety+code & della & 0.40 & 100.00 $\pm$ 0.00 & 43.75 $\pm$ 36.43 & 59.27 $\pm$ 30.27 \\
safety+code & della & 0.60 & 100.00 $\pm$ 0.00 & 33.65 $\pm$ 8.21 & 67.19 $\pm$ 5.42 \\
safety+code & della & 0.80 & 100.00 $\pm$ 0.00 & 34.17 $\pm$ 25.33 & 66.56 $\pm$ 23.72 \\
safety+code & della & 1.00 & 1.98 $\pm$ 1.78 & 0.21 $\pm$ 0.18 & 1.67 $\pm$ 1.48 \\
safety+code & diagonal\_sst\_main & 0.00 & 98.44 $\pm$ 0.00 & 74.48 $\pm$ 0.18 & 21.88 $\pm$ 0.83 \\
safety+code & diagonal\_sst\_main & 0.20 & 98.23 $\pm$ 0.65 & 74.90 $\pm$ 1.41 & 32.92 $\pm$ 3.78 \\
safety+code & diagonal\_sst\_main & 0.40 & 98.23 $\pm$ 0.48 & 63.96 $\pm$ 10.31 & 32.29 $\pm$ 3.31 \\
safety+code & diagonal\_sst\_main & 0.60 & 98.12 $\pm$ 1.65 & 76.77 $\pm$ 3.31 & 19.58 $\pm$ 2.90 \\
safety+code & diagonal\_sst\_main & 0.80 & 92.71 $\pm$ 2.95 & 69.38 $\pm$ 3.80 & 25.94 $\pm$ 2.72 \\
safety+code & diagonal\_sst\_main & 1.00 & 77.40 $\pm$ 9.23 & 59.58 $\pm$ 5.24 & 30.21 $\pm$ 5.56 \\
safety+code & led\_merging & N/A & 87.19 $\pm$ 0.00 & 95.94 $\pm$ 0.00 & 5.00 $\pm$ 0.00 \\
safety+code & matena\_fisher & N/A & 95.62 $\pm$ 1.13 & 71.98 $\pm$ 1.57 & 28.33 $\pm$ 2.22 \\
safety+code & mergealign & N/A & 100.00 $\pm$ 0.00 & 100.00 $\pm$ 0.00 & 0.00 $\pm$ 0.00 \\
safety+code & safemerge & N/A & 4.58 $\pm$ 5.81 & 6.25 $\pm$ 10.56 & 2.81 $\pm$ 2.19 \\
safety+code & task\_arithmetic & 0.00 & 71.46 $\pm$ 0.72 & 39.17 $\pm$ 0.36 & 39.17 $\pm$ 0.36 \\
safety+code & task\_arithmetic & 0.20 & 87.71 $\pm$ 0.65 & 46.56 $\pm$ 2.25 & 42.60 $\pm$ 2.80 \\
safety+code & task\_arithmetic & 0.40 & 76.77 $\pm$ 2.08 & 41.98 $\pm$ 2.37 & 33.65 $\pm$ 0.48 \\
safety+code & task\_arithmetic & 0.60 & 82.40 $\pm$ 2.73 & 56.67 $\pm$ 4.07 & 24.38 $\pm$ 3.29 \\
safety+code & task\_arithmetic & 0.80 & 21.67 $\pm$ 4.55 & 5.00 $\pm$ 2.86 & 15.42 $\pm$ 3.93 \\
safety+code & task\_arithmetic & 1.00 & 0.83 $\pm$ 0.79 & 0.10 $\pm$ 0.18 & 1.46 $\pm$ 0.79 \\
safety+code & ties & 0.00 & 63.85 $\pm$ 0.72 & 57.08 $\pm$ 1.44 & 18.02 $\pm$ 0.36 \\
safety+code & ties & 0.20 & 92.50 $\pm$ 2.48 & 77.60 $\pm$ 4.03 & 20.10 $\pm$ 4.42 \\
safety+code & ties & 0.40 & 96.98 $\pm$ 1.26 & 62.60 $\pm$ 4.61 & 26.77 $\pm$ 4.02 \\
safety+code & ties & 0.60 & 99.38 $\pm$ 0.31 & 62.29 $\pm$ 6.59 & 34.58 $\pm$ 2.37 \\
safety+code & ties & 0.80 & 100.00 $\pm$ 0.00 & 68.33 $\pm$ 23.31 & 31.35 $\pm$ 23.08 \\
safety+code & ties & 1.00 & 0.73 $\pm$ 1.00 & 0.00 $\pm$ 0.00 & 1.46 $\pm$ 0.95 \\
\bottomrule
\end{tabularx}
\end{table}


% --- 予備実験 シード別詳細 - Method: task_arithmetic, Pattern: safety+code ---
\begin{table}[H]
\centering
\caption{予備実験 シード別詳細 - Method: task_arithmetic, Pattern: safety+code}
\label{tab:prelim_seeds_taskarithmetic_safety_code}
\small
\begin{tabular}{\textwidth}{llllRRR}
\toprule
\textbf{Pattern} & \textbf{Method} & \textbf{Alpha} & \textbf{Seed} & \textbf{TrustLLM Raw ASR (ASR$\downarrow$ \%)} & \textbf{Gibberish Ratio (崩壊率$\uparrow$ \%)} & \textbf{BeaverTails Utility (Score$\uparrow$ \%)} \\
\midrule
safety+code & task\_arithmetic & 0.00 & 42 & 70.62 & 38.75 & 38.75 \\
safety+code & task\_arithmetic & 0.00 & 43 & 71.88 & 39.38 & 39.38 \\
safety+code & task\_arithmetic & 0.00 & 44 & 71.88 & 39.38 & 39.38 \\
safety+code & task\_arithmetic & 0.20 & 42 & 88.44 & 48.44 & 44.06 \\
safety+code & task\_arithmetic & 0.20 & 43 & 87.50 & 44.06 & 44.38 \\
safety+code & task\_arithmetic & 0.20 & 44 & 87.19 & 47.19 & 39.38 \\
safety+code & task\_arithmetic & 0.40 & 42 & 75.00 & 44.69 & 34.06 \\
safety+code & task\_arithmetic & 0.40 & 43 & 79.06 & 40.31 & 33.75 \\
safety+code & task\_arithmetic & 0.40 & 44 & 76.25 & 40.94 & 33.12 \\
safety+code & task\_arithmetic & 0.60 & 42 & 83.12 & 56.88 & 24.69 \\
safety+code & task\_arithmetic & 0.60 & 43 & 84.69 & 52.50 & 27.50 \\
safety+code & task\_arithmetic & 0.60 & 44 & 79.38 & 60.62 & 20.94 \\
safety+code & task\_arithmetic & 0.80 & 42 & 19.69 & 4.38 & 15.94 \\
safety+code & task\_arithmetic & 0.80 & 43 & 26.88 & 2.50 & 19.06 \\
safety+code & task\_arithmetic & 0.80 & 44 & 18.44 & 8.12 & 11.25 \\
safety+code & task\_arithmetic & 1.00 & 42 & 0.94 & 0.31 & 1.56 \\
safety+code & task\_arithmetic & 1.00 & 43 & 1.56 & 0.00 & 2.19 \\
safety+code & task\_arithmetic & 1.00 & 44 & 0.00 & 0.00 & 0.62 \\
\bottomrule
\end{tabularx}
\end{table}


% --- 予備実験 シード別詳細 - Method: della, Pattern: safety+code ---
\begin{table}[H]
\centering
\caption{予備実験 シード別詳細 - Method: della, Pattern: safety+code}
\label{tab:prelim_seeds_della_safety_code}
\small
\begin{tabular}{\textwidth}{llllRRR}
\toprule
\textbf{Pattern} & \textbf{Method} & \textbf{Alpha} & \textbf{Seed} & \textbf{TrustLLM Raw ASR (ASR$\downarrow$ \%)} & \textbf{Gibberish Ratio (崩壊率$\uparrow$ \%)} & \textbf{BeaverTails Utility (Score$\uparrow$ \%)} \\
\midrule
safety+code & della & 0.00 & 42 & 100.00 & 34.38 & 65.94 \\
safety+code & della & 0.00 & 43 & 100.00 & 53.44 & 45.31 \\
safety+code & della & 0.00 & 44 & 100.00 & 38.75 & 63.75 \\
safety+code & della & 0.20 & 42 & 100.00 & 68.12 & 36.56 \\
safety+code & della & 0.20 & 43 & 100.00 & 72.81 & 26.88 \\
safety+code & della & 0.20 & 44 & 100.00 & 60.31 & 35.62 \\
safety+code & della & 0.40 & 42 & 100.00 & 79.38 & 30.63 \\
safety+code & della & 0.40 & 43 & 100.00 & 6.56 & 90.94 \\
safety+code & della & 0.40 & 44 & 100.00 & 45.31 & 56.25 \\
safety+code & della & 0.60 & 42 & 100.00 & 43.12 & 60.94 \\
safety+code & della & 0.60 & 43 & 100.00 & 28.75 & 70.00 \\
safety+code & della & 0.60 & 44 & 100.00 & 29.06 & 70.62 \\
safety+code & della & 0.80 & 42 & 100.00 & 33.12 & 69.38 \\
safety+code & della & 0.80 & 43 & 100.00 & 60.00 & 41.56 \\
safety+code & della & 0.80 & 44 & 100.00 & 9.38 & 88.75 \\
safety+code & della & 1.00 & 42 & 3.44 & 0.31 & 2.81 \\
safety+code & della & 1.00 & 43 & 2.50 & 0.00 & 2.19 \\
safety+code & della & 1.00 & 44 & 0.00 & 0.31 & 0.00 \\
\bottomrule
\end{tabularx}
\end{table}


% --- 予備実験 シード別詳細 - Method: dare, Pattern: safety+code ---
\begin{table}[H]
\centering
\caption{予備実験 シード別詳細 - Method: dare, Pattern: safety+code}
\label{tab:prelim_seeds_dare_safety_code}
\small
\begin{tabular}{\textwidth}{llllRRR}
\toprule
\textbf{Pattern} & \textbf{Method} & \textbf{Alpha} & \textbf{Seed} & \textbf{TrustLLM Raw ASR (ASR$\downarrow$ \%)} & \textbf{Gibberish Ratio (崩壊率$\uparrow$ \%)} & \textbf{BeaverTails Utility (Score$\uparrow$ \%)} \\
\midrule
safety+code & dare & 0.00 & 42 & 100.00 & 69.06 & 38.12 \\
safety+code & dare & 0.00 & 43 & 100.00 & 20.94 & 75.62 \\
safety+code & dare & 0.00 & 44 & 100.00 & 80.31 & 20.31 \\
safety+code & dare & 0.20 & 42 & 100.00 & 78.12 & 17.19 \\
safety+code & dare & 0.20 & 43 & 100.00 & 66.25 & 32.50 \\
safety+code & dare & 0.20 & 44 & 100.00 & 70.94 & 26.25 \\
safety+code & dare & 0.40 & 42 & 100.00 & 55.94 & 46.88 \\
safety+code & dare & 0.40 & 43 & 100.00 & 32.19 & 64.06 \\
safety+code & dare & 0.40 & 44 & 100.00 & 15.94 & 80.94 \\
safety+code & dare & 0.60 & 42 & 100.00 & 21.88 & 77.50 \\
safety+code & dare & 0.60 & 43 & 100.00 & 41.88 & 55.00 \\
safety+code & dare & 0.60 & 44 & 100.00 & 45.31 & 59.69 \\
safety+code & dare & 0.80 & 42 & 100.00 & 62.81 & 38.75 \\
safety+code & dare & 0.80 & 43 & 100.00 & 66.25 & 30.00 \\
safety+code & dare & 0.80 & 44 & 100.00 & 25.31 & 78.12 \\
safety+code & dare & 1.00 & 42 & 1.25 & 0.00 & 1.25 \\
safety+code & dare & 1.00 & 43 & 2.50 & 0.00 & 2.81 \\
safety+code & dare & 1.00 & 44 & 0.31 & 0.00 & 0.00 \\
\bottomrule
\end{tabularx}
\end{table}


% --- 予備実験 シード別詳細 - Method: ties, Pattern: safety+code ---
\begin{table}[H]
\centering
\caption{予備実験 シード別詳細 - Method: ties, Pattern: safety+code}
\label{tab:prelim_seeds_ties_safety_code}
\small
\begin{tabular}{\textwidth}{llllRRR}
\toprule
\textbf{Pattern} & \textbf{Method} & \textbf{Alpha} & \textbf{Seed} & \textbf{TrustLLM Raw ASR (ASR$\downarrow$ \%)} & \textbf{Gibberish Ratio (崩壊率$\uparrow$ \%)} & \textbf{BeaverTails Utility (Score$\uparrow$ \%)} \\
\midrule
safety+code & ties & 0.00 & 42 & 63.44 & 56.25 & 17.81 \\
safety+code & ties & 0.00 & 43 & 64.69 & 58.75 & 18.44 \\
safety+code & ties & 0.00 & 44 & 63.44 & 56.25 & 17.81 \\
safety+code & ties & 0.20 & 42 & 93.44 & 80.94 & 15.00 \\
safety+code & ties & 0.20 & 43 & 89.69 & 73.12 & 22.81 \\
safety+code & ties & 0.20 & 44 & 94.38 & 78.75 & 22.50 \\
safety+code & ties & 0.40 & 42 & 98.12 & 59.06 & 29.69 \\
safety+code & ties & 0.40 & 43 & 95.62 & 60.94 & 22.19 \\
safety+code & ties & 0.40 & 44 & 97.19 & 67.81 & 28.44 \\
safety+code & ties & 0.60 & 42 & 99.69 & 65.94 & 35.62 \\
safety+code & ties & 0.60 & 43 & 99.38 & 54.69 & 36.25 \\
safety+code & ties & 0.60 & 44 & 99.06 & 66.25 & 31.87 \\
safety+code & ties & 0.80 & 42 & 100.00 & 44.38 & 56.56 \\
safety+code & ties & 0.80 & 43 & 100.00 & 69.69 & 26.25 \\
safety+code & ties & 0.80 & 44 & 100.00 & 90.94 & 11.25 \\
safety+code & ties & 1.00 & 42 & 0.31 & 0.00 & 1.25 \\
safety+code & ties & 1.00 & 43 & 1.88 & 0.00 & 2.50 \\
safety+code & ties & 1.00 & 44 & 0.00 & 0.00 & 0.62 \\
\bottomrule
\end{tabularx}
\end{table}


% --- 予備実験 シード別詳細 - Method: diagonal_sst_main, Pattern: safety+code ---
\begin{table}[H]
\centering
\caption{予備実験 シード別詳細 - Method: diagonal_sst_main, Pattern: safety+code}
\label{tab:prelim_seeds_diagonalsstmain_safety_code}
\small
\begin{tabular}{\textwidth}{llllRRR}
\toprule
\textbf{Pattern} & \textbf{Method} & \textbf{Alpha} & \textbf{Seed} & \textbf{TrustLLM Raw ASR (ASR$\downarrow$ \%)} & \textbf{Gibberish Ratio (崩壊率$\uparrow$ \%)} & \textbf{BeaverTails Utility (Score$\uparrow$ \%)} \\
\midrule
safety+code & diagonal\_sst\_main & 0.00 & 42 & 98.44 & 74.69 & 22.81 \\
safety+code & diagonal\_sst\_main & 0.00 & 43 & 98.44 & 74.38 & 21.25 \\
safety+code & diagonal\_sst\_main & 0.00 & 44 & 98.44 & 74.38 & 21.56 \\
safety+code & diagonal\_sst\_main & 0.20 & 42 & 97.50 & 73.44 & 37.19 \\
safety+code & diagonal\_sst\_main & 0.20 & 43 & 98.44 & 75.00 & 31.56 \\
safety+code & diagonal\_sst\_main & 0.20 & 44 & 98.75 & 76.25 & 30.00 \\
safety+code & diagonal\_sst\_main & 0.40 & 42 & 97.81 & 53.75 & 35.31 \\
safety+code & diagonal\_sst\_main & 0.40 & 43 & 98.75 & 63.75 & 28.75 \\
safety+code & diagonal\_sst\_main & 0.40 & 44 & 98.12 & 74.38 & 32.81 \\
safety+code & diagonal\_sst\_main & 0.60 & 42 & 96.25 & 76.25 & 17.19 \\
safety+code & diagonal\_sst\_main & 0.60 & 43 & 98.75 & 73.75 & 22.81 \\
safety+code & diagonal\_sst\_main & 0.60 & 44 & 99.38 & 80.31 & 18.75 \\
safety+code & diagonal\_sst\_main & 0.80 & 42 & 89.38 & 66.88 & 27.19 \\
safety+code & diagonal\_sst\_main & 0.80 & 43 & 93.75 & 73.75 & 22.81 \\
safety+code & diagonal\_sst\_main & 0.80 & 44 & 95.00 & 67.50 & 27.81 \\
safety+code & diagonal\_sst\_main & 1.00 & 42 & 68.44 & 56.88 & 29.06 \\
safety+code & diagonal\_sst\_main & 1.00 & 43 & 76.88 & 65.62 & 25.31 \\
safety+code & diagonal\_sst\_main & 1.00 & 44 & 86.88 & 56.25 & 36.25 \\
\bottomrule
\end{tabularx}
\end{table}


% --- 予備実験 シード別詳細 - Method: safemerge, Pattern: safety+code ---
\begin{table}[H]
\centering
\caption{予備実験 シード別詳細 - Method: safemerge, Pattern: safety+code}
\label{tab:prelim_seeds_safemerge_safety_code}
\small
\begin{tabular}{\textwidth}{llllRRR}
\toprule
\textbf{Pattern} & \textbf{Method} & \textbf{Alpha} & \textbf{Seed} & \textbf{TrustLLM Raw ASR (ASR$\downarrow$ \%)} & \textbf{Gibberish Ratio (崩壊率$\uparrow$ \%)} & \textbf{BeaverTails Utility (Score$\uparrow$ \%)} \\
\midrule
safety+code & safemerge & N/A & 42 & 11.25 & 18.44 & 5.31 \\
safety+code & safemerge & N/A & 43 & 0.62 & 0.31 & 1.88 \\
safety+code & safemerge & N/A & 44 & 1.88 & 0.00 & 1.25 \\
\bottomrule
\end{tabularx}
\end{table}


% --- 予備実験 シード別詳細 - Method: data_free_sst_main, Pattern: safety+code ---
\begin{table}[H]
\centering
\caption{予備実験 シード別詳細 - Method: data_free_sst_main, Pattern: safety+code}
\label{tab:prelim_seeds_datafreesstmain_safety_code}
\small
\begin{tabular}{\textwidth}{llllRRR}
\toprule
\textbf{Pattern} & \textbf{Method} & \textbf{Alpha} & \textbf{Seed} & \textbf{TrustLLM Raw ASR (ASR$\downarrow$ \%)} & \textbf{Gibberish Ratio (崩壊率$\uparrow$ \%)} & \textbf{BeaverTails Utility (Score$\uparrow$ \%)} \\
\midrule
safety+code & data\_free\_sst\_main & 0.00 & 42 & 98.44 & 74.38 & 21.25 \\
safety+code & data\_free\_sst\_main & 0.00 & 43 & 98.44 & 74.38 & 21.25 \\
safety+code & data\_free\_sst\_main & 0.00 & 44 & 98.44 & 74.69 & 22.81 \\
safety+code & data\_free\_sst\_main & 0.20 & 42 & 98.12 & 74.69 & 23.75 \\
safety+code & data\_free\_sst\_main & 0.20 & 43 & 97.50 & 79.69 & 26.88 \\
safety+code & data\_free\_sst\_main & 0.20 & 44 & 98.44 & 79.38 & 29.38 \\
safety+code & data\_free\_sst\_main & 0.40 & 42 & 95.62 & 69.06 & 32.19 \\
safety+code & data\_free\_sst\_main & 0.40 & 43 & 92.50 & 67.19 & 30.63 \\
safety+code & data\_free\_sst\_main & 0.40 & 44 & 95.00 & 69.06 & 25.94 \\
safety+code & data\_free\_sst\_main & 0.60 & 42 & 90.62 & 63.44 & 30.94 \\
safety+code & data\_free\_sst\_main & 0.60 & 43 & 91.25 & 65.62 & 23.44 \\
safety+code & data\_free\_sst\_main & 0.60 & 44 & 95.94 & 64.06 & 27.81 \\
safety+code & data\_free\_sst\_main & 0.80 & 42 & 93.44 & 65.00 & 34.69 \\
safety+code & data\_free\_sst\_main & 0.80 & 43 & 89.69 & 60.62 & 30.00 \\
safety+code & data\_free\_sst\_main & 0.80 & 44 & 97.19 & 62.81 & 38.44 \\
safety+code & data\_free\_sst\_main & 1.00 & 42 & 100.00 & 61.25 & 32.81 \\
safety+code & data\_free\_sst\_main & 1.00 & 43 & 99.06 & 58.13 & 46.56 \\
safety+code & data\_free\_sst\_main & 1.00 & 44 & 100.00 & 45.94 & 52.81 \\
\bottomrule
\end{tabularx}
\end{table}


% --- 予備実験 シード別詳細 - Method: mergealign, Pattern: safety+code ---
\begin{table}[H]
\centering
\caption{予備実験 シード別詳細 - Method: mergealign, Pattern: safety+code}
\label{tab:prelim_seeds_mergealign_safety_code}
\small
\begin{tabular}{\textwidth}{llllRRR}
\toprule
\textbf{Pattern} & \textbf{Method} & \textbf{Alpha} & \textbf{Seed} & \textbf{TrustLLM Raw ASR (ASR$\downarrow$ \%)} & \textbf{Gibberish Ratio (崩壊率$\uparrow$ \%)} & \textbf{BeaverTails Utility (Score$\uparrow$ \%)} \\
\midrule
safety+code & mergealign & N/A & 42 & 100.00 & 100.00 & 0.00 \\
safety+code & mergealign & N/A & 43 & 100.00 & 100.00 & 0.00 \\
safety+code & mergealign & N/A & 44 & 100.00 & 100.00 & 0.00 \\
\bottomrule
\end{tabularx}
\end{table}


% --- 予備実験 シード別詳細 - Method: led_merging, Pattern: safety+code ---
\begin{table}[H]
\centering
\caption{予備実験 シード別詳細 - Method: led_merging, Pattern: safety+code}
\label{tab:prelim_seeds_ledmerging_safety_code}
\small
\begin{tabular}{\textwidth}{llllRRR}
\toprule
\textbf{Pattern} & \textbf{Method} & \textbf{Alpha} & \textbf{Seed} & \textbf{TrustLLM Raw ASR (ASR$\downarrow$ \%)} & \textbf{Gibberish Ratio (崩壊率$\uparrow$ \%)} & \textbf{BeaverTails Utility (Score$\uparrow$ \%)} \\
\midrule
safety+code & led\_merging & N/A & 42 & 87.19 & 95.94 & 5.00 \\
safety+code & led\_merging & N/A & 43 & 87.19 & 95.94 & 5.00 \\
safety+code & led\_merging & N/A & 44 & 87.19 & 95.94 & 5.00 \\
\bottomrule
\end{tabularx}
\end{table}


% --- 予備実験 シード別詳細 - Method: matena_fisher, Pattern: safety+code ---
\begin{table}[H]
\centering
\caption{予備実験 シード別詳細 - Method: matena_fisher, Pattern: safety+code}
\label{tab:prelim_seeds_matenafisher_safety_code}
\small
\begin{tabular}{\textwidth}{llllRRR}
\toprule
\textbf{Pattern} & \textbf{Method} & \textbf{Alpha} & \textbf{Seed} & \textbf{TrustLLM Raw ASR (ASR$\downarrow$ \%)} & \textbf{Gibberish Ratio (崩壊率$\uparrow$ \%)} & \textbf{BeaverTails Utility (Score$\uparrow$ \%)} \\
\midrule
safety+code & matena\_fisher & N/A & 42 & 96.56 & 72.19 & 25.94 \\
safety+code & matena\_fisher & N/A & 43 & 94.38 & 73.44 & 28.75 \\
safety+code & matena\_fisher & N/A & 44 & 95.94 & 70.31 & 30.31 \\
\bottomrule
\end{tabularx}
\end{table}


% --- 予備実験 集計 (Mean \pm Std) - Pattern: safety+math ---
\begin{table}[H]
\centering
\caption{予備実験 集計 (Mean \pm Std) - Pattern: safety+math}
\label{tab:prelim_safety_math}
\small
\begin{tabular}{\textwidth}{lllRRR}
\toprule
\textbf{Pattern} & \textbf{Method} & \textbf{Alpha} & \textbf{TrustLLM Raw ASR (ASR$\downarrow$ \%)} & \textbf{Gibberish Ratio (崩壊率$\uparrow$ \%)} & \textbf{BeaverTails Utility (Score$\uparrow$ \%)} \\
\midrule
safety+math & dare & 0.00 & 70.00 $\pm$ 2.72 & 19.17 $\pm$ 2.08 & 35.31 $\pm$ 1.13 \\
safety+math & dare & 0.20 & 24.27 $\pm$ 3.70 & 30.42 $\pm$ 5.08 & 16.46 $\pm$ 1.48 \\
safety+math & dare & 0.40 & 20.42 $\pm$ 2.01 & 27.08 $\pm$ 2.95 & 12.19 $\pm$ 0.31 \\
safety+math & dare & 0.60 & 17.19 $\pm$ 0.83 & 18.65 $\pm$ 0.95 & 9.69 $\pm$ 2.86 \\
safety+math & dare & 0.80 & 12.19 $\pm$ 4.73 & 21.15 $\pm$ 4.07 & 7.71 $\pm$ 3.13 \\
safety+math & dare & 1.00 & 1.15 $\pm$ 1.26 & 0.31 $\pm$ 0.31 & 1.04 $\pm$ 1.00 \\
safety+math & data\_free\_sst\_main & 0.00 & 70.42 $\pm$ 0.18 & 17.08 $\pm$ 0.36 & 35.42 $\pm$ 0.36 \\
safety+math & data\_free\_sst\_main & 0.20 & 60.94 $\pm$ 2.72 & 18.65 $\pm$ 1.72 & 31.15 $\pm$ 1.10 \\
safety+math & data\_free\_sst\_main & 0.40 & 47.08 $\pm$ 1.54 & 20.00 $\pm$ 5.14 & 25.10 $\pm$ 1.88 \\
safety+math & data\_free\_sst\_main & 0.60 & 36.88 $\pm$ 3.25 & 36.04 $\pm$ 4.24 & 21.15 $\pm$ 0.36 \\
safety+math & data\_free\_sst\_main & 0.80 & 22.92 $\pm$ 0.18 & 37.29 $\pm$ 0.72 & 16.56 $\pm$ 1.56 \\
safety+math & data\_free\_sst\_main & 1.00 & 22.50 $\pm$ 3.79 & 17.29 $\pm$ 5.20 & 14.06 $\pm$ 1.65 \\
safety+math & della & 0.00 & 68.96 $\pm$ 1.30 & 20.83 $\pm$ 1.41 & 35.62 $\pm$ 1.25 \\
safety+math & della & 0.20 & 24.58 $\pm$ 3.34 & 28.33 $\pm$ 4.70 & 17.60 $\pm$ 2.51 \\
safety+math & della & 0.40 & 21.56 $\pm$ 0.62 & 23.12 $\pm$ 11.56 & 10.94 $\pm$ 0.00 \\
safety+math & della & 0.60 & 19.06 $\pm$ 3.17 & 20.62 $\pm$ 5.20 & 7.81 $\pm$ 2.05 \\
safety+math & della & 0.80 & 13.54 $\pm$ 3.82 & 17.71 $\pm$ 2.69 & 8.65 $\pm$ 2.95 \\
safety+math & della & 1.00 & 1.15 $\pm$ 1.48 & 0.31 $\pm$ 0.31 & 1.46 $\pm$ 1.10 \\
safety+math & diagonal\_sst\_main & 0.00 & 70.42 $\pm$ 0.18 & 17.08 $\pm$ 0.36 & 35.42 $\pm$ 0.36 \\
safety+math & diagonal\_sst\_main & 0.20 & 56.98 $\pm$ 1.72 & 21.77 $\pm$ 3.07 & 26.88 $\pm$ 1.36 \\
safety+math & diagonal\_sst\_main & 0.40 & 41.25 $\pm$ 2.17 & 32.50 $\pm$ 7.04 & 20.21 $\pm$ 0.18 \\
safety+math & diagonal\_sst\_main & 0.60 & 24.38 $\pm$ 1.74 & 58.44 $\pm$ 5.94 & 12.81 $\pm$ 2.44 \\
safety+math & diagonal\_sst\_main & 0.80 & 11.25 $\pm$ 0.54 & 38.33 $\pm$ 26.46 & 6.77 $\pm$ 0.48 \\
safety+math & diagonal\_sst\_main & 1.00 & 2.71 $\pm$ 1.00 & 3.12 $\pm$ 2.56 & 2.08 $\pm$ 0.18 \\
safety+math & led\_merging & N/A & 87.19 $\pm$ 0.00 & 96.15 $\pm$ 0.18 & 5.00 $\pm$ 0.00 \\
safety+math & matena\_fisher & N/A & 19.79 $\pm$ 1.18 & 67.19 $\pm$ 3.52 & 10.00 $\pm$ 1.65 \\
safety+math & mergealign & N/A & 100.00 $\pm$ 0.00 & 100.00 $\pm$ 0.00 & 0.00 $\pm$ 0.00 \\
safety+math & safemerge & N/A & 14.27 $\pm$ 22.55 & 31.15 $\pm$ 53.95 & 3.33 $\pm$ 3.34 \\
safety+math & task\_arithmetic & 0.00 & 66.88 $\pm$ 0.00 & 19.69 $\pm$ 0.00 & 35.62 $\pm$ 0.00 \\
safety+math & task\_arithmetic & 0.20 & 52.81 $\pm$ 2.72 & 26.25 $\pm$ 1.90 & 23.65 $\pm$ 1.30 \\
safety+math & task\_arithmetic & 0.40 & 31.15 $\pm$ 4.08 & 48.33 $\pm$ 5.39 & 16.46 $\pm$ 0.36 \\
safety+math & task\_arithmetic & 0.60 & 9.69 $\pm$ 0.31 & 47.92 $\pm$ 18.35 & 6.35 $\pm$ 1.91 \\
safety+math & task\_arithmetic & 0.80 & 1.67 $\pm$ 1.30 & 0.94 $\pm$ 0.83 & 1.67 $\pm$ 1.00 \\
safety+math & task\_arithmetic & 1.00 & 0.83 $\pm$ 0.79 & 0.10 $\pm$ 0.18 & 1.46 $\pm$ 0.79 \\
safety+math & ties & 0.00 & 72.50 $\pm$ 0.00 & 30.31 $\pm$ 0.00 & 29.38 $\pm$ 0.00 \\
safety+math & ties & 0.20 & 25.00 $\pm$ 4.23 & 38.75 $\pm$ 2.78 & 14.48 $\pm$ 2.55 \\
safety+math & ties & 0.40 & 26.77 $\pm$ 3.19 & 39.79 $\pm$ 3.28 & 14.27 $\pm$ 0.79 \\
safety+math & ties & 0.60 & 18.65 $\pm$ 4.77 & 36.67 $\pm$ 4.77 & 11.25 $\pm$ 1.08 \\
safety+math & ties & 0.80 & 9.79 $\pm$ 5.56 & 19.58 $\pm$ 19.59 & 4.90 $\pm$ 2.37 \\
safety+math & ties & 1.00 & 0.73 $\pm$ 1.00 & 0.00 $\pm$ 0.00 & 1.46 $\pm$ 0.95 \\
\bottomrule
\end{tabularx}
\end{table}


% --- 予備実験 シード別詳細 - Method: data_free_sst_main, Pattern: safety+math ---
\begin{table}[H]
\centering
\caption{予備実験 シード別詳細 - Method: data_free_sst_main, Pattern: safety+math}
\label{tab:prelim_seeds_datafreesstmain_safety_math}
\small
\begin{tabular}{\textwidth}{llllRRR}
\toprule
\textbf{Pattern} & \textbf{Method} & \textbf{Alpha} & \textbf{Seed} & \textbf{TrustLLM Raw ASR (ASR$\downarrow$ \%)} & \textbf{Gibberish Ratio (崩壊率$\uparrow$ \%)} & \textbf{BeaverTails Utility (Score$\uparrow$ \%)} \\
\midrule
safety+math & data\_free\_sst\_main & 0.00 & 42 & 70.31 & 16.88 & 35.62 \\
safety+math & data\_free\_sst\_main & 0.00 & 43 & 70.62 & 17.50 & 35.00 \\
safety+math & data\_free\_sst\_main & 0.00 & 44 & 70.31 & 16.88 & 35.62 \\
safety+math & data\_free\_sst\_main & 0.20 & 42 & 59.69 & 18.75 & 32.19 \\
safety+math & data\_free\_sst\_main & 0.20 & 43 & 64.06 & 16.88 & 31.25 \\
safety+math & data\_free\_sst\_main & 0.20 & 44 & 59.06 & 20.31 & 30.00 \\
safety+math & data\_free\_sst\_main & 0.40 & 42 & 48.12 & 22.81 & 25.31 \\
safety+math & data\_free\_sst\_main & 0.40 & 43 & 47.81 & 23.12 & 23.12 \\
safety+math & data\_free\_sst\_main & 0.40 & 44 & 45.31 & 14.06 & 26.88 \\
safety+math & data\_free\_sst\_main & 0.60 & 42 & 33.12 & 36.56 & 20.94 \\
safety+math & data\_free\_sst\_main & 0.60 & 43 & 38.75 & 40.00 & 21.56 \\
safety+math & data\_free\_sst\_main & 0.60 & 44 & 38.75 & 31.56 & 20.94 \\
safety+math & data\_free\_sst\_main & 0.80 & 42 & 22.81 & 36.88 & 16.56 \\
safety+math & data\_free\_sst\_main & 0.80 & 43 & 22.81 & 36.88 & 15.00 \\
safety+math & data\_free\_sst\_main & 0.80 & 44 & 23.12 & 38.12 & 18.12 \\
safety+math & data\_free\_sst\_main & 1.00 & 42 & 18.44 & 13.12 & 12.81 \\
safety+math & data\_free\_sst\_main & 1.00 & 43 & 23.12 & 15.62 & 13.44 \\
safety+math & data\_free\_sst\_main & 1.00 & 44 & 25.94 & 23.12 & 15.94 \\
\bottomrule
\end{tabularx}
\end{table}


% --- 予備実験 シード別詳細 - Method: della, Pattern: safety+math ---
\begin{table}[H]
\centering
\caption{予備実験 シード別詳細 - Method: della, Pattern: safety+math}
\label{tab:prelim_seeds_della_safety_math}
\small
\begin{tabular}{\textwidth}{llllRRR}
\toprule
\textbf{Pattern} & \textbf{Method} & \textbf{Alpha} & \textbf{Seed} & \textbf{TrustLLM Raw ASR (ASR$\downarrow$ \%)} & \textbf{Gibberish Ratio (崩壊率$\uparrow$ \%)} & \textbf{BeaverTails Utility (Score$\uparrow$ \%)} \\
\midrule
safety+math & della & 0.00 & 42 & 70.00 & 22.19 & 35.62 \\
safety+math & della & 0.00 & 43 & 67.50 & 20.94 & 36.88 \\
safety+math & della & 0.00 & 44 & 69.38 & 19.38 & 34.38 \\
safety+math & della & 0.20 & 42 & 22.81 & 25.31 & 15.00 \\
safety+math & della & 0.20 & 43 & 22.50 & 33.75 & 17.81 \\
safety+math & della & 0.20 & 44 & 28.44 & 25.94 & 20.00 \\
safety+math & della & 0.40 & 42 & 20.94 & 20.94 & 10.94 \\
safety+math & della & 0.40 & 43 & 21.56 & 35.62 & 10.94 \\
safety+math & della & 0.40 & 44 & 22.19 & 12.81 & 10.94 \\
safety+math & della & 0.60 & 42 & 16.25 & 26.56 & 5.94 \\
safety+math & della & 0.60 & 43 & 18.44 & 18.44 & 7.50 \\
safety+math & della & 0.60 & 44 & 22.50 & 16.88 & 10.00 \\
safety+math & della & 0.80 & 42 & 9.38 & 15.31 & 5.31 \\
safety+math & della & 0.80 & 43 & 14.37 & 17.19 & 9.69 \\
safety+math & della & 0.80 & 44 & 16.88 & 20.62 & 10.94 \\
safety+math & della & 1.00 & 42 & 0.62 & 0.62 & 2.50 \\
safety+math & della & 1.00 & 43 & 2.81 & 0.31 & 1.56 \\
safety+math & della & 1.00 & 44 & 0.00 & 0.00 & 0.31 \\
\bottomrule
\end{tabularx}
\end{table}


% --- 予備実験 シード別詳細 - Method: safemerge, Pattern: safety+math ---
\begin{table}[H]
\centering
\caption{予備実験 シード別詳細 - Method: safemerge, Pattern: safety+math}
\label{tab:prelim_seeds_safemerge_safety_math}
\small
\begin{tabular}{\textwidth}{llllRRR}
\toprule
\textbf{Pattern} & \textbf{Method} & \textbf{Alpha} & \textbf{Seed} & \textbf{TrustLLM Raw ASR (ASR$\downarrow$ \%)} & \textbf{Gibberish Ratio (崩壊率$\uparrow$ \%)} & \textbf{BeaverTails Utility (Score$\uparrow$ \%)} \\
\midrule
safety+math & safemerge & N/A & 42 & 40.31 & 93.44 & 7.19 \\
safety+math & safemerge & N/A & 43 & 1.25 & 0.00 & 1.56 \\
safety+math & safemerge & N/A & 44 & 1.25 & 0.00 & 1.25 \\
\bottomrule
\end{tabularx}
\end{table}


% --- 予備実験 シード別詳細 - Method: ties, Pattern: safety+math ---
\begin{table}[H]
\centering
\caption{予備実験 シード別詳細 - Method: ties, Pattern: safety+math}
\label{tab:prelim_seeds_ties_safety_math}
\small
\begin{tabular}{\textwidth}{llllRRR}
\toprule
\textbf{Pattern} & \textbf{Method} & \textbf{Alpha} & \textbf{Seed} & \textbf{TrustLLM Raw ASR (ASR$\downarrow$ \%)} & \textbf{Gibberish Ratio (崩壊率$\uparrow$ \%)} & \textbf{BeaverTails Utility (Score$\uparrow$ \%)} \\
\midrule
safety+math & ties & 0.00 & 42 & 72.50 & 30.31 & 29.38 \\
safety+math & ties & 0.00 & 43 & 72.50 & 30.31 & 29.38 \\
safety+math & ties & 0.00 & 44 & 72.50 & 30.31 & 29.38 \\
safety+math & ties & 0.20 & 42 & 20.62 & 40.94 & 11.56 \\
safety+math & ties & 0.20 & 43 & 29.06 & 39.69 & 15.62 \\
safety+math & ties & 0.20 & 44 & 25.31 & 35.62 & 16.25 \\
safety+math & ties & 0.40 & 42 & 23.12 & 43.12 & 13.44 \\
safety+math & ties & 0.40 & 43 & 28.12 & 39.69 & 15.00 \\
safety+math & ties & 0.40 & 44 & 29.06 & 36.56 & 14.37 \\
safety+math & ties & 0.60 & 42 & 13.44 & 41.88 & 11.88 \\
safety+math & ties & 0.60 & 43 & 19.69 & 32.50 & 10.00 \\
safety+math & ties & 0.60 & 44 & 22.81 & 35.62 & 11.88 \\
safety+math & ties & 0.80 & 42 & 3.75 & 1.88 & 2.19 \\
safety+math & ties & 0.80 & 43 & 10.94 & 16.25 & 5.94 \\
safety+math & ties & 0.80 & 44 & 14.69 & 40.62 & 6.56 \\
safety+math & ties & 1.00 & 42 & 0.31 & 0.00 & 1.25 \\
safety+math & ties & 1.00 & 43 & 1.88 & 0.00 & 2.50 \\
safety+math & ties & 1.00 & 44 & 0.00 & 0.00 & 0.62 \\
\bottomrule
\end{tabularx}
\end{table}


% --- 予備実験 シード別詳細 - Method: task_arithmetic, Pattern: safety+math ---
\begin{table}[H]
\centering
\caption{予備実験 シード別詳細 - Method: task_arithmetic, Pattern: safety+math}
\label{tab:prelim_seeds_taskarithmetic_safety_math}
\small
\begin{tabular}{\textwidth}{llllRRR}
\toprule
\textbf{Pattern} & \textbf{Method} & \textbf{Alpha} & \textbf{Seed} & \textbf{TrustLLM Raw ASR (ASR$\downarrow$ \%)} & \textbf{Gibberish Ratio (崩壊率$\uparrow$ \%)} & \textbf{BeaverTails Utility (Score$\uparrow$ \%)} \\
\midrule
safety+math & task\_arithmetic & 0.00 & 42 & 66.88 & 19.69 & 35.62 \\
safety+math & task\_arithmetic & 0.00 & 43 & 66.88 & 19.69 & 35.62 \\
safety+math & task\_arithmetic & 0.00 & 44 & 66.88 & 19.69 & 35.62 \\
safety+math & task\_arithmetic & 0.20 & 42 & 55.94 & 27.50 & 24.69 \\
safety+math & task\_arithmetic & 0.20 & 43 & 51.56 & 27.19 & 24.06 \\
safety+math & task\_arithmetic & 0.20 & 44 & 50.94 & 24.06 & 22.19 \\
safety+math & task\_arithmetic & 0.40 & 42 & 26.88 & 49.38 & 16.25 \\
safety+math & task\_arithmetic & 0.40 & 43 & 31.56 & 53.12 & 16.25 \\
safety+math & task\_arithmetic & 0.40 & 44 & 35.00 & 42.50 & 16.88 \\
safety+math & task\_arithmetic & 0.60 & 42 & 10.00 & 26.88 & 4.69 \\
safety+math & task\_arithmetic & 0.60 & 43 & 9.38 & 60.62 & 8.44 \\
safety+math & task\_arithmetic & 0.60 & 44 & 9.69 & 56.25 & 5.94 \\
safety+math & task\_arithmetic & 0.80 & 42 & 1.25 & 0.00 & 0.94 \\
safety+math & task\_arithmetic & 0.80 & 43 & 3.12 & 1.25 & 2.81 \\
safety+math & task\_arithmetic & 0.80 & 44 & 0.62 & 1.56 & 1.25 \\
safety+math & task\_arithmetic & 1.00 & 42 & 0.94 & 0.31 & 1.56 \\
safety+math & task\_arithmetic & 1.00 & 43 & 1.56 & 0.00 & 2.19 \\
safety+math & task\_arithmetic & 1.00 & 44 & 0.00 & 0.00 & 0.62 \\
\bottomrule
\end{tabularx}
\end{table}


% --- 予備実験 シード別詳細 - Method: dare, Pattern: safety+math ---
\begin{table}[H]
\centering
\caption{予備実験 シード別詳細 - Method: dare, Pattern: safety+math}
\label{tab:prelim_seeds_dare_safety_math}
\small
\begin{tabular}{\textwidth}{llllRRR}
\toprule
\textbf{Pattern} & \textbf{Method} & \textbf{Alpha} & \textbf{Seed} & \textbf{TrustLLM Raw ASR (ASR$\downarrow$ \%)} & \textbf{Gibberish Ratio (崩壊率$\uparrow$ \%)} & \textbf{BeaverTails Utility (Score$\uparrow$ \%)} \\
\midrule
safety+math & dare & 0.00 & 42 & 71.25 & 16.88 & 36.56 \\
safety+math & dare & 0.00 & 43 & 71.88 & 20.94 & 35.00 \\
safety+math & dare & 0.00 & 44 & 66.88 & 19.69 & 34.38 \\
safety+math & dare & 0.20 & 42 & 20.00 & 34.38 & 15.31 \\
safety+math & dare & 0.20 & 43 & 26.25 & 32.19 & 18.12 \\
safety+math & dare & 0.20 & 44 & 26.56 & 24.69 & 15.94 \\
safety+math & dare & 0.40 & 42 & 21.25 & 23.75 & 12.19 \\
safety+math & dare & 0.40 & 43 & 18.12 & 29.38 & 12.50 \\
safety+math & dare & 0.40 & 44 & 21.88 & 28.12 & 11.88 \\
safety+math & dare & 0.60 & 42 & 17.50 & 19.69 & 12.19 \\
safety+math & dare & 0.60 & 43 & 16.25 & 17.81 & 6.56 \\
safety+math & dare & 0.60 & 44 & 17.81 & 18.44 & 10.31 \\
safety+math & dare & 0.80 & 42 & 6.88 & 20.94 & 4.69 \\
safety+math & dare & 0.80 & 43 & 13.75 & 17.19 & 10.94 \\
safety+math & dare & 0.80 & 44 & 15.94 & 25.31 & 7.50 \\
safety+math & dare & 1.00 & 42 & 0.94 & 0.31 & 0.62 \\
safety+math & dare & 1.00 & 43 & 2.50 & 0.62 & 2.19 \\
safety+math & dare & 1.00 & 44 & 0.00 & 0.00 & 0.31 \\
\bottomrule
\end{tabularx}
\end{table}


% --- 予備実験 シード別詳細 - Method: mergealign, Pattern: safety+math ---
\begin{table}[H]
\centering
\caption{予備実験 シード別詳細 - Method: mergealign, Pattern: safety+math}
\label{tab:prelim_seeds_mergealign_safety_math}
\small
\begin{tabular}{\textwidth}{llllRRR}
\toprule
\textbf{Pattern} & \textbf{Method} & \textbf{Alpha} & \textbf{Seed} & \textbf{TrustLLM Raw ASR (ASR$\downarrow$ \%)} & \textbf{Gibberish Ratio (崩壊率$\uparrow$ \%)} & \textbf{BeaverTails Utility (Score$\uparrow$ \%)} \\
\midrule
safety+math & mergealign & N/A & 42 & 100.00 & 100.00 & 0.00 \\
safety+math & mergealign & N/A & 43 & 100.00 & 100.00 & 0.00 \\
safety+math & mergealign & N/A & 44 & 100.00 & 100.00 & 0.00 \\
\bottomrule
\end{tabularx}
\end{table}


% --- 予備実験 シード別詳細 - Method: diagonal_sst_main, Pattern: safety+math ---
\begin{table}[H]
\centering
\caption{予備実験 シード別詳細 - Method: diagonal_sst_main, Pattern: safety+math}
\label{tab:prelim_seeds_diagonalsstmain_safety_math}
\small
\begin{tabular}{\textwidth}{llllRRR}
\toprule
\textbf{Pattern} & \textbf{Method} & \textbf{Alpha} & \textbf{Seed} & \textbf{TrustLLM Raw ASR (ASR$\downarrow$ \%)} & \textbf{Gibberish Ratio (崩壊率$\uparrow$ \%)} & \textbf{BeaverTails Utility (Score$\uparrow$ \%)} \\
\midrule
safety+math & diagonal\_sst\_main & 0.00 & 42 & 70.62 & 17.50 & 35.00 \\
safety+math & diagonal\_sst\_main & 0.00 & 43 & 70.31 & 16.88 & 35.62 \\
safety+math & diagonal\_sst\_main & 0.00 & 44 & 70.31 & 16.88 & 35.62 \\
safety+math & diagonal\_sst\_main & 0.20 & 42 & 58.75 & 25.31 & 27.50 \\
safety+math & diagonal\_sst\_main & 0.20 & 43 & 55.31 & 20.00 & 25.31 \\
safety+math & diagonal\_sst\_main & 0.20 & 44 & 56.88 & 20.00 & 27.81 \\
safety+math & diagonal\_sst\_main & 0.40 & 42 & 42.50 & 28.44 & 20.31 \\
safety+math & diagonal\_sst\_main & 0.40 & 43 & 38.75 & 40.62 & 20.00 \\
safety+math & diagonal\_sst\_main & 0.40 & 44 & 42.50 & 28.44 & 20.31 \\
safety+math & diagonal\_sst\_main & 0.60 & 42 & 26.25 & 52.50 & 14.06 \\
safety+math & diagonal\_sst\_main & 0.60 & 43 & 22.81 & 64.38 & 14.37 \\
safety+math & diagonal\_sst\_main & 0.60 & 44 & 24.06 & 58.44 & 10.00 \\
safety+math & diagonal\_sst\_main & 0.80 & 42 & 10.62 & 61.25 & 6.25 \\
safety+math & diagonal\_sst\_main & 0.80 & 43 & 11.56 & 9.38 & 7.19 \\
safety+math & diagonal\_sst\_main & 0.80 & 44 & 11.56 & 44.38 & 6.88 \\
safety+math & diagonal\_sst\_main & 1.00 & 42 & 3.12 & 2.50 & 2.19 \\
safety+math & diagonal\_sst\_main & 1.00 & 43 & 3.44 & 0.94 & 1.88 \\
safety+math & diagonal\_sst\_main & 1.00 & 44 & 1.56 & 5.94 & 2.19 \\
\bottomrule
\end{tabularx}
\end{table}


% --- 予備実験 シード別詳細 - Method: led_merging, Pattern: safety+math ---
\begin{table}[H]
\centering
\caption{予備実験 シード別詳細 - Method: led_merging, Pattern: safety+math}
\label{tab:prelim_seeds_ledmerging_safety_math}
\small
\begin{tabular}{\textwidth}{llllRRR}
\toprule
\textbf{Pattern} & \textbf{Method} & \textbf{Alpha} & \textbf{Seed} & \textbf{TrustLLM Raw ASR (ASR$\downarrow$ \%)} & \textbf{Gibberish Ratio (崩壊率$\uparrow$ \%)} & \textbf{BeaverTails Utility (Score$\uparrow$ \%)} \\
\midrule
safety+math & led\_merging & N/A & 42 & 87.19 & 96.25 & 5.00 \\
safety+math & led\_merging & N/A & 43 & 87.19 & 96.25 & 5.00 \\
safety+math & led\_merging & N/A & 44 & 87.19 & 95.94 & 5.00 \\
\bottomrule
\end{tabularx}
\end{table}


% --- 予備実験 シード別詳細 - Method: matena_fisher, Pattern: safety+math ---
\begin{table}[H]
\centering
\caption{予備実験 シード別詳細 - Method: matena_fisher, Pattern: safety+math}
\label{tab:prelim_seeds_matenafisher_safety_math}
\small
\begin{tabular}{\textwidth}{llllRRR}
\toprule
\textbf{Pattern} & \textbf{Method} & \textbf{Alpha} & \textbf{Seed} & \textbf{TrustLLM Raw ASR (ASR$\downarrow$ \%)} & \textbf{Gibberish Ratio (崩壊率$\uparrow$ \%)} & \textbf{BeaverTails Utility (Score$\uparrow$ \%)} \\
\midrule
safety+math & matena\_fisher & N/A & 42 & 18.44 & 65.31 & 8.75 \\
safety+math & matena\_fisher & N/A & 43 & 20.31 & 71.25 & 11.88 \\
safety+math & matena\_fisher & N/A & 44 & 20.62 & 65.00 & 9.38 \\
\bottomrule
\end{tabularx}
\end{table}


% --------------------------------------------------------------------------
% 2. メイン実験 (Main Experiments)
% --------------------------------------------------------------------------

% --- メイン実験 集計 (Mean \pm Std) - Pattern: safety+math+code+medical ---
\begin{table}[H]
\centering
\caption{メイン実験 集計 (Mean \pm Std) - Pattern: safety+math+code+medical}
\label{tab:main_safety_math_code_medical}
\small
\begin{tabular}{\textwidth}{lllRRRRRRR}
\toprule
\textbf{Pattern} & \textbf{Method} & \textbf{Alpha} & \textbf{Safety Ave [Harmful Content] (ASR$\downarrow$ \%)} & \textbf{Math Ave (\%)} & \textbf{Code Ave (\%)} & \textbf{Medical Ave (\%)} & \textbf{General/Inst Ave (\%)} & \textbf{EvolCode (PPL$\downarrow$)} & \textbf{MedAlpaca (PPL$\downarrow$)} \\
\midrule
safety+math+code+medical & dare & 0.60 & 0.08 $\pm$ 0.08 & 0.00 $\pm$ 0.00 & 0.00 $\pm$ 0.00 & 38.02 $\pm$ 41.77 & 11.53 $\pm$ 0.23 & 5368837.19 $\pm$ 6493560.14 & 6550135.45 $\pm$ 3035748.48 \\
safety+math+code+medical & dare & 1.00 & 0.00 $\pm$ 0.00 & 0.68 $\pm$ 0.70 & 5.36 $\pm$ 2.13 & 55.83 $\pm$ 25.27 & 17.94 $\pm$ 12.11 & 2.35 $\pm$ 0.03 & 7.15 $\pm$ 0.10 \\
safety+math+code+medical & data\_free\_sst\_main & 0.60 & 0.08 $\pm$ 0.00 & 0.00 $\pm$ 0.00 & 0.00 $\pm$ 0.00 & 26.67 $\pm$ 24.56 & 11.85 $\pm$ 0.59 & 8840.75 $\pm$ 1530.06 & 17110.70 $\pm$ 5533.52 \\
safety+math+code+medical & data\_free\_sst\_main & 1.00 & 0.03 $\pm$ 0.05 & 0.00 $\pm$ 0.00 & 0.00 $\pm$ 0.00 & 62.40 $\pm$ 27.79 & 11.48 $\pm$ 0.20 & 6377.05 $\pm$ 1134.31 & 5459.66 $\pm$ 490.26 \\
safety+math+code+medical & della & 0.60 & 0.08 $\pm$ 0.00 & 0.00 $\pm$ 0.00 & 0.00 $\pm$ 0.00 & 15.42 $\pm$ 16.18 & 12.53 $\pm$ 0.78 & 1556853.48 $\pm$ 450599.15 & 2115706.72 $\pm$ 117978.22 \\
safety+math+code+medical & della & 1.00 & 0.00 $\pm$ 0.00 & 0.62 $\pm$ 0.31 & 4.58 $\pm$ 3.10 & 57.60 $\pm$ 23.98 & 18.90 $\pm$ 12.88 & 2.39 $\pm$ 0.04 & 7.38 $\pm$ 0.06 \\
safety+math+code+medical & diagonal\_sst\_main & 0.60 & 0.00 $\pm$ 0.00 & 0.42 $\pm$ 0.18 & 0.00 $\pm$ 0.00 & 57.29 $\pm$ 17.67 & 10.16 $\pm$ 0.49 & 319.70 $\pm$ 69.81 & 139.54 $\pm$ 27.87 \\
safety+math+code+medical & diagonal\_sst\_main & 1.00 & 0.16 $\pm$ 0.00 & 0.26 $\pm$ 0.24 & 0.00 $\pm$ 0.00 & 80.83 $\pm$ 8.57 & 11.10 $\pm$ 1.30 & 78.49 $\pm$ 39.87 & 78.49 $\pm$ 32.65 \\
safety+math+code+medical & matena\_fisher & N/A & 0.00 $\pm$ 0.00 & 0.99 $\pm$ 0.24 & 0.00 $\pm$ 0.00 & 27.81 $\pm$ 3.29 & 13.12 $\pm$ 0.12 & 63.52 $\pm$ 2.39 & 128.84 $\pm$ 4.62 \\
safety+math+code+medical & task\_arithmetic & 0.60 & 5.88 $\pm$ 0.55 & 2.55 $\pm$ 0.39 & 0.05 $\pm$ 0.09 & 65.42 $\pm$ 9.24 & 13.26 $\pm$ 1.65 & 3.22 $\pm$ 0.04 & 6.16 $\pm$ 0.11 \\
safety+math+code+medical & task\_arithmetic & 1.00 & 0.00 $\pm$ 0.00 & 0.57 $\pm$ 0.50 & 4.22 $\pm$ 3.39 & 58.12 $\pm$ 23.35 & 18.28 $\pm$ 12.09 & 2.34 $\pm$ 0.03 & 7.15 $\pm$ 0.13 \\
safety+math+code+medical & ties & 0.60 & 0.05 $\pm$ 0.05 & 0.00 $\pm$ 0.00 & 0.00 $\pm$ 0.00 & 9.17 $\pm$ 7.94 & 12.28 $\pm$ 0.40 & 233905.82 $\pm$ 257072.89 & 291983.62 $\pm$ 249888.17 \\
safety+math+code+medical & ties & 1.00 & 0.00 $\pm$ 0.00 & 0.94 $\pm$ 0.68 & 6.56 $\pm$ 2.46 & 54.58 $\pm$ 25.92 & 23.30 $\pm$ 18.12 & 2.27 $\pm$ 0.03 & 6.72 $\pm$ 0.13 \\
\bottomrule
\end{tabularx}
\end{table}


% --- メイン実験 シード別詳細 - Method: matena_fisher, Pattern: safety+math+code+medical ---
\begin{table}[H]
\centering
\caption{メイン実験 シード別詳細 - Method: matena_fisher, Pattern: safety+math+code+medical}
\label{tab:main_seeds_matenafisher_safety_math_code_medical}
\small
\begin{tabular}{\textwidth}{llllRRRRRRR}
\toprule
\textbf{Pattern} & \textbf{Method} & \textbf{Alpha} & \textbf{Seed} & \textbf{Safety Ave [Harmful Content] (ASR$\downarrow$ \%)} & \textbf{Math Ave (\%)} & \textbf{Code Ave (\%)} & \textbf{Medical Ave (\%)} & \textbf{General/Inst Ave (\%)} & \textbf{EvolCode (PPL$\downarrow$)} & \textbf{MedAlpaca (PPL$\downarrow$)} \\
\midrule
safety+math+code+medical & matena\_fisher & N/A & 42.0 & 0.00 & 0.94 & 0.00 & 30.94 & 12.98 & 65.01 & 124.79 \\
safety+math+code+medical & matena\_fisher & N/A & 43.0 & 0.00 & 0.78 & 0.00 & 28.12 & 13.15 & 60.76 & 127.88 \\
safety+math+code+medical & matena\_fisher & N/A & 44.0 & 0.00 & 1.25 & 0.00 & 24.38 & 13.22 & 64.78 & 133.87 \\
\bottomrule
\end{tabularx}
\end{table}


% --- メイン実験 シード別詳細 - Method: diagonal_sst_main, Pattern: safety+math+code+medical ---
\begin{table}[H]
\centering
\caption{メイン実験 シード別詳細 - Method: diagonal_sst_main, Pattern: safety+math+code+medical}
\label{tab:main_seeds_diagonalsstmain_safety_math_code_medical}
\small
\begin{tabular}{\textwidth}{llllRRRRRRR}
\toprule
\textbf{Pattern} & \textbf{Method} & \textbf{Alpha} & \textbf{Seed} & \textbf{Safety Ave [Harmful Content] (ASR$\downarrow$ \%)} & \textbf{Math Ave (\%)} & \textbf{Code Ave (\%)} & \textbf{Medical Ave (\%)} & \textbf{General/Inst Ave (\%)} & \textbf{EvolCode (PPL$\downarrow$)} & \textbf{MedAlpaca (PPL$\downarrow$)} \\
\midrule
safety+math+code+medical & diagonal\_sst\_main & 0.60 & 42.0 & 0.00 & 0.62 & 0.00 & 40.00 & 10.41 & 392.09 & 168.34 \\
safety+math+code+medical & diagonal\_sst\_main & 0.60 & 43.0 & 0.00 & 0.31 & 0.00 & 75.31 & 10.48 & 252.79 & 112.69 \\
safety+math+code+medical & diagonal\_sst\_main & 0.60 & 44.0 & 0.00 & 0.31 & 0.00 & 56.56 & 9.59 & 314.23 & 137.59 \\
safety+math+code+medical & diagonal\_sst\_main & 1.00 & 42.0 & 0.16 & 0.47 & 0.00 & 85.94 & 12.28 & 119.75 & 108.05 \\
safety+math+code+medical & diagonal\_sst\_main & 1.00 & 43.0 & 0.16 & 0.31 & 0.00 & 85.62 & 11.31 & 40.18 & 43.44 \\
safety+math+code+medical & diagonal\_sst\_main & 1.00 & 44.0 & 0.16 & 0.00 & 0.00 & 70.94 & 9.71 & 75.56 & 83.97 \\
\bottomrule
\end{tabularx}
\end{table}


% --- メイン実験 シード別詳細 - Method: dare, Pattern: safety+math+code+medical ---
\begin{table}[H]
\centering
\caption{メイン実験 シード別詳細 - Method: dare, Pattern: safety+math+code+medical}
\label{tab:main_seeds_dare_safety_math_code_medical}
\small
\begin{tabular}{\textwidth}{llllRRRRRRR}
\toprule
\textbf{Pattern} & \textbf{Method} & \textbf{Alpha} & \textbf{Seed} & \textbf{Safety Ave [Harmful Content] (ASR$\downarrow$ \%)} & \textbf{Math Ave (\%)} & \textbf{Code Ave (\%)} & \textbf{Medical Ave (\%)} & \textbf{General/Inst Ave (\%)} & \textbf{EvolCode (PPL$\downarrow$)} & \textbf{MedAlpaca (PPL$\downarrow$)} \\
\midrule
safety+math+code+medical & dare & 0.60 & 42.0 & 0.16 & 0.00 & 0.00 & 86.25 & 11.42 & 12857479.52 & 10037855.54 \\
safety+math+code+medical & dare & 0.60 & 43.0 & 0.08 & 0.00 & 0.00 & 14.06 & 11.38 & 1950859.88 & 4501931.51 \\
safety+math+code+medical & dare & 0.60 & 44.0 & 0.00 & 0.00 & 0.00 & 13.75 & 11.79 & 1298172.17 & 5110619.29 \\
safety+math+code+medical & dare & 1.00 & 42.0 & 0.00 & 1.41 & 7.81 & 84.06 & 31.92 & 2.31 & 7.26 \\
safety+math+code+medical & dare & 1.00 & 43.0 & 0.00 & 0.62 & 3.91 & 48.12 & 10.60 & 2.36 & 7.08 \\
safety+math+code+medical & dare & 1.00 & 44.0 & 0.00 & 0.00 & 4.38 & 35.31 & 11.30 & 2.37 & 7.10 \\
\bottomrule
\end{tabularx}
\end{table}


% --- メイン実験 シード別詳細 - Method: della, Pattern: safety+math+code+medical ---
\begin{table}[H]
\centering
\caption{メイン実験 シード別詳細 - Method: della, Pattern: safety+math+code+medical}
\label{tab:main_seeds_della_safety_math_code_medical}
\small
\begin{tabular}{\textwidth}{llllRRRRRRR}
\toprule
\textbf{Pattern} & \textbf{Method} & \textbf{Alpha} & \textbf{Seed} & \textbf{Safety Ave [Harmful Content] (ASR$\downarrow$ \%)} & \textbf{Math Ave (\%)} & \textbf{Code Ave (\%)} & \textbf{Medical Ave (\%)} & \textbf{General/Inst Ave (\%)} & \textbf{EvolCode (PPL$\downarrow$)} & \textbf{MedAlpaca (PPL$\downarrow$)} \\
\midrule
safety+math+code+medical & della & 0.60 & 42.0 & 0.08 & 0.00 & 0.00 & 13.44 & 11.64 & 1547222.80 & 1979585.04 \\
safety+math+code+medical & della & 0.60 & 43.0 & 0.08 & 0.00 & 0.00 & 0.31 & 12.97 & 1111146.86 & 2188460.85 \\
safety+math+code+medical & della & 0.60 & 44.0 & 0.08 & 0.00 & 0.00 & 32.50 & 13.00 & 2012190.77 & 2179074.27 \\
safety+math+code+medical & della & 1.00 & 42.0 & 0.00 & 0.94 & 8.12 & 84.69 & 33.77 & 2.34 & 7.44 \\
safety+math+code+medical & della & 1.00 & 43.0 & 0.00 & 0.62 & 2.34 & 49.06 & 11.21 & 2.41 & 7.34 \\
safety+math+code+medical & della & 1.00 & 44.0 & 0.00 & 0.31 & 3.28 & 39.06 & 11.72 & 2.42 & 7.35 \\
\bottomrule
\end{tabularx}
\end{table}


% --- メイン実験 シード別詳細 - Method: data_free_sst_main, Pattern: safety+math+code+medical ---
\begin{table}[H]
\centering
\caption{メイン実験 シード別詳細 - Method: data_free_sst_main, Pattern: safety+math+code+medical}
\label{tab:main_seeds_datafreesstmain_safety_math_code_medical}
\small
\begin{tabular}{\textwidth}{llllRRRRRRR}
\toprule
\textbf{Pattern} & \textbf{Method} & \textbf{Alpha} & \textbf{Seed} & \textbf{Safety Ave [Harmful Content] (ASR$\downarrow$ \%)} & \textbf{Math Ave (\%)} & \textbf{Code Ave (\%)} & \textbf{Medical Ave (\%)} & \textbf{General/Inst Ave (\%)} & \textbf{EvolCode (PPL$\downarrow$)} & \textbf{MedAlpaca (PPL$\downarrow$)} \\
\midrule
safety+math+code+medical & data\_free\_sst\_main & 0.60 & 42.0 & 0.08 & 0.00 & 0.00 & 13.44 & 12.53 & 8906.95 & 19263.61 \\
safety+math+code+medical & data\_free\_sst\_main & 0.60 & 43.0 & 0.08 & 0.00 & 0.00 & 11.56 & 11.50 & 7278.66 & 10824.29 \\
safety+math+code+medical & data\_free\_sst\_main & 0.60 & 44.0 & 0.08 & 0.00 & 0.00 & 55.00 & 11.51 & 10336.63 & 21244.18 \\
safety+math+code+medical & data\_free\_sst\_main & 1.00 & 42.0 & 0.08 & 0.00 & 0.00 & 31.87 & 11.30 & 7255.67 & 4971.49 \\
safety+math+code+medical & data\_free\_sst\_main & 1.00 & 43.0 & 0.00 & 0.00 & 0.00 & 69.06 & 11.69 & 6778.98 & 5455.51 \\
safety+math+code+medical & data\_free\_sst\_main & 1.00 & 44.0 & 0.00 & 0.00 & 0.00 & 86.25 & 11.44 & 5096.50 & 5951.99 \\
\bottomrule
\end{tabularx}
\end{table}


% --- メイン実験 シード別詳細 - Method: ties, Pattern: safety+math+code+medical ---
\begin{table}[H]
\centering
\caption{メイン実験 シード別詳細 - Method: ties, Pattern: safety+math+code+medical}
\label{tab:main_seeds_ties_safety_math_code_medical}
\small
\begin{tabular}{\textwidth}{llllRRRRRRR}
\toprule
\textbf{Pattern} & \textbf{Method} & \textbf{Alpha} & \textbf{Seed} & \textbf{Safety Ave [Harmful Content] (ASR$\downarrow$ \%)} & \textbf{Math Ave (\%)} & \textbf{Code Ave (\%)} & \textbf{Medical Ave (\%)} & \textbf{General/Inst Ave (\%)} & \textbf{EvolCode (PPL$\downarrow$)} & \textbf{MedAlpaca (PPL$\downarrow$)} \\
\midrule
safety+math+code+medical & ties & 0.60 & 42.0 & 0.08 & 0.00 & 0.00 & 13.75 & 12.22 & 104380.52 & 193731.31 \\
safety+math+code+medical & ties & 0.60 & 43.0 & 0.08 & 0.00 & 0.00 & 0.00 & 11.91 & 529977.51 & 576065.00 \\
safety+math+code+medical & ties & 0.60 & 44.0 & 0.00 & 0.00 & 0.00 & 13.75 & 12.71 & 67359.44 & 106154.54 \\
safety+math+code+medical & ties & 1.00 & 42.0 & 0.00 & 1.72 & 8.75 & 84.38 & 44.16 & 2.24 & 6.85 \\
safety+math+code+medical & ties & 1.00 & 43.0 & 0.00 & 0.62 & 3.91 & 37.19 & 11.49 & 2.28 & 6.71 \\
safety+math+code+medical & ties & 1.00 & 44.0 & 0.00 & 0.47 & 7.03 & 42.19 & 14.25 & 2.29 & 6.60 \\
\bottomrule
\end{tabularx}
\end{table}


% --- メイン実験 シード別詳細 - Method: task_arithmetic, Pattern: safety+math+code+medical ---
\begin{table}[H]
\centering
\caption{メイン実験 シード別詳細 - Method: task_arithmetic, Pattern: safety+math+code+medical}
\label{tab:main_seeds_taskarithmetic_safety_math_code_medical}
\small
\begin{tabular}{\textwidth}{llllRRRRRRR}
\toprule
\textbf{Pattern} & \textbf{Method} & \textbf{Alpha} & \textbf{Seed} & \textbf{Safety Ave [Harmful Content] (ASR$\downarrow$ \%)} & \textbf{Math Ave (\%)} & \textbf{Code Ave (\%)} & \textbf{Medical Ave (\%)} & \textbf{General/Inst Ave (\%)} & \textbf{EvolCode (PPL$\downarrow$)} & \textbf{MedAlpaca (PPL$\downarrow$)} \\
\midrule
safety+math+code+medical & task\_arithmetic & 0.60 & 42.0 & 5.77 & 2.19 & 0.16 & 73.44 & 14.72 & 3.17 & 6.10 \\
safety+math+code+medical & task\_arithmetic & 0.60 & 43.0 & 6.48 & 2.50 & 0.00 & 67.50 & 13.59 & 3.25 & 6.28 \\
safety+math+code+medical & task\_arithmetic & 0.60 & 44.0 & 5.40 & 2.97 & 0.00 & 55.31 & 11.48 & 3.24 & 6.10 \\
safety+math+code+medical & task\_arithmetic & 1.00 & 42.0 & 0.00 & 0.94 & 8.12 & 84.38 & 32.24 & 2.31 & 7.26 \\
safety+math+code+medical & task\_arithmetic & 1.00 & 43.0 & 0.00 & 0.78 & 2.03 & 50.31 & 11.23 & 2.36 & 7.18 \\
safety+math+code+medical & task\_arithmetic & 1.00 & 44.0 & 0.00 & 0.00 & 2.50 & 39.69 & 11.37 & 2.36 & 7.00 \\
\bottomrule
\end{tabularx}
\end{table}


% --- メイン実験 集計 (Mean \pm Std) - Pattern: safety+math ---
\begin{table}[H]
\centering
\caption{メイン実験 集計 (Mean \pm Std) - Pattern: safety+math}
\label{tab:main_safety_math}
\small
\begin{tabular}{\textwidth}{lllRRRRRRR}
\toprule
\textbf{Pattern} & \textbf{Method} & \textbf{Alpha} & \textbf{Safety Ave [Harmful Content] (ASR$\downarrow$ \%)} & \textbf{Math Ave (\%)} & \textbf{Code Ave (\%)} & \textbf{Medical Ave (\%)} & \textbf{General/Inst Ave (\%)} & \textbf{EvolCode (PPL$\downarrow$)} & \textbf{MedAlpaca (PPL$\downarrow$)} \\
\midrule
safety+math & dare & 0.60 & 6.14 $\pm$ 5.14 & 5.42 $\pm$ 0.79 & 5.00 $\pm$ 1.02 & 84.27 $\pm$ 2.35 & 22.85 $\pm$ 8.57 & 2.43 $\pm$ 0.07 & 4.44 $\pm$ 0.28 \\
safety+math & dare & 1.00 & 0.00 $\pm$ 0.00 & 0.57 $\pm$ 0.63 & 5.68 $\pm$ 3.61 & 55.94 $\pm$ 25.88 & 11.22 $\pm$ 0.57 & 2.35 $\pm$ 0.03 & 7.18 $\pm$ 0.22 \\
safety+math & data\_free\_sst\_main & 0.60 & 14.73 $\pm$ 3.41 & 5.57 $\pm$ 0.24 & 7.19 $\pm$ 4.33 & 69.38 $\pm$ 13.52 & 24.69 $\pm$ 0.38 & 2.02 $\pm$ 0.02 & 3.02 $\pm$ 0.02 \\
safety+math & data\_free\_sst\_main & 1.00 & 11.49 $\pm$ 9.25 & 5.89 $\pm$ 0.79 & 5.57 $\pm$ 0.24 & 56.77 $\pm$ 26.43 & 21.14 $\pm$ 7.66 & 2.19 $\pm$ 0.04 & 4.64 $\pm$ 0.10 \\
safety+math & della & 0.60 & 6.66 $\pm$ 5.52 & 5.52 $\pm$ 0.63 & 5.31 $\pm$ 1.36 & 85.62 $\pm$ 0.54 & 31.03 $\pm$ 4.06 & 2.46 $\pm$ 0.06 & 4.50 $\pm$ 0.30 \\
safety+math & della & 1.00 & 0.00 $\pm$ 0.00 & 0.62 $\pm$ 0.41 & 4.69 $\pm$ 3.29 & 58.54 $\pm$ 23.27 & 17.78 $\pm$ 11.72 & 2.38 $\pm$ 0.03 & 7.33 $\pm$ 0.14 \\
safety+math & diagonal\_sst\_main & 0.60 & 17.09 $\pm$ 5.01 & 6.82 $\pm$ 0.86 & 6.72 $\pm$ 0.95 & 72.55 $\pm$ 15.91 & 22.48 $\pm$ 8.92 & 1.84 $\pm$ 0.01 & 2.75 $\pm$ 0.05 \\
safety+math & diagonal\_sst\_main & 1.00 & 0.00 $\pm$ 0.00 & 9.84 $\pm$ 0.47 & 7.76 $\pm$ 3.63 & 65.42 $\pm$ 31.57 & 20.65 $\pm$ 6.27 & 1.98 $\pm$ 0.07 & 4.36 $\pm$ 0.49 \\
safety+math & led\_merging & N/A & 2.80 $\pm$ 0.00 & 1.41 $\pm$ 0.00 & 5.62 $\pm$ 0.00 & 83.75 $\pm$ 0.00 & 22.51 $\pm$ 0.02 & 1.75 $\pm$ 0.00 & 2.72 $\pm$ 0.00 \\
safety+math & matena\_fisher & N/A & 4.29 $\pm$ 0.43 & 6.04 $\pm$ 0.86 & 6.61 $\pm$ 0.39 & 85.73 $\pm$ 0.65 & 26.72 $\pm$ 11.70 & 1.82 $\pm$ 0.01 & 2.88 $\pm$ 0.02 \\
safety+math & mergealign & N/A & 0.00 $\pm$ 0.00 & 0.00 $\pm$ 0.00 & 0.00 $\pm$ 0.00 & 86.25 $\pm$ 0.00 & 15.73 $\pm$ 3.41 & N/A & N/A \\
safety+math & safemerge & N/A & 4.61 $\pm$ 7.99 & 1.30 $\pm$ 0.86 & 6.67 $\pm$ 4.86 & 82.50 $\pm$ 1.74 & 13.09 $\pm$ 1.77 & 2.22 $\pm$ 0.48 & 6.41 $\pm$ 3.54 \\
safety+math & task\_arithmetic & 0.60 & 2.48 $\pm$ 2.89 & 6.88 $\pm$ 1.56 & 7.50 $\pm$ 0.68 & 75.21 $\pm$ 15.93 & 27.18 $\pm$ 6.92 & 1.88 $\pm$ 0.02 & 3.47 $\pm$ 0.12 \\
safety+math & task\_arithmetic & 1.00 & 0.00 $\pm$ 0.00 & 0.62 $\pm$ 0.54 & 6.93 $\pm$ 8.08 & 48.80 $\pm$ 8.77 & 18.36 $\pm$ 12.02 & 2.34 $\pm$ 0.03 & 7.15 $\pm$ 0.14 \\
safety+math & ties & 0.60 & 9.03 $\pm$ 5.55 & 5.94 $\pm$ 0.83 & 6.15 $\pm$ 0.18 & 85.00 $\pm$ 2.44 & 27.65 $\pm$ 3.07 & 2.15 $\pm$ 0.05 & 3.57 $\pm$ 0.10 \\
safety+math & ties & 1.00 & 0.00 $\pm$ 0.00 & 0.94 $\pm$ 0.68 & 6.56 $\pm$ 2.46 & 54.58 $\pm$ 25.92 & 25.52 $\pm$ 16.78 & 2.27 $\pm$ 0.03 & 6.72 $\pm$ 0.13 \\
\bottomrule
\end{tabularx}
\end{table}


% --- メイン実験 シード別詳細 - Method: della, Pattern: safety+math ---
\begin{table}[H]
\centering
\caption{メイン実験 シード別詳細 - Method: della, Pattern: safety+math}
\label{tab:main_seeds_della_safety_math}
\small
\begin{tabular}{\textwidth}{llllRRRRRRR}
\toprule
\textbf{Pattern} & \textbf{Method} & \textbf{Alpha} & \textbf{Seed} & \textbf{Safety Ave [Harmful Content] (ASR$\downarrow$ \%)} & \textbf{Math Ave (\%)} & \textbf{Code Ave (\%)} & \textbf{Medical Ave (\%)} & \textbf{General/Inst Ave (\%)} & \textbf{EvolCode (PPL$\downarrow$)} & \textbf{MedAlpaca (PPL$\downarrow$)} \\
\midrule
safety+math & della & 0.60 & 42.0 & 0.63 & 4.84 & 6.25 & 85.31 & 29.00 & 2.39 & 4.26 \\
safety+math & della & 0.60 & 43.0 & 11.45 & 5.62 & 5.94 & 85.31 & 28.39 & 2.47 & 4.41 \\
safety+math & della & 0.60 & 44.0 & 7.92 & 6.09 & 3.75 & 86.25 & 35.70 & 2.51 & 4.84 \\
safety+math & della & 1.00 & 42.0 & 0.00 & 0.78 & 7.81 & 85.00 & 31.29 & 2.35 & 7.45 \\
safety+math & della & 1.00 & 43.0 & 0.00 & 0.94 & 1.25 & 49.38 & 10.38 & 2.41 & 7.37 \\
safety+math & della & 1.00 & 44.0 & 0.00 & 0.16 & 5.00 & 41.25 & 11.68 & 2.40 & 7.17 \\
\bottomrule
\end{tabularx}
\end{table}


% --- メイン実験 シード別詳細 - Method: matena_fisher, Pattern: safety+math ---
\begin{table}[H]
\centering
\caption{メイン実験 シード別詳細 - Method: matena_fisher, Pattern: safety+math}
\label{tab:main_seeds_matenafisher_safety_math}
\small
\begin{tabular}{\textwidth}{llllRRRRRRR}
\toprule
\textbf{Pattern} & \textbf{Method} & \textbf{Alpha} & \textbf{Seed} & \textbf{Safety Ave [Harmful Content] (ASR$\downarrow$ \%)} & \textbf{Math Ave (\%)} & \textbf{Code Ave (\%)} & \textbf{Medical Ave (\%)} & \textbf{General/Inst Ave (\%)} & \textbf{EvolCode (PPL$\downarrow$)} & \textbf{MedAlpaca (PPL$\downarrow$)} \\
\midrule
safety+math & matena\_fisher & N/A & 42.0 & 4.71 & 6.09 & 6.25 & 85.00 & 39.12 & 1.81 & 2.90 \\
safety+math & matena\_fisher & N/A & 43.0 & 4.30 & 5.16 & 7.03 & 86.25 & 15.87 & 1.83 & 2.86 \\
safety+math & matena\_fisher & N/A & 44.0 & 3.85 & 6.88 & 6.56 & 85.94 & 25.17 & 1.84 & 2.88 \\
\bottomrule
\end{tabularx}
\end{table}


% --- メイン実験 シード別詳細 - Method: safemerge, Pattern: safety+math ---
\begin{table}[H]
\centering
\caption{メイン実験 シード別詳細 - Method: safemerge, Pattern: safety+math}
\label{tab:main_seeds_safemerge_safety_math}
\small
\begin{tabular}{\textwidth}{llllRRRRRRR}
\toprule
\textbf{Pattern} & \textbf{Method} & \textbf{Alpha} & \textbf{Seed} & \textbf{Safety Ave [Harmful Content] (ASR$\downarrow$ \%)} & \textbf{Math Ave (\%)} & \textbf{Code Ave (\%)} & \textbf{Medical Ave (\%)} & \textbf{General/Inst Ave (\%)} & \textbf{EvolCode (PPL$\downarrow$)} & \textbf{MedAlpaca (PPL$\downarrow$)} \\
\midrule
safety+math & safemerge & N/A & 42.0 & 13.84 & 2.19 & 8.91 & 80.62 & 15.06 & 1.71 & 2.51 \\
safety+math & safemerge & N/A & 43.0 & 0.00 & 0.47 & 1.09 & 84.06 & 11.63 & 2.66 & 9.43 \\
safety+math & safemerge & N/A & 44.0 & 0.00 & 1.25 & 10.00 & 82.81 & 12.60 & 2.30 & 7.29 \\
\bottomrule
\end{tabularx}
\end{table}


% --- メイン実験 シード別詳細 - Method: dare, Pattern: safety+math ---
\begin{table}[H]
\centering
\caption{メイン実験 シード別詳細 - Method: dare, Pattern: safety+math}
\label{tab:main_seeds_dare_safety_math}
\small
\begin{tabular}{\textwidth}{llllRRRRRRR}
\toprule
\textbf{Pattern} & \textbf{Method} & \textbf{Alpha} & \textbf{Seed} & \textbf{Safety Ave [Harmful Content] (ASR$\downarrow$ \%)} & \textbf{Math Ave (\%)} & \textbf{Code Ave (\%)} & \textbf{Medical Ave (\%)} & \textbf{General/Inst Ave (\%)} & \textbf{EvolCode (PPL$\downarrow$)} & \textbf{MedAlpaca (PPL$\downarrow$)} \\
\midrule
safety+math & dare & 0.60 & 42.0 & 0.32 & 6.25 & 5.16 & 84.38 & 15.04 & 2.35 & 4.18 \\
safety+math & dare & 0.60 & 43.0 & 8.06 & 4.69 & 5.94 & 81.88 & 21.51 & 2.46 & 4.40 \\
safety+math & dare & 0.60 & 44.0 & 10.04 & 5.31 & 3.91 & 86.56 & 32.02 & 2.48 & 4.73 \\
safety+math & dare & 1.00 & 42.0 & 0.00 & 1.25 & 9.06 & 85.62 & 11.49 & 2.32 & 7.37 \\
safety+math & dare & 1.00 & 43.0 & 0.00 & 0.47 & 1.88 & 44.06 & 10.56 & 2.37 & 7.23 \\
safety+math & dare & 1.00 & 44.0 & 0.00 & 0.00 & 6.09 & 38.12 & 11.61 & 2.36 & 6.93 \\
\bottomrule
\end{tabularx}
\end{table}


% --- メイン実験 シード別詳細 - Method: data_free_sst_main, Pattern: safety+math ---
\begin{table}[H]
\centering
\caption{メイン実験 シード別詳細 - Method: data_free_sst_main, Pattern: safety+math}
\label{tab:main_seeds_datafreesstmain_safety_math}
\small
\begin{tabular}{\textwidth}{llllRRRRRRR}
\toprule
\textbf{Pattern} & \textbf{Method} & \textbf{Alpha} & \textbf{Seed} & \textbf{Safety Ave [Harmful Content] (ASR$\downarrow$ \%)} & \textbf{Math Ave (\%)} & \textbf{Code Ave (\%)} & \textbf{Medical Ave (\%)} & \textbf{General/Inst Ave (\%)} & \textbf{EvolCode (PPL$\downarrow$)} & \textbf{MedAlpaca (PPL$\downarrow$)} \\
\midrule
safety+math & data\_free\_sst\_main & 0.60 & 42.0 & 10.89 & 5.31 & 12.19 & 54.06 & 24.30 & 1.99 & 3.03 \\
safety+math & data\_free\_sst\_main & 0.60 & 43.0 & 15.92 & 5.78 & 4.84 & 74.38 & 25.07 & 2.02 & 3.04 \\
safety+math & data\_free\_sst\_main & 0.60 & 44.0 & 17.39 & 5.62 & 4.53 & 79.69 & 24.70 & 2.04 & 3.00 \\
safety+math & data\_free\_sst\_main & 1.00 & 42.0 & 0.95 & 5.78 & 5.31 & 28.75 & 13.26 & 2.14 & 4.52 \\
safety+math & data\_free\_sst\_main & 1.00 & 43.0 & 18.26 & 6.72 & 5.78 & 60.31 & 28.55 & 2.23 & 4.71 \\
safety+math & data\_free\_sst\_main & 1.00 & 44.0 & 15.26 & 5.16 & 5.62 & 81.25 & 21.63 & 2.21 & 4.69 \\
\bottomrule
\end{tabularx}
\end{table}


% --- メイン実験 シード別詳細 - Method: ties, Pattern: safety+math ---
\begin{table}[H]
\centering
\caption{メイン実験 シード別詳細 - Method: ties, Pattern: safety+math}
\label{tab:main_seeds_ties_safety_math}
\small
\begin{tabular}{\textwidth}{llllRRRRRRR}
\toprule
\textbf{Pattern} & \textbf{Method} & \textbf{Alpha} & \textbf{Seed} & \textbf{Safety Ave [Harmful Content] (ASR$\downarrow$ \%)} & \textbf{Math Ave (\%)} & \textbf{Code Ave (\%)} & \textbf{Medical Ave (\%)} & \textbf{General/Inst Ave (\%)} & \textbf{EvolCode (PPL$\downarrow$)} & \textbf{MedAlpaca (PPL$\downarrow$)} \\
\midrule
safety+math & ties & 0.60 & 42.0 & 3.28 & 6.88 & 6.25 & 86.25 & 24.12 & 2.10 & 3.46 \\
safety+math & ties & 0.60 & 43.0 & 9.45 & 5.62 & 6.25 & 82.19 & 29.17 & 2.16 & 3.61 \\
safety+math & ties & 0.60 & 44.0 & 14.36 & 5.31 & 5.94 & 86.56 & 29.67 & 2.19 & 3.64 \\
safety+math & ties & 1.00 & 42.0 & 0.00 & 1.72 & 8.75 & 84.38 & 44.14 & 2.24 & 6.85 \\
safety+math & ties & 1.00 & 43.0 & 0.00 & 0.62 & 3.91 & 37.19 & 11.60 & 2.28 & 6.71 \\
safety+math & ties & 1.00 & 44.0 & 0.00 & 0.47 & 7.03 & 42.19 & 20.81 & 2.29 & 6.60 \\
\bottomrule
\end{tabularx}
\end{table}


% --- メイン実験 シード別詳細 - Method: led_merging, Pattern: safety+math ---
\begin{table}[H]
\centering
\caption{メイン実験 シード別詳細 - Method: led_merging, Pattern: safety+math}
\label{tab:main_seeds_ledmerging_safety_math}
\small
\begin{tabular}{\textwidth}{llllRRRRRRR}
\toprule
\textbf{Pattern} & \textbf{Method} & \textbf{Alpha} & \textbf{Seed} & \textbf{Safety Ave [Harmful Content] (ASR$\downarrow$ \%)} & \textbf{Math Ave (\%)} & \textbf{Code Ave (\%)} & \textbf{Medical Ave (\%)} & \textbf{General/Inst Ave (\%)} & \textbf{EvolCode (PPL$\downarrow$)} & \textbf{MedAlpaca (PPL$\downarrow$)} \\
\midrule
safety+math & led\_merging & N/A & 42.0 & 2.80 & 1.41 & 5.62 & 83.75 & 22.50 & 1.75 & 2.72 \\
safety+math & led\_merging & N/A & 43.0 & 2.80 & 1.41 & 5.62 & 83.75 & 22.52 & 1.75 & 2.72 \\
safety+math & led\_merging & N/A & 44.0 & 2.80 & 1.41 & 5.62 & 83.75 & 22.49 & 1.75 & 2.72 \\
\bottomrule
\end{tabularx}
\end{table}


% --- メイン実験 シード別詳細 - Method: mergealign, Pattern: safety+math ---
\begin{table}[H]
\centering
\caption{メイン実験 シード別詳細 - Method: mergealign, Pattern: safety+math}
\label{tab:main_seeds_mergealign_safety_math}
\small
\begin{tabular}{\textwidth}{llllRRRRRRR}
\toprule
\textbf{Pattern} & \textbf{Method} & \textbf{Alpha} & \textbf{Seed} & \textbf{Safety Ave [Harmful Content] (ASR$\downarrow$ \%)} & \textbf{Math Ave (\%)} & \textbf{Code Ave (\%)} & \textbf{Medical Ave (\%)} & \textbf{General/Inst Ave (\%)} & \textbf{EvolCode (PPL$\downarrow$)} & \textbf{MedAlpaca (PPL$\downarrow$)} \\
\midrule
safety+math & mergealign & N/A & 42.0 & 0.00 & 0.00 & 0.00 & 86.25 & 17.71 & N/A & N/A \\
safety+math & mergealign & N/A & 43.0 & 0.00 & 0.00 & 0.00 & 86.25 & 11.79 & N/A & N/A \\
safety+math & mergealign & N/A & 44.0 & 0.00 & 0.00 & 0.00 & 86.25 & 17.71 & N/A & N/A \\
\bottomrule
\end{tabularx}
\end{table}


% --- メイン実験 シード別詳細 - Method: task_arithmetic, Pattern: safety+math ---
\begin{table}[H]
\centering
\caption{メイン実験 シード別詳細 - Method: task_arithmetic, Pattern: safety+math}
\label{tab:main_seeds_taskarithmetic_safety_math}
\small
\begin{tabular}{\textwidth}{llllRRRRRRR}
\toprule
\textbf{Pattern} & \textbf{Method} & \textbf{Alpha} & \textbf{Seed} & \textbf{Safety Ave [Harmful Content] (ASR$\downarrow$ \%)} & \textbf{Math Ave (\%)} & \textbf{Code Ave (\%)} & \textbf{Medical Ave (\%)} & \textbf{General/Inst Ave (\%)} & \textbf{EvolCode (PPL$\downarrow$)} & \textbf{MedAlpaca (PPL$\downarrow$)} \\
\midrule
safety+math & task\_arithmetic & 0.60 & 42.0 & 0.16 & 8.44 & 6.72 & 56.88 & 35.05 & 1.86 & 3.48 \\
safety+math & task\_arithmetic & 0.60 & 43.0 & 1.57 & 5.31 & 7.97 & 83.12 & 22.05 & 1.89 & 3.34 \\
safety+math & task\_arithmetic & 0.60 & 44.0 & 5.72 & 6.88 & 7.81 & 85.62 & 24.45 & 1.89 & 3.57 \\
safety+math & task\_arithmetic & 1.00 & 42.0 & 0.00 & 0.94 & 16.25 & 56.72 & 32.24 & 2.31 & 7.26 \\
safety+math & task\_arithmetic & 1.00 & 43.0 & 0.00 & 0.94 & 2.03 & 50.31 & 11.23 & 2.36 & 7.18 \\
safety+math & task\_arithmetic & 1.00 & 44.0 & 0.00 & 0.00 & 2.50 & 39.38 & 11.62 & 2.36 & 7.00 \\
\bottomrule
\end{tabularx}
\end{table}


% --- メイン実験 シード別詳細 - Method: diagonal_sst_main, Pattern: safety+math ---
\begin{table}[H]
\centering
\caption{メイン実験 シード別詳細 - Method: diagonal_sst_main, Pattern: safety+math}
\label{tab:main_seeds_diagonalsstmain_safety_math}
\small
\begin{tabular}{\textwidth}{llllRRRRRRR}
\toprule
\textbf{Pattern} & \textbf{Method} & \textbf{Alpha} & \textbf{Seed} & \textbf{Safety Ave [Harmful Content] (ASR$\downarrow$ \%)} & \textbf{Math Ave (\%)} & \textbf{Code Ave (\%)} & \textbf{Medical Ave (\%)} & \textbf{General/Inst Ave (\%)} & \textbf{EvolCode (PPL$\downarrow$)} & \textbf{MedAlpaca (PPL$\downarrow$)} \\
\midrule
safety+math & diagonal\_sst\_main & 0.60 & 42.0 & 22.60 & 6.25 & 7.34 & 54.22 & 13.40 & 1.83 & 2.69 \\
safety+math & diagonal\_sst\_main & 0.60 & 43.0 & 15.85 & 6.41 & 7.19 & 80.62 & 31.23 & 1.84 & 2.79 \\
safety+math & diagonal\_sst\_main & 0.60 & 44.0 & 12.81 & 7.81 & 5.62 & 82.81 & 22.82 & 1.85 & 2.78 \\
safety+math & diagonal\_sst\_main & 1.00 & 42.0 & 0.00 & 9.38 & 11.88 & 29.06 & 13.45 & 1.90 & 3.82 \\
safety+math & diagonal\_sst\_main & 1.00 & 43.0 & 0.00 & 10.31 & 6.41 & 81.25 & 24.94 & 2.03 & 4.79 \\
safety+math & diagonal\_sst\_main & 1.00 & 44.0 & 0.00 & 9.84 & 5.00 & 85.94 & 23.56 & 2.00 & 4.46 \\
\bottomrule
\end{tabularx}
\end{table}


% --- メイン実験 集計 (Mean \pm Std) - Pattern: safety+code ---
\begin{table}[H]
\centering
\caption{メイン実験 集計 (Mean \pm Std) - Pattern: safety+code}
\label{tab:main_safety_code}
\small
\begin{tabular}{\textwidth}{lllRRRRRRR}
\toprule
\textbf{Pattern} & \textbf{Method} & \textbf{Alpha} & \textbf{Safety Ave [Harmful Content] (ASR$\downarrow$ \%)} & \textbf{Math Ave (\%)} & \textbf{Code Ave (\%)} & \textbf{Medical Ave (\%)} & \textbf{General/Inst Ave (\%)} & \textbf{EvolCode (PPL$\downarrow$)} & \textbf{MedAlpaca (PPL$\downarrow$)} \\
\midrule
safety+code & dare & 0.60 & 0.00 $\pm$ 0.00 & 0.05 $\pm$ 0.09 & 0.00 $\pm$ 0.00 & 29.58 $\pm$ 12.27 & 11.73 $\pm$ 0.70 & 202107.17 $\pm$ 205144.66 & 343746.27 $\pm$ 163477.65 \\
safety+code & dare & 1.00 & 0.00 $\pm$ 0.00 & 0.73 $\pm$ 0.50 & 5.21 $\pm$ 2.67 & 57.81 $\pm$ 22.81 & 21.54 $\pm$ 10.74 & 2.35 $\pm$ 0.04 & 7.07 $\pm$ 0.07 \\
safety+code & data\_free\_sst\_main & 0.60 & 0.03 $\pm$ 0.05 & 0.21 $\pm$ 0.24 & 0.00 $\pm$ 0.00 & 78.75 $\pm$ 2.48 & 14.07 $\pm$ 0.05 & 218.68 $\pm$ 118.96 & 49.90 $\pm$ 16.95 \\
safety+code & data\_free\_sst\_main & 1.00 & 0.03 $\pm$ 0.05 & 0.05 $\pm$ 0.09 & 0.00 $\pm$ 0.00 & 28.85 $\pm$ 16.70 & 11.52 $\pm$ 0.25 & 17917.37 $\pm$ 18265.91 & 2247.19 $\pm$ 1603.28 \\
safety+code & della & 0.60 & 0.00 $\pm$ 0.00 & 0.62 $\pm$ 0.41 & 0.00 $\pm$ 0.00 & 43.54 $\pm$ 14.08 & 12.43 $\pm$ 0.46 & 107944.55 $\pm$ 59575.64 & 211408.90 $\pm$ 76715.67 \\
safety+code & della & 1.00 & 0.00 $\pm$ 0.00 & 0.62 $\pm$ 0.72 & 4.43 $\pm$ 3.60 & 53.54 $\pm$ 27.30 & 11.04 $\pm$ 0.62 & 2.38 $\pm$ 0.04 & 7.38 $\pm$ 0.07 \\
safety+code & diagonal\_sst\_main & 0.60 & 0.97 $\pm$ 0.42 & 0.89 $\pm$ 0.24 & 0.00 $\pm$ 0.00 & 86.25 $\pm$ 0.31 & 12.26 $\pm$ 0.55 & 191.50 $\pm$ 74.10 & 31.40 $\pm$ 8.25 \\
safety+code & diagonal\_sst\_main & 1.00 & 0.00 $\pm$ 0.00 & 0.42 $\pm$ 0.24 & 0.00 $\pm$ 0.00 & 15.21 $\pm$ 0.95 & 10.63 $\pm$ 0.27 & 187.93 $\pm$ 30.20 & 249.26 $\pm$ 144.36 \\
safety+code & led\_merging & N/A & 2.88 $\pm$ 0.14 & 1.41 $\pm$ 0.00 & 5.62 $\pm$ 0.00 & 83.75 $\pm$ 0.00 & 22.49 $\pm$ 0.04 & 1.75 $\pm$ 0.00 & 2.72 $\pm$ 0.00 \\
safety+code & matena\_fisher & N/A & 3.37 $\pm$ 1.40 & 1.09 $\pm$ 0.00 & 0.00 $\pm$ 0.00 & 83.12 $\pm$ 0.31 & 12.77 $\pm$ 0.37 & 69.53 $\pm$ 21.34 & 84.66 $\pm$ 27.77 \\
safety+code & mergealign & N/A & 0.00 $\pm$ 0.00 & 0.00 $\pm$ 0.00 & 0.00 $\pm$ 0.00 & 86.25 $\pm$ 0.00 & 15.73 $\pm$ 3.41 & N/A & N/A \\
safety+code & safemerge & N/A & 0.10 $\pm$ 0.18 & 1.51 $\pm$ 1.33 & 5.99 $\pm$ 5.39 & 73.75 $\pm$ 2.98 & 19.84 $\pm$ 8.02 & 2.35 $\pm$ 0.78 & 7.03 $\pm$ 4.96 \\
safety+code & task\_arithmetic & 0.60 & 29.17 $\pm$ 3.07 & 2.97 $\pm$ 0.68 & 4.63 $\pm$ 0.09 & 88.85 $\pm$ 0.79 & 17.39 $\pm$ 0.39 & 2.47 $\pm$ 0.01 & 4.84 $\pm$ 0.03 \\
safety+code & task\_arithmetic & 1.00 & 0.00 $\pm$ 0.00 & 0.57 $\pm$ 0.50 & 4.22 $\pm$ 3.39 & 58.02 $\pm$ 23.47 & 18.37 $\pm$ 12.03 & 2.34 $\pm$ 0.03 & 7.15 $\pm$ 0.13 \\
safety+code & ties & 0.60 & 0.11 $\pm$ 0.04 & 1.41 $\pm$ 0.47 & 0.00 $\pm$ 0.00 & 38.65 $\pm$ 26.32 & 17.85 $\pm$ 9.81 & 92.27 $\pm$ 111.77 & 25.62 $\pm$ 27.16 \\
safety+code & ties & 1.00 & 0.00 $\pm$ 0.00 & 0.94 $\pm$ 0.68 & 6.56 $\pm$ 2.46 & 54.58 $\pm$ 25.92 & 25.53 $\pm$ 16.91 & 2.27 $\pm$ 0.03 & 6.72 $\pm$ 0.13 \\
\bottomrule
\end{tabularx}
\end{table}


% --- メイン実験 シード別詳細 - Method: task_arithmetic, Pattern: safety+code ---
\begin{table}[H]
\centering
\caption{メイン実験 シード別詳細 - Method: task_arithmetic, Pattern: safety+code}
\label{tab:main_seeds_taskarithmetic_safety_code}
\small
\begin{tabular}{\textwidth}{llllRRRRRRR}
\toprule
\textbf{Pattern} & \textbf{Method} & \textbf{Alpha} & \textbf{Seed} & \textbf{Safety Ave [Harmful Content] (ASR$\downarrow$ \%)} & \textbf{Math Ave (\%)} & \textbf{Code Ave (\%)} & \textbf{Medical Ave (\%)} & \textbf{General/Inst Ave (\%)} & \textbf{EvolCode (PPL$\downarrow$)} & \textbf{MedAlpaca (PPL$\downarrow$)} \\
\midrule
safety+code & task\_arithmetic & 0.60 & 42.0 & 32.56 & 3.28 & 4.53 & 89.69 & 17.60 & 2.47 & 4.80 \\
safety+code & task\_arithmetic & 0.60 & 43.0 & 26.56 & 3.44 & 4.68 & 88.12 & 17.62 & 2.47 & 4.84 \\
safety+code & task\_arithmetic & 0.60 & 44.0 & 28.39 & 2.19 & 4.68 & 88.75 & 16.93 & 2.48 & 4.87 \\
safety+code & task\_arithmetic & 1.00 & 42.0 & 0.00 & 0.94 & 8.12 & 84.38 & 32.26 & 2.31 & 7.26 \\
safety+code & task\_arithmetic & 1.00 & 43.0 & 0.00 & 0.78 & 2.03 & 50.31 & 11.23 & 2.36 & 7.18 \\
safety+code & task\_arithmetic & 1.00 & 44.0 & 0.00 & 0.00 & 2.50 & 39.38 & 11.62 & 2.36 & 7.00 \\
\bottomrule
\end{tabularx}
\end{table}


% --- メイン実験 シード別詳細 - Method: dare, Pattern: safety+code ---
\begin{table}[H]
\centering
\caption{メイン実験 シード別詳細 - Method: dare, Pattern: safety+code}
\label{tab:main_seeds_dare_safety_code}
\small
\begin{tabular}{\textwidth}{llllRRRRRRR}
\toprule
\textbf{Pattern} & \textbf{Method} & \textbf{Alpha} & \textbf{Seed} & \textbf{Safety Ave [Harmful Content] (ASR$\downarrow$ \%)} & \textbf{Math Ave (\%)} & \textbf{Code Ave (\%)} & \textbf{Medical Ave (\%)} & \textbf{General/Inst Ave (\%)} & \textbf{EvolCode (PPL$\downarrow$)} & \textbf{MedAlpaca (PPL$\downarrow$)} \\
\midrule
safety+code & dare & 0.60 & 42.0 & 0.00 & 0.00 & 0.00 & 31.25 & 11.10 & 84728.55 & 334537.08 \\
safety+code & dare & 0.60 & 43.0 & 0.00 & 0.16 & 0.00 & 16.56 & 11.60 & 82608.30 & 185067.88 \\
safety+code & dare & 0.60 & 44.0 & 0.00 & 0.00 & 0.00 & 40.94 & 12.48 & 438984.66 & 511633.87 \\
safety+code & dare & 1.00 & 42.0 & 0.00 & 0.94 & 8.28 & 84.06 & 32.23 & 2.30 & 7.08 \\
safety+code & dare & 1.00 & 43.0 & 0.00 & 1.09 & 3.44 & 42.81 & 10.74 & 2.37 & 7.13 \\
safety+code & dare & 1.00 & 44.0 & 0.00 & 0.16 & 3.91 & 46.56 & 21.65 & 2.37 & 6.99 \\
\bottomrule
\end{tabularx}
\end{table}


% --- メイン実験 シード別詳細 - Method: matena_fisher, Pattern: safety+code ---
\begin{table}[H]
\centering
\caption{メイン実験 シード別詳細 - Method: matena_fisher, Pattern: safety+code}
\label{tab:main_seeds_matenafisher_safety_code}
\small
\begin{tabular}{\textwidth}{llllRRRRRRR}
\toprule
\textbf{Pattern} & \textbf{Method} & \textbf{Alpha} & \textbf{Seed} & \textbf{Safety Ave [Harmful Content] (ASR$\downarrow$ \%)} & \textbf{Math Ave (\%)} & \textbf{Code Ave (\%)} & \textbf{Medical Ave (\%)} & \textbf{General/Inst Ave (\%)} & \textbf{EvolCode (PPL$\downarrow$)} & \textbf{MedAlpaca (PPL$\downarrow$)} \\
\midrule
safety+code & matena\_fisher & N/A & 42.0 & 4.96 & 1.09 & 0.00 & 83.44 & 12.87 & 47.93 & 54.86 \\
safety+code & matena\_fisher & N/A & 43.0 & 2.85 & 1.09 & 0.00 & 82.81 & 13.07 & 70.07 & 89.29 \\
safety+code & matena\_fisher & N/A & 44.0 & 2.31 & 1.09 & 0.00 & 83.12 & 12.36 & 90.59 & 109.82 \\
\bottomrule
\end{tabularx}
\end{table}


% --- メイン実験 シード別詳細 - Method: diagonal_sst_main, Pattern: safety+code ---
\begin{table}[H]
\centering
\caption{メイン実験 シード別詳細 - Method: diagonal_sst_main, Pattern: safety+code}
\label{tab:main_seeds_diagonalsstmain_safety_code}
\small
\begin{tabular}{\textwidth}{llllRRRRRRR}
\toprule
\textbf{Pattern} & \textbf{Method} & \textbf{Alpha} & \textbf{Seed} & \textbf{Safety Ave [Harmful Content] (ASR$\downarrow$ \%)} & \textbf{Math Ave (\%)} & \textbf{Code Ave (\%)} & \textbf{Medical Ave (\%)} & \textbf{General/Inst Ave (\%)} & \textbf{EvolCode (PPL$\downarrow$)} & \textbf{MedAlpaca (PPL$\downarrow$)} \\
\midrule
safety+code & diagonal\_sst\_main & 0.60 & 42.0 & 1.21 & 1.09 & 0.00 & 86.25 & 12.88 & 135.96 & 21.94 \\
safety+code & diagonal\_sst\_main & 0.60 & 43.0 & 1.20 & 0.94 & 0.00 & 85.94 & 11.85 & 275.64 & 35.11 \\
safety+code & diagonal\_sst\_main & 0.60 & 44.0 & 0.49 & 0.62 & 0.00 & 86.56 & 12.05 & 162.89 & 37.14 \\
safety+code & diagonal\_sst\_main & 1.00 & 42.0 & 0.00 & 0.47 & 0.00 & 16.25 & 10.94 & 221.80 & 409.75 \\
safety+code & diagonal\_sst\_main & 1.00 & 43.0 & 0.00 & 0.62 & 0.00 & 14.37 & 10.47 & 178.18 & 130.00 \\
safety+code & diagonal\_sst\_main & 1.00 & 44.0 & 0.00 & 0.16 & 0.00 & 15.00 & 10.49 & 163.80 & 208.02 \\
\bottomrule
\end{tabularx}
\end{table}


% --- メイン実験 シード別詳細 - Method: safemerge, Pattern: safety+code ---
\begin{table}[H]
\centering
\caption{メイン実験 シード別詳細 - Method: safemerge, Pattern: safety+code}
\label{tab:main_seeds_safemerge_safety_code}
\small
\begin{tabular}{\textwidth}{llllRRRRRRR}
\toprule
\textbf{Pattern} & \textbf{Method} & \textbf{Alpha} & \textbf{Seed} & \textbf{Safety Ave [Harmful Content] (ASR$\downarrow$ \%)} & \textbf{Math Ave (\%)} & \textbf{Code Ave (\%)} & \textbf{Medical Ave (\%)} & \textbf{General/Inst Ave (\%)} & \textbf{EvolCode (PPL$\downarrow$)} & \textbf{MedAlpaca (PPL$\downarrow$)} \\
\midrule
safety+code & safemerge & N/A & 42.0 & 0.31 & 2.50 & 7.03 & 75.62 & 27.76 & 1.75 & 2.87 \\
safety+code & safemerge & N/A & 43.0 & 0.00 & 2.03 & 10.78 & 70.31 & 20.03 & 2.06 & 5.72 \\
safety+code & safemerge & N/A & 44.0 & 0.00 & 0.00 & 0.16 & 75.31 & 11.72 & 3.23 & 12.51 \\
\bottomrule
\end{tabularx}
\end{table}


% --- メイン実験 シード別詳細 - Method: led_merging, Pattern: safety+code ---
\begin{table}[H]
\centering
\caption{メイン実験 シード別詳細 - Method: led_merging, Pattern: safety+code}
\label{tab:main_seeds_ledmerging_safety_code}
\small
\begin{tabular}{\textwidth}{llllRRRRRRR}
\toprule
\textbf{Pattern} & \textbf{Method} & \textbf{Alpha} & \textbf{Seed} & \textbf{Safety Ave [Harmful Content] (ASR$\downarrow$ \%)} & \textbf{Math Ave (\%)} & \textbf{Code Ave (\%)} & \textbf{Medical Ave (\%)} & \textbf{General/Inst Ave (\%)} & \textbf{EvolCode (PPL$\downarrow$)} & \textbf{MedAlpaca (PPL$\downarrow$)} \\
\midrule
safety+code & led\_merging & N/A & 42.0 & 3.04 & 1.41 & 5.62 & 83.75 & 22.49 & 1.75 & 2.72 \\
safety+code & led\_merging & N/A & 43.0 & 2.80 & 1.41 & 5.62 & 83.75 & 22.52 & 1.75 & 2.72 \\
safety+code & led\_merging & N/A & 44.0 & 2.80 & 1.41 & 5.62 & 83.75 & 22.45 & 1.75 & 2.72 \\
\bottomrule
\end{tabularx}
\end{table}


% --- メイン実験 シード別詳細 - Method: della, Pattern: safety+code ---
\begin{table}[H]
\centering
\caption{メイン実験 シード別詳細 - Method: della, Pattern: safety+code}
\label{tab:main_seeds_della_safety_code}
\small
\begin{tabular}{\textwidth}{llllRRRRRRR}
\toprule
\textbf{Pattern} & \textbf{Method} & \textbf{Alpha} & \textbf{Seed} & \textbf{Safety Ave [Harmful Content] (ASR$\downarrow$ \%)} & \textbf{Math Ave (\%)} & \textbf{Code Ave (\%)} & \textbf{Medical Ave (\%)} & \textbf{General/Inst Ave (\%)} & \textbf{EvolCode (PPL$\downarrow$)} & \textbf{MedAlpaca (PPL$\downarrow$)} \\
\midrule
safety+code & della & 0.60 & 42.0 & 0.00 & 0.94 & 0.00 & 57.19 & 12.96 & 173731.57 & 126277.55 \\
safety+code & della & 0.60 & 43.0 & 0.00 & 0.16 & 0.00 & 29.06 & 12.10 & 92466.80 & 275182.68 \\
safety+code & della & 0.60 & 44.0 & 0.00 & 0.78 & 0.00 & 44.38 & 12.23 & 57635.29 & 232766.47 \\
safety+code & della & 1.00 & 42.0 & 0.00 & 1.41 & 8.12 & 84.69 & 11.16 & 2.34 & 7.41 \\
safety+code & della & 1.00 & 43.0 & 0.00 & 0.47 & 0.94 & 42.19 & 10.37 & 2.40 & 7.43 \\
safety+code & della & 1.00 & 44.0 & 0.00 & 0.00 & 4.22 & 33.75 & 11.59 & 2.40 & 7.29 \\
\bottomrule
\end{tabularx}
\end{table}


% --- メイン実験 シード別詳細 - Method: mergealign, Pattern: safety+code ---
\begin{table}[H]
\centering
\caption{メイン実験 シード別詳細 - Method: mergealign, Pattern: safety+code}
\label{tab:main_seeds_mergealign_safety_code}
\small
\begin{tabular}{\textwidth}{llllRRRRRRR}
\toprule
\textbf{Pattern} & \textbf{Method} & \textbf{Alpha} & \textbf{Seed} & \textbf{Safety Ave [Harmful Content] (ASR$\downarrow$ \%)} & \textbf{Math Ave (\%)} & \textbf{Code Ave (\%)} & \textbf{Medical Ave (\%)} & \textbf{General/Inst Ave (\%)} & \textbf{EvolCode (PPL$\downarrow$)} & \textbf{MedAlpaca (PPL$\downarrow$)} \\
\midrule
safety+code & mergealign & N/A & 42.0 & 0.00 & 0.00 & 0.00 & 86.25 & 17.71 & N/A & N/A \\
safety+code & mergealign & N/A & 43.0 & 0.00 & 0.00 & 0.00 & 86.25 & 11.79 & N/A & N/A \\
safety+code & mergealign & N/A & 44.0 & 0.00 & 0.00 & 0.00 & 86.25 & 17.71 & N/A & N/A \\
\bottomrule
\end{tabularx}
\end{table}


% --- メイン実験 シード別詳細 - Method: ties, Pattern: safety+code ---
\begin{table}[H]
\centering
\caption{メイン実験 シード別詳細 - Method: ties, Pattern: safety+code}
\label{tab:main_seeds_ties_safety_code}
\small
\begin{tabular}{\textwidth}{llllRRRRRRR}
\toprule
\textbf{Pattern} & \textbf{Method} & \textbf{Alpha} & \textbf{Seed} & \textbf{Safety Ave [Harmful Content] (ASR$\downarrow$ \%)} & \textbf{Math Ave (\%)} & \textbf{Code Ave (\%)} & \textbf{Medical Ave (\%)} & \textbf{General/Inst Ave (\%)} & \textbf{EvolCode (PPL$\downarrow$)} & \textbf{MedAlpaca (PPL$\downarrow$)} \\
\midrule
safety+code & ties & 0.60 & 42.0 & 0.08 & 1.41 & 0.00 & 68.75 & 12.90 & 50.93 & 56.91 \\
safety+code & ties & 0.60 & 43.0 & 0.16 & 1.88 & 0.00 & 20.00 & 11.50 & 7.06 & 8.18 \\
safety+code & ties & 0.60 & 44.0 & 0.08 & 0.94 & 0.00 & 27.19 & 29.15 & 218.81 & 11.76 \\
safety+code & ties & 1.00 & 42.0 & 0.00 & 1.72 & 8.75 & 84.38 & 44.31 & 2.24 & 6.85 \\
safety+code & ties & 1.00 & 43.0 & 0.00 & 0.62 & 3.91 & 37.19 & 11.49 & 2.28 & 6.71 \\
safety+code & ties & 1.00 & 44.0 & 0.00 & 0.47 & 7.03 & 42.19 & 20.80 & 2.29 & 6.60 \\
\bottomrule
\end{tabularx}
\end{table}


% --- メイン実験 シード別詳細 - Method: data_free_sst_main, Pattern: safety+code ---
\begin{table}[H]
\centering
\caption{メイン実験 シード別詳細 - Method: data_free_sst_main, Pattern: safety+code}
\label{tab:main_seeds_datafreesstmain_safety_code}
\small
\begin{tabular}{\textwidth}{llllRRRRRRR}
\toprule
\textbf{Pattern} & \textbf{Method} & \textbf{Alpha} & \textbf{Seed} & \textbf{Safety Ave [Harmful Content] (ASR$\downarrow$ \%)} & \textbf{Math Ave (\%)} & \textbf{Code Ave (\%)} & \textbf{Medical Ave (\%)} & \textbf{General/Inst Ave (\%)} & \textbf{EvolCode (PPL$\downarrow$)} & \textbf{MedAlpaca (PPL$\downarrow$)} \\
\midrule
safety+code & data\_free\_sst\_main & 0.60 & 42.0 & 0.00 & 0.16 & 0.00 & 80.62 & 14.01 & 242.36 & 51.52 \\
safety+code & data\_free\_sst\_main & 0.60 & 43.0 & 0.00 & 0.47 & 0.00 & 79.69 & 14.10 & 89.66 & 32.20 \\
safety+code & data\_free\_sst\_main & 0.60 & 44.0 & 0.08 & 0.00 & 0.00 & 75.94 & 14.09 & 324.02 & 65.98 \\
safety+code & data\_free\_sst\_main & 1.00 & 42.0 & 0.00 & 0.00 & 0.00 & 48.12 & 11.80 & 13963.92 & 2780.62 \\
safety+code & data\_free\_sst\_main & 1.00 & 43.0 & 0.00 & 0.16 & 0.00 & 19.69 & 11.40 & 1951.94 & 445.19 \\
safety+code & data\_free\_sst\_main & 1.00 & 44.0 & 0.08 & 0.00 & 0.00 & 18.75 & 11.35 & 37836.26 & 3515.76 \\
\bottomrule
\end{tabularx}
\end{table}


% --- メイン実験 集計 (Mean \pm Std) - Pattern: safety+medical ---
\begin{table}[H]
\centering
\caption{メイン実験 集計 (Mean \pm Std) - Pattern: safety+medical}
\label{tab:main_safety_medical}
\small
\begin{tabular}{\textwidth}{lllRRRRRRR}
\toprule
\textbf{Pattern} & \textbf{Method} & \textbf{Alpha} & \textbf{Safety Ave [Harmful Content] (ASR$\downarrow$ \%)} & \textbf{Math Ave (\%)} & \textbf{Code Ave (\%)} & \textbf{Medical Ave (\%)} & \textbf{General/Inst Ave (\%)} & \textbf{EvolCode (PPL$\downarrow$)} & \textbf{MedAlpaca (PPL$\downarrow$)} \\
\midrule
safety+medical & dare & 0.60 & 0.03 $\pm$ 0.05 & 0.00 $\pm$ 0.00 & 0.00 $\pm$ 0.00 & 22.08 $\pm$ 10.81 & 12.39 $\pm$ 0.44 & 12863372.97 $\pm$ 9716067.34 & 13423154.39 $\pm$ 11060518.87 \\
safety+medical & dare & 1.00 & 0.00 $\pm$ 0.00 & 0.47 $\pm$ 0.56 & 4.38 $\pm$ 3.53 & 62.40 $\pm$ 23.88 & 23.44 $\pm$ 10.58 & 2.34 $\pm$ 0.04 & 7.13 $\pm$ 0.18 \\
safety+medical & data\_free\_sst\_main & 0.60 & 0.05 $\pm$ 0.05 & 0.00 $\pm$ 0.00 & 0.00 $\pm$ 0.00 & 61.67 $\pm$ 41.50 & 12.27 $\pm$ 0.07 & 151079.31 $\pm$ 79294.36 & 109978.81 $\pm$ 39827.45 \\
safety+medical & data\_free\_sst\_main & 1.00 & 0.08 $\pm$ 0.08 & 0.00 $\pm$ 0.00 & 0.00 $\pm$ 0.00 & 32.92 $\pm$ 46.61 & 11.96 $\pm$ 0.22 & 216810.71 $\pm$ 115193.54 & 136969.65 $\pm$ 56251.67 \\
safety+medical & della & 0.60 & 0.03 $\pm$ 0.05 & 0.00 $\pm$ 0.00 & 0.00 $\pm$ 0.00 & 27.29 $\pm$ 23.45 & 12.10 $\pm$ 1.01 & 10754022.73 $\pm$ 9470200.24 & 8631379.06 $\pm$ 4139779.72 \\
safety+medical & della & 1.00 & 0.00 $\pm$ 0.00 & 0.62 $\pm$ 0.68 & 4.01 $\pm$ 3.43 & 56.56 $\pm$ 22.05 & 17.08 $\pm$ 10.75 & 2.38 $\pm$ 0.04 & 7.38 $\pm$ 0.14 \\
safety+medical & diagonal\_sst\_main & 0.60 & 0.05 $\pm$ 0.05 & 0.00 $\pm$ 0.00 & 0.00 $\pm$ 0.00 & 9.69 $\pm$ 8.46 & 12.06 $\pm$ 0.13 & 75380.55 $\pm$ 8006.67 & 78424.43 $\pm$ 9712.01 \\
safety+medical & diagonal\_sst\_main & 1.00 & 0.05 $\pm$ 0.05 & 0.00 $\pm$ 0.00 & 0.00 $\pm$ 0.00 & 26.25 $\pm$ 22.89 & 11.87 $\pm$ 0.98 & 138717.43 $\pm$ 145875.51 & 95092.19 $\pm$ 65045.99 \\
safety+medical & matena\_fisher & N/A & 0.16 $\pm$ 0.00 & 0.00 $\pm$ 0.00 & 0.00 $\pm$ 0.00 & 0.00 $\pm$ 0.00 & 12.99 $\pm$ 0.18 & 576870.77 $\pm$ 100574.92 & 439254.45 $\pm$ 88250.04 \\
safety+medical & mergealign & N/A & 0.00 $\pm$ 0.00 & 0.00 $\pm$ 0.00 & 0.00 $\pm$ 0.00 & 86.25 $\pm$ 0.00 & 13.76 $\pm$ 3.41 & N/A & N/A \\
safety+medical & safemerge & N/A & 0.03 $\pm$ 0.05 & 1.09 $\pm$ 1.89 & 3.91 $\pm$ 6.77 & 53.44 $\pm$ 10.84 & 12.81 $\pm$ 1.63 & 3.84 $\pm$ 1.84 & 15.66 $\pm$ 10.55 \\
safety+medical & task\_arithmetic & 0.60 & 0.16 $\pm$ 0.00 & 0.00 $\pm$ 0.00 & 0.00 $\pm$ 0.00 & 13.85 $\pm$ 0.18 & 12.72 $\pm$ 0.06 & 217753.69 $\pm$ 12313.53 & 198661.81 $\pm$ 12470.98 \\
safety+medical & task\_arithmetic & 1.00 & 0.00 $\pm$ 0.00 & 0.57 $\pm$ 0.50 & 4.22 $\pm$ 3.39 & 58.02 $\pm$ 23.47 & 18.37 $\pm$ 12.03 & 2.34 $\pm$ 0.03 & 7.15 $\pm$ 0.13 \\
safety+medical & ties & 0.60 & 0.05 $\pm$ 0.05 & 0.00 $\pm$ 0.00 & 0.00 $\pm$ 0.00 & 13.54 $\pm$ 0.36 & 11.90 $\pm$ 0.62 & 34172.15 $\pm$ 4640.21 & 67899.54 $\pm$ 37744.22 \\
safety+medical & ties & 1.00 & 0.00 $\pm$ 0.00 & 0.94 $\pm$ 0.68 & 6.56 $\pm$ 2.46 & 54.58 $\pm$ 25.92 & 20.66 $\pm$ 9.09 & 2.27 $\pm$ 0.03 & 6.72 $\pm$ 0.13 \\
\bottomrule
\end{tabularx}
\end{table}


% --- メイン実験 シード別詳細 - Method: safemerge, Pattern: safety+medical ---
\begin{table}[H]
\centering
\caption{メイン実験 シード別詳細 - Method: safemerge, Pattern: safety+medical}
\label{tab:main_seeds_safemerge_safety_medical}
\small
\begin{tabular}{\textwidth}{llllRRRRRRR}
\toprule
\textbf{Pattern} & \textbf{Method} & \textbf{Alpha} & \textbf{Seed} & \textbf{Safety Ave [Harmful Content] (ASR$\downarrow$ \%)} & \textbf{Math Ave (\%)} & \textbf{Code Ave (\%)} & \textbf{Medical Ave (\%)} & \textbf{General/Inst Ave (\%)} & \textbf{EvolCode (PPL$\downarrow$)} & \textbf{MedAlpaca (PPL$\downarrow$)} \\
\midrule
safety+medical & safemerge & N/A & 42.0 & 0.00 & 3.28 & 11.72 & 60.31 & 14.69 & 1.87 & 4.37 \\
safety+medical & safemerge & N/A & 43.0 & 0.08 & 0.00 & 0.00 & 40.94 & 11.81 & 5.53 & 25.25 \\
safety+medical & safemerge & N/A & 44.0 & 0.00 & 0.00 & 0.00 & 59.06 & 11.93 & 4.11 & 17.37 \\
\bottomrule
\end{tabularx}
\end{table}


% --- メイン実験 シード別詳細 - Method: task_arithmetic, Pattern: safety+medical ---
\begin{table}[H]
\centering
\caption{メイン実験 シード別詳細 - Method: task_arithmetic, Pattern: safety+medical}
\label{tab:main_seeds_taskarithmetic_safety_medical}
\small
\begin{tabular}{\textwidth}{llllRRRRRRR}
\toprule
\textbf{Pattern} & \textbf{Method} & \textbf{Alpha} & \textbf{Seed} & \textbf{Safety Ave [Harmful Content] (ASR$\downarrow$ \%)} & \textbf{Math Ave (\%)} & \textbf{Code Ave (\%)} & \textbf{Medical Ave (\%)} & \textbf{General/Inst Ave (\%)} & \textbf{EvolCode (PPL$\downarrow$)} & \textbf{MedAlpaca (PPL$\downarrow$)} \\
\midrule
safety+medical & task\_arithmetic & 0.60 & 42.0 & 0.16 & 0.00 & 0.00 & 13.75 & 12.69 & 231670.29 & 211313.97 \\
safety+medical & task\_arithmetic & 0.60 & 43.0 & 0.16 & 0.00 & 0.00 & 14.06 & 12.69 & 213319.17 & 186380.28 \\
safety+medical & task\_arithmetic & 0.60 & 44.0 & 0.16 & 0.00 & 0.00 & 13.75 & 12.79 & 208271.62 & 198291.16 \\
safety+medical & task\_arithmetic & 1.00 & 42.0 & 0.00 & 0.94 & 8.12 & 84.38 & 32.26 & 2.31 & 7.26 \\
safety+medical & task\_arithmetic & 1.00 & 43.0 & 0.00 & 0.78 & 2.03 & 50.31 & 11.23 & 2.36 & 7.18 \\
safety+medical & task\_arithmetic & 1.00 & 44.0 & 0.00 & 0.00 & 2.50 & 39.38 & 11.62 & 2.36 & 7.00 \\
\bottomrule
\end{tabularx}
\end{table}


% --- メイン実験 シード別詳細 - Method: dare, Pattern: safety+medical ---
\begin{table}[H]
\centering
\caption{メイン実験 シード別詳細 - Method: dare, Pattern: safety+medical}
\label{tab:main_seeds_dare_safety_medical}
\small
\begin{tabular}{\textwidth}{llllRRRRRRR}
\toprule
\textbf{Pattern} & \textbf{Method} & \textbf{Alpha} & \textbf{Seed} & \textbf{Safety Ave [Harmful Content] (ASR$\downarrow$ \%)} & \textbf{Math Ave (\%)} & \textbf{Code Ave (\%)} & \textbf{Medical Ave (\%)} & \textbf{General/Inst Ave (\%)} & \textbf{EvolCode (PPL$\downarrow$)} & \textbf{MedAlpaca (PPL$\downarrow$)} \\
\midrule
safety+medical & dare & 0.60 & 42.0 & 0.00 & 0.00 & 0.00 & 17.81 & 12.81 & 23273464.09 & 25783754.01 \\
safety+medical & dare & 0.60 & 43.0 & 0.08 & 0.00 & 0.00 & 14.06 & 11.93 & 4035556.82 & 10026154.84 \\
safety+medical & dare & 0.60 & 44.0 & 0.00 & 0.00 & 0.00 & 34.38 & 12.42 & 11281098.00 & 4459554.31 \\
safety+medical & dare & 1.00 & 42.0 & 0.00 & 1.09 & 8.44 & 84.69 & 30.19 & 2.30 & 7.33 \\
safety+medical & dare & 1.00 & 43.0 & 0.00 & 0.31 & 2.66 & 65.31 & 28.88 & 2.36 & 7.10 \\
safety+medical & dare & 1.00 & 44.0 & 0.00 & 0.00 & 2.03 & 37.19 & 11.24 & 2.37 & 6.97 \\
\bottomrule
\end{tabularx}
\end{table}


% --- メイン実験 シード別詳細 - Method: matena_fisher, Pattern: safety+medical ---
\begin{table}[H]
\centering
\caption{メイン実験 シード別詳細 - Method: matena_fisher, Pattern: safety+medical}
\label{tab:main_seeds_matenafisher_safety_medical}
\small
\begin{tabular}{\textwidth}{llllRRRRRRR}
\toprule
\textbf{Pattern} & \textbf{Method} & \textbf{Alpha} & \textbf{Seed} & \textbf{Safety Ave [Harmful Content] (ASR$\downarrow$ \%)} & \textbf{Math Ave (\%)} & \textbf{Code Ave (\%)} & \textbf{Medical Ave (\%)} & \textbf{General/Inst Ave (\%)} & \textbf{EvolCode (PPL$\downarrow$)} & \textbf{MedAlpaca (PPL$\downarrow$)} \\
\midrule
safety+medical & matena\_fisher & N/A & 42.0 & 0.16 & 0.00 & 0.00 & 0.00 & 13.09 & 691154.34 & 541088.28 \\
safety+medical & matena\_fisher & N/A & 43.0 & 0.16 & 0.00 & 0.00 & 0.00 & 12.78 & 501847.06 & 385101.33 \\
safety+medical & matena\_fisher & N/A & 44.0 & 0.16 & 0.00 & 0.00 & 0.00 & 13.09 & 537610.90 & 391573.73 \\
\bottomrule
\end{tabularx}
\end{table}


% --- メイン実験 シード別詳細 - Method: ties, Pattern: safety+medical ---
\begin{table}[H]
\centering
\caption{メイン実験 シード別詳細 - Method: ties, Pattern: safety+medical}
\label{tab:main_seeds_ties_safety_medical}
\small
\begin{tabular}{\textwidth}{llllRRRRRRR}
\toprule
\textbf{Pattern} & \textbf{Method} & \textbf{Alpha} & \textbf{Seed} & \textbf{Safety Ave [Harmful Content] (ASR$\downarrow$ \%)} & \textbf{Math Ave (\%)} & \textbf{Code Ave (\%)} & \textbf{Medical Ave (\%)} & \textbf{General/Inst Ave (\%)} & \textbf{EvolCode (PPL$\downarrow$)} & \textbf{MedAlpaca (PPL$\downarrow$)} \\
\midrule
safety+medical & ties & 0.60 & 42.0 & 0.08 & 0.00 & 0.00 & 13.75 & 12.33 & 38946.69 & 110844.44 \\
safety+medical & ties & 0.60 & 43.0 & 0.00 & 0.00 & 0.00 & 13.12 & 12.17 & 33890.70 & 39990.65 \\
safety+medical & ties & 0.60 & 44.0 & 0.08 & 0.00 & 0.00 & 13.75 & 11.19 & 29679.08 & 52863.53 \\
safety+medical & ties & 1.00 & 42.0 & 0.00 & 1.72 & 8.75 & 84.38 & 29.68 & 2.24 & 6.85 \\
safety+medical & ties & 1.00 & 43.0 & 0.00 & 0.62 & 3.91 & 37.19 & 11.49 & 2.28 & 6.71 \\
safety+medical & ties & 1.00 & 44.0 & 0.00 & 0.47 & 7.03 & 42.19 & 20.80 & 2.29 & 6.60 \\
\bottomrule
\end{tabularx}
\end{table}


% --- メイン実験 シード別詳細 - Method: data_free_sst_main, Pattern: safety+medical ---
\begin{table}[H]
\centering
\caption{メイン実験 シード別詳細 - Method: data_free_sst_main, Pattern: safety+medical}
\label{tab:main_seeds_datafreesstmain_safety_medical}
\small
\begin{tabular}{\textwidth}{llllRRRRRRR}
\toprule
\textbf{Pattern} & \textbf{Method} & \textbf{Alpha} & \textbf{Seed} & \textbf{Safety Ave [Harmful Content] (ASR$\downarrow$ \%)} & \textbf{Math Ave (\%)} & \textbf{Code Ave (\%)} & \textbf{Medical Ave (\%)} & \textbf{General/Inst Ave (\%)} & \textbf{EvolCode (PPL$\downarrow$)} & \textbf{MedAlpaca (PPL$\downarrow$)} \\
\midrule
safety+medical & data\_free\_sst\_main & 0.60 & 42.0 & 0.00 & 0.00 & 0.00 & 85.00 & 12.19 & 204517.26 & 155544.43 \\
safety+medical & data\_free\_sst\_main & 0.60 & 43.0 & 0.08 & 0.00 & 0.00 & 86.25 & 12.31 & 59971.79 & 81805.58 \\
safety+medical & data\_free\_sst\_main & 0.60 & 44.0 & 0.08 & 0.00 & 0.00 & 13.75 & 12.31 & 188748.87 & 92586.41 \\
safety+medical & data\_free\_sst\_main & 1.00 & 42.0 & 0.08 & 0.00 & 0.00 & 86.25 & 12.08 & 236077.90 & 149930.00 \\
safety+medical & data\_free\_sst\_main & 1.00 & 43.0 & 0.00 & 0.00 & 0.00 & 0.00 & 11.70 & 321155.77 & 185610.00 \\
safety+medical & data\_free\_sst\_main & 1.00 & 44.0 & 0.16 & 0.00 & 0.00 & 12.50 & 12.09 & 93198.47 & 75368.95 \\
\bottomrule
\end{tabularx}
\end{table}


% --- メイン実験 シード別詳細 - Method: della, Pattern: safety+medical ---
\begin{table}[H]
\centering
\caption{メイン実験 シード別詳細 - Method: della, Pattern: safety+medical}
\label{tab:main_seeds_della_safety_medical}
\small
\begin{tabular}{\textwidth}{llllRRRRRRR}
\toprule
\textbf{Pattern} & \textbf{Method} & \textbf{Alpha} & \textbf{Seed} & \textbf{Safety Ave [Harmful Content] (ASR$\downarrow$ \%)} & \textbf{Math Ave (\%)} & \textbf{Code Ave (\%)} & \textbf{Medical Ave (\%)} & \textbf{General/Inst Ave (\%)} & \textbf{EvolCode (PPL$\downarrow$)} & \textbf{MedAlpaca (PPL$\downarrow$)} \\
\midrule
safety+medical & della & 0.60 & 42.0 & 0.08 & 0.00 & 0.00 & 13.75 & 12.84 & 21562404.57 & 13042229.15 \\
safety+medical & della & 0.60 & 43.0 & 0.00 & 0.00 & 0.00 & 54.37 & 12.52 & 6788182.90 & 4830310.81 \\
safety+medical & della & 0.60 & 44.0 & 0.00 & 0.00 & 0.00 & 13.75 & 10.95 & 3911480.70 & 8021597.21 \\
safety+medical & della & 1.00 & 42.0 & 0.00 & 1.41 & 7.97 & 81.88 & 29.49 & 2.33 & 7.51 \\
safety+medical & della & 1.00 & 43.0 & 0.00 & 0.31 & 2.19 & 46.25 & 10.77 & 2.39 & 7.41 \\
safety+medical & della & 1.00 & 44.0 & 0.00 & 0.16 & 1.88 & 41.56 & 10.97 & 2.41 & 7.23 \\
\bottomrule
\end{tabularx}
\end{table}


% --- メイン実験 シード別詳細 - Method: mergealign, Pattern: safety+medical ---
\begin{table}[H]
\centering
\caption{メイン実験 シード別詳細 - Method: mergealign, Pattern: safety+medical}
\label{tab:main_seeds_mergealign_safety_medical}
\small
\begin{tabular}{\textwidth}{llllRRRRRRR}
\toprule
\textbf{Pattern} & \textbf{Method} & \textbf{Alpha} & \textbf{Seed} & \textbf{Safety Ave [Harmful Content] (ASR$\downarrow$ \%)} & \textbf{Math Ave (\%)} & \textbf{Code Ave (\%)} & \textbf{Medical Ave (\%)} & \textbf{General/Inst Ave (\%)} & \textbf{EvolCode (PPL$\downarrow$)} & \textbf{MedAlpaca (PPL$\downarrow$)} \\
\midrule
safety+medical & mergealign & N/A & 42.0 & 0.00 & 0.00 & 0.00 & 86.25 & 17.71 & N/A & N/A \\
safety+medical & mergealign & N/A & 43.0 & 0.00 & 0.00 & 0.00 & 86.25 & 11.79 & N/A & N/A \\
safety+medical & mergealign & N/A & 44.0 & 0.00 & 0.00 & 0.00 & 86.25 & 11.79 & N/A & N/A \\
\bottomrule
\end{tabularx}
\end{table}


% --- メイン実験 シード別詳細 - Method: diagonal_sst_main, Pattern: safety+medical ---
\begin{table}[H]
\centering
\caption{メイン実験 シード別詳細 - Method: diagonal_sst_main, Pattern: safety+medical}
\label{tab:main_seeds_diagonalsstmain_safety_medical}
\small
\begin{tabular}{\textwidth}{llllRRRRRRR}
\toprule
\textbf{Pattern} & \textbf{Method} & \textbf{Alpha} & \textbf{Seed} & \textbf{Safety Ave [Harmful Content] (ASR$\downarrow$ \%)} & \textbf{Math Ave (\%)} & \textbf{Code Ave (\%)} & \textbf{Medical Ave (\%)} & \textbf{General/Inst Ave (\%)} & \textbf{EvolCode (PPL$\downarrow$)} & \textbf{MedAlpaca (PPL$\downarrow$)} \\
\midrule
safety+medical & diagonal\_sst\_main & 0.60 & 42.0 & 0.08 & 0.00 & 0.00 & 13.44 & 11.91 & 72783.78 & 67362.99 \\
safety+medical & diagonal\_sst\_main & 0.60 & 43.0 & 0.00 & 0.00 & 0.00 & 0.00 & 12.16 & 84363.30 & 85554.10 \\
safety+medical & diagonal\_sst\_main & 0.60 & 44.0 & 0.08 & 0.00 & 0.00 & 15.62 & 12.11 & 68994.57 & 82356.19 \\
safety+medical & diagonal\_sst\_main & 1.00 & 42.0 & 0.00 & 0.00 & 0.00 & 7.81 & 11.46 & 11982.52 & 20544.71 \\
safety+medical & diagonal\_sst\_main & 1.00 & 43.0 & 0.08 & 0.00 & 0.00 & 51.88 & 11.16 & 298175.26 & 140302.29 \\
safety+medical & diagonal\_sst\_main & 1.00 & 44.0 & 0.08 & 0.00 & 0.00 & 19.06 & 13.00 & 105994.50 & 124429.57 \\
\bottomrule
\end{tabularx}
\end{table}


% --------------------------------------------------------------------------
% 3. ベースモデル (Base Models)
% --------------------------------------------------------------------------

% --- ベースモデル評価 集計 (Mean \pm Std) ---
\begin{table}[H]
\centering
\caption{ベースモデル評価 集計 (Mean \pm Std)}
\label{tab:base_models}
\small
\begin{tabular}{\textwidth}{lllRRRRRRR}
\toprule
\textbf{Pattern} & \textbf{Method} & \textbf{Alpha} & \textbf{Safety Ave [Harmful Content] (ASR$\downarrow$ \%)} & \textbf{Math Ave (\%)} & \textbf{Code Ave (\%)} & \textbf{Medical Ave (\%)} & \textbf{General/Inst Ave (\%)} & \textbf{EvolCode (PPL$\downarrow$)} & \textbf{MedAlpaca (PPL$\downarrow$)} \\
\midrule
Base (MedAlpaca) & Base (MedAlpaca) & N/A & 5.17 $\pm$ 0.00 & 4.53 $\pm$ 0.00 & 3.59 $\pm$ 0.00 & 52.81 $\pm$ 0.00 & 20.28 $\pm$ 0.00 & 2.32 $\pm$ 0.00 & 5.71 $\pm$ 2.88 \\
Base (SafetyFT) & Base (SafetyFT) & N/A & 0.00 $\pm$ 0.00 & 0.52 $\pm$ 0.48 & 4.22 $\pm$ 3.39 & 43.49 $\pm$ 11.80 & 11.38 $\pm$ 0.41 & 2.34 $\pm$ 0.03 & N/A \\
Base (WizardCoder) & Base (WizardCoder) & N/A & 24.69 $\pm$ 0.00 & 4.53 $\pm$ 0.00 & 21.03 $\pm$ 0.00 & 54.53 $\pm$ 0.00 & 18.41 $\pm$ 0.00 & 1.54 $\pm$ 0.00 & 3.72 $\pm$ 0.00 \\
Base (WizardMath) & Base (WizardMath) & N/A & 31.01 $\pm$ 0.00 & 6.41 $\pm$ 0.00 & 3.12 $\pm$ 0.00 & 55.00 $\pm$ 0.00 & 14.02 $\pm$ 0.00 & 2.03 $\pm$ 0.00 & 2.77 $\pm$ 0.00 \\
\bottomrule
\end{tabularx}
\end{table}


% --- ベースモデル評価 シード別詳細 ---
\begin{table}[H]
\centering
\caption{ベースモデル評価 シード別詳細}
\label{tab:base_models_seeds}
\small
\begin{tabular}{\textwidth}{llllRRRRRRR}
\toprule
\textbf{Pattern} & \textbf{Method} & \textbf{Alpha} & \textbf{Seed} & \textbf{Safety Ave [Harmful Content] (ASR$\downarrow$ \%)} & \textbf{Math Ave (\%)} & \textbf{Code Ave (\%)} & \textbf{Medical Ave (\%)} & \textbf{General/Inst Ave (\%)} & \textbf{EvolCode (PPL$\downarrow$)} & \textbf{MedAlpaca (PPL$\downarrow$)} \\
\midrule
Base (MedAlpaca) & Base (MedAlpaca) & N/A & 42.0 & N/A & N/A & N/A & N/A & N/A & N/A & 7.26 \\
Base (MedAlpaca) & Base (MedAlpaca) & N/A & 43.0 & N/A & N/A & N/A & N/A & N/A & N/A & 7.18 \\
Base (MedAlpaca) & Base (MedAlpaca) & N/A & 44.0 & N/A & N/A & N/A & N/A & N/A & N/A & 7.00 \\
Base (MedAlpaca) & Base (MedAlpaca) & N/A & nan & 5.17 & 4.53 & 3.59 & 52.81 & 20.28 & 2.32 & 1.40 \\
Base (SafetyFT) & Base (SafetyFT) & N/A & 42.0 & 0.00 & 0.94 & 8.12 & 56.72 & 11.86 & 2.31 & N/A \\
Base (SafetyFT) & Base (SafetyFT) & N/A & 43.0 & 0.00 & 0.62 & 2.03 & 39.69 & 11.13 & 2.36 & N/A \\
Base (SafetyFT) & Base (SafetyFT) & N/A & 44.0 & 0.00 & 0.00 & 2.50 & 34.06 & 11.16 & 2.36 & N/A \\
Base (WizardCoder) & Base (WizardCoder) & N/A & nan & 24.69 & 4.53 & 21.03 & 54.53 & 18.41 & 1.54 & 3.72 \\
Base (WizardMath) & Base (WizardMath) & N/A & nan & 31.01 & 6.41 & 3.12 & 55.00 & 14.02 & 2.03 & 2.77 \\
\bottomrule
\end{tabularx}
\end{table}


% --------------------------------------------------------------------------
% 4. Alpha スイープ (Alpha Sweep)
% --------------------------------------------------------------------------

% --- Alpha スイープ 集計 - Method: matena_fisher, Pattern: safety+math+code+medical ---
\begin{table}[H]
\centering
\caption{Alpha スイープ 集計 - Method: matena_fisher, Pattern: safety+math+code+medical}
\label{tab:sweep_matenafisher_safety_math_code_medical}
\small
\begin{tabular}{\textwidth}{lllRRRRRRR}
\toprule
\textbf{Pattern} & \textbf{Method} & \textbf{Alpha} & \textbf{Safety Ave [Harmful Content] (ASR$\downarrow$ \%)} & \textbf{Math Ave (\%)} & \textbf{Code Ave (\%)} & \textbf{Medical Ave (\%)} & \textbf{General/Inst Ave (\%)} & \textbf{EvolCode (PPL$\downarrow$)} & \textbf{MedAlpaca (PPL$\downarrow$)} \\
\midrule
safety+math+code+medical & matena\_fisher & N/A & 0.00 $\pm$ 0.00 & 0.99 $\pm$ 0.24 & 0.00 $\pm$ 0.00 & 27.81 $\pm$ 3.29 & 13.12 $\pm$ 0.12 & 63.52 $\pm$ 2.39 & 128.84 $\pm$ 4.62 \\
\bottomrule
\end{tabularx}
\end{table}


% --- Alpha スイープ シード別詳細 - Method: matena_fisher, Pattern: safety+math+code+medical ---
\begin{table}[H]
\centering
\caption{Alpha スイープ シード別詳細 - Method: matena_fisher, Pattern: safety+math+code+medical}
\label{tab:sweep_seeds_matenafisher_safety_math_code_medical}
\small
\begin{tabular}{\textwidth}{llllRRRRRRR}
\toprule
\textbf{Pattern} & \textbf{Method} & \textbf{Alpha} & \textbf{Seed} & \textbf{Safety Ave [Harmful Content] (ASR$\downarrow$ \%)} & \textbf{Math Ave (\%)} & \textbf{Code Ave (\%)} & \textbf{Medical Ave (\%)} & \textbf{General/Inst Ave (\%)} & \textbf{EvolCode (PPL$\downarrow$)} & \textbf{MedAlpaca (PPL$\downarrow$)} \\
\midrule
safety+math+code+medical & matena\_fisher & N/A & 42.0 & 0.00 & 0.94 & 0.00 & 30.94 & 12.98 & 65.01 & 124.79 \\
safety+math+code+medical & matena\_fisher & N/A & 43.0 & 0.00 & 0.78 & 0.00 & 28.12 & 13.15 & 60.76 & 127.88 \\
safety+math+code+medical & matena\_fisher & N/A & 44.0 & 0.00 & 1.25 & 0.00 & 24.38 & 13.22 & 64.78 & 133.87 \\
\bottomrule
\end{tabularx}
\end{table}


% --- Alpha スイープ 集計 - Method: matena_fisher, Pattern: safety+math ---
\begin{table}[H]
\centering
\caption{Alpha スイープ 集計 - Method: matena_fisher, Pattern: safety+math}
\label{tab:sweep_matenafisher_safety_math}
\small
\begin{tabular}{\textwidth}{lllRRRRRRR}
\toprule
\textbf{Pattern} & \textbf{Method} & \textbf{Alpha} & \textbf{Safety Ave [Harmful Content] (ASR$\downarrow$ \%)} & \textbf{Math Ave (\%)} & \textbf{Code Ave (\%)} & \textbf{Medical Ave (\%)} & \textbf{General/Inst Ave (\%)} & \textbf{EvolCode (PPL$\downarrow$)} & \textbf{MedAlpaca (PPL$\downarrow$)} \\
\midrule
safety+math & matena\_fisher & N/A & 4.29 $\pm$ 0.43 & 6.04 $\pm$ 0.86 & 6.61 $\pm$ 0.39 & 85.73 $\pm$ 0.65 & 26.72 $\pm$ 11.70 & 1.82 $\pm$ 0.01 & 2.88 $\pm$ 0.02 \\
\bottomrule
\end{tabularx}
\end{table}


% --- Alpha スイープ シード別詳細 - Method: matena_fisher, Pattern: safety+math ---
\begin{table}[H]
\centering
\caption{Alpha スイープ シード別詳細 - Method: matena_fisher, Pattern: safety+math}
\label{tab:sweep_seeds_matenafisher_safety_math}
\small
\begin{tabular}{\textwidth}{llllRRRRRRR}
\toprule
\textbf{Pattern} & \textbf{Method} & \textbf{Alpha} & \textbf{Seed} & \textbf{Safety Ave [Harmful Content] (ASR$\downarrow$ \%)} & \textbf{Math Ave (\%)} & \textbf{Code Ave (\%)} & \textbf{Medical Ave (\%)} & \textbf{General/Inst Ave (\%)} & \textbf{EvolCode (PPL$\downarrow$)} & \textbf{MedAlpaca (PPL$\downarrow$)} \\
\midrule
safety+math & matena\_fisher & N/A & 42.0 & 4.71 & 6.09 & 6.25 & 85.00 & 39.12 & 1.81 & 2.90 \\
safety+math & matena\_fisher & N/A & 43.0 & 4.30 & 5.16 & 7.03 & 86.25 & 15.87 & 1.83 & 2.86 \\
safety+math & matena\_fisher & N/A & 44.0 & 3.85 & 6.88 & 6.56 & 85.94 & 25.17 & 1.84 & 2.88 \\
\bottomrule
\end{tabularx}
\end{table}


% --- Alpha スイープ 集計 - Method: matena_fisher, Pattern: safety+code ---
\begin{table}[H]
\centering
\caption{Alpha スイープ 集計 - Method: matena_fisher, Pattern: safety+code}
\label{tab:sweep_matenafisher_safety_code}
\small
\begin{tabular}{\textwidth}{lllRRRRRRR}
\toprule
\textbf{Pattern} & \textbf{Method} & \textbf{Alpha} & \textbf{Safety Ave [Harmful Content] (ASR$\downarrow$ \%)} & \textbf{Math Ave (\%)} & \textbf{Code Ave (\%)} & \textbf{Medical Ave (\%)} & \textbf{General/Inst Ave (\%)} & \textbf{EvolCode (PPL$\downarrow$)} & \textbf{MedAlpaca (PPL$\downarrow$)} \\
\midrule
safety+code & matena\_fisher & N/A & 3.37 $\pm$ 1.40 & 1.09 $\pm$ 0.00 & 0.00 $\pm$ 0.00 & 83.12 $\pm$ 0.31 & 12.77 $\pm$ 0.37 & 69.53 $\pm$ 21.34 & 84.66 $\pm$ 27.77 \\
\bottomrule
\end{tabularx}
\end{table}


% --- Alpha スイープ シード別詳細 - Method: matena_fisher, Pattern: safety+code ---
\begin{table}[H]
\centering
\caption{Alpha スイープ シード別詳細 - Method: matena_fisher, Pattern: safety+code}
\label{tab:sweep_seeds_matenafisher_safety_code}
\small
\begin{tabular}{\textwidth}{llllRRRRRRR}
\toprule
\textbf{Pattern} & \textbf{Method} & \textbf{Alpha} & \textbf{Seed} & \textbf{Safety Ave [Harmful Content] (ASR$\downarrow$ \%)} & \textbf{Math Ave (\%)} & \textbf{Code Ave (\%)} & \textbf{Medical Ave (\%)} & \textbf{General/Inst Ave (\%)} & \textbf{EvolCode (PPL$\downarrow$)} & \textbf{MedAlpaca (PPL$\downarrow$)} \\
\midrule
safety+code & matena\_fisher & N/A & 42.0 & 4.96 & 1.09 & 0.00 & 83.44 & 12.87 & 47.93 & 54.86 \\
safety+code & matena\_fisher & N/A & 43.0 & 2.85 & 1.09 & 0.00 & 82.81 & 13.07 & 70.07 & 89.29 \\
safety+code & matena\_fisher & N/A & 44.0 & 2.31 & 1.09 & 0.00 & 83.12 & 12.36 & 90.59 & 109.82 \\
\bottomrule
\end{tabularx}
\end{table}


% --- Alpha スイープ 集計 - Method: matena_fisher, Pattern: safety+medical ---
\begin{table}[H]
\centering
\caption{Alpha スイープ 集計 - Method: matena_fisher, Pattern: safety+medical}
\label{tab:sweep_matenafisher_safety_medical}
\small
\begin{tabular}{\textwidth}{lllRRRRRRR}
\toprule
\textbf{Pattern} & \textbf{Method} & \textbf{Alpha} & \textbf{Safety Ave [Harmful Content] (ASR$\downarrow$ \%)} & \textbf{Math Ave (\%)} & \textbf{Code Ave (\%)} & \textbf{Medical Ave (\%)} & \textbf{General/Inst Ave (\%)} & \textbf{EvolCode (PPL$\downarrow$)} & \textbf{MedAlpaca (PPL$\downarrow$)} \\
\midrule
safety+medical & matena\_fisher & N/A & 0.16 $\pm$ 0.00 & 0.00 $\pm$ 0.00 & 0.00 $\pm$ 0.00 & 0.00 $\pm$ 0.00 & 12.99 $\pm$ 0.18 & 576870.77 $\pm$ 100574.92 & 439254.45 $\pm$ 88250.04 \\
\bottomrule
\end{tabularx}
\end{table}


% --- Alpha スイープ シード別詳細 - Method: matena_fisher, Pattern: safety+medical ---
\begin{table}[H]
\centering
\caption{Alpha スイープ シード別詳細 - Method: matena_fisher, Pattern: safety+medical}
\label{tab:sweep_seeds_matenafisher_safety_medical}
\small
\begin{tabular}{\textwidth}{llllRRRRRRR}
\toprule
\textbf{Pattern} & \textbf{Method} & \textbf{Alpha} & \textbf{Seed} & \textbf{Safety Ave [Harmful Content] (ASR$\downarrow$ \%)} & \textbf{Math Ave (\%)} & \textbf{Code Ave (\%)} & \textbf{Medical Ave (\%)} & \textbf{General/Inst Ave (\%)} & \textbf{EvolCode (PPL$\downarrow$)} & \textbf{MedAlpaca (PPL$\downarrow$)} \\
\midrule
safety+medical & matena\_fisher & N/A & 42.0 & 0.16 & 0.00 & 0.00 & 0.00 & 13.09 & 691154.34 & 541088.28 \\
safety+medical & matena\_fisher & N/A & 43.0 & 0.16 & 0.00 & 0.00 & 0.00 & 12.78 & 501847.06 & 385101.33 \\
safety+medical & matena\_fisher & N/A & 44.0 & 0.16 & 0.00 & 0.00 & 0.00 & 13.09 & 537610.90 & 391573.73 \\
\bottomrule
\end{tabularx}
\end{table}


% --- Alpha スイープ 集計 - Method: diagonal_sst_main, Pattern: safety+math+code+medical ---
\begin{table}[H]
\centering
\caption{Alpha スイープ 集計 - Method: diagonal_sst_main, Pattern: safety+math+code+medical}
\label{tab:sweep_diagonalsstmain_safety_math_code_medical}
\small
\begin{tabular}{\textwidth}{lllRRRRRRR}
\toprule
\textbf{Pattern} & \textbf{Method} & \textbf{Alpha} & \textbf{Safety Ave [Harmful Content] (ASR$\downarrow$ \%)} & \textbf{Math Ave (\%)} & \textbf{Code Ave (\%)} & \textbf{Medical Ave (\%)} & \textbf{General/Inst Ave (\%)} & \textbf{EvolCode (PPL$\downarrow$)} & \textbf{MedAlpaca (PPL$\downarrow$)} \\
\midrule
safety+math+code+medical & diagonal\_sst\_main & 0.00 & 0.16 $\pm$ 0.00 & 0.00 $\pm$ 0.00 & 0.00 $\pm$ 0.00 & 12.19 $\pm$ 0.00 & 10.72 $\pm$ 0.06 & 48514.43 $\pm$ 0.00 & 60540.72 $\pm$ 0.00 \\
safety+math+code+medical & diagonal\_sst\_main & 0.20 & 0.05 $\pm$ 0.05 & 0.31 $\pm$ 0.54 & 0.00 $\pm$ 0.00 & 13.75 $\pm$ 0.00 & 12.21 $\pm$ 0.86 & 3890.44 $\pm$ 579.86 & 11358.40 $\pm$ 1880.12 \\
safety+math+code+medical & diagonal\_sst\_main & 0.40 & 0.00 $\pm$ 0.00 & 0.99 $\pm$ 0.36 & 0.00 $\pm$ 0.00 & 13.75 $\pm$ 0.00 & 11.72 $\pm$ 0.25 & 2876.79 $\pm$ 1018.70 & 2210.66 $\pm$ 378.54 \\
safety+math+code+medical & diagonal\_sst\_main & 0.60 & 0.00 $\pm$ 0.00 & 0.42 $\pm$ 0.18 & 0.00 $\pm$ 0.00 & 57.29 $\pm$ 17.67 & 10.16 $\pm$ 0.49 & 319.70 $\pm$ 69.81 & 139.54 $\pm$ 27.87 \\
safety+math+code+medical & diagonal\_sst\_main & 0.80 & 0.13 $\pm$ 0.05 & 0.36 $\pm$ 0.18 & 0.00 $\pm$ 0.00 & 82.81 $\pm$ 3.79 & 9.54 $\pm$ 0.92 & 119.22 $\pm$ 47.28 & 75.00 $\pm$ 24.40 \\
safety+math+code+medical & diagonal\_sst\_main & 1.00 & 0.16 $\pm$ 0.00 & 0.26 $\pm$ 0.24 & 0.00 $\pm$ 0.00 & 80.83 $\pm$ 8.57 & 11.10 $\pm$ 1.30 & 78.49 $\pm$ 39.87 & 78.49 $\pm$ 32.65 \\
\bottomrule
\end{tabularx}
\end{table}


% --- Alpha スイープ シード別詳細 - Method: diagonal_sst_main, Pattern: safety+math+code+medical ---
\begin{table}[H]
\centering
\caption{Alpha スイープ シード別詳細 - Method: diagonal_sst_main, Pattern: safety+math+code+medical}
\label{tab:sweep_seeds_diagonalsstmain_safety_math_code_medical}
\small
\begin{tabular}{\textwidth}{llllRRRRRRR}
\toprule
\textbf{Pattern} & \textbf{Method} & \textbf{Alpha} & \textbf{Seed} & \textbf{Safety Ave [Harmful Content] (ASR$\downarrow$ \%)} & \textbf{Math Ave (\%)} & \textbf{Code Ave (\%)} & \textbf{Medical Ave (\%)} & \textbf{General/Inst Ave (\%)} & \textbf{EvolCode (PPL$\downarrow$)} & \textbf{MedAlpaca (PPL$\downarrow$)} \\
\midrule
safety+math+code+medical & diagonal\_sst\_main & 0.00 & 42.0 & 0.16 & 0.00 & 0.00 & 12.19 & 10.79 & 48514.43 & 60540.72 \\
safety+math+code+medical & diagonal\_sst\_main & 0.00 & 43.0 & 0.16 & 0.00 & 0.00 & 12.19 & 10.68 & 48514.43 & 60540.72 \\
safety+math+code+medical & diagonal\_sst\_main & 0.00 & 44.0 & 0.16 & 0.00 & 0.00 & 12.19 & 10.68 & 48514.43 & 60540.72 \\
safety+math+code+medical & diagonal\_sst\_main & 0.20 & 42.0 & 0.08 & 0.00 & 0.00 & 13.75 & 11.83 & 3837.57 & 11701.23 \\
safety+math+code+medical & diagonal\_sst\_main & 0.20 & 43.0 & 0.00 & 0.00 & 0.00 & 13.75 & 13.18 & 3338.82 & 9330.45 \\
safety+math+code+medical & diagonal\_sst\_main & 0.20 & 44.0 & 0.08 & 0.94 & 0.00 & 13.75 & 11.60 & 4494.92 & 13043.51 \\
safety+math+code+medical & diagonal\_sst\_main & 0.40 & 42.0 & 0.00 & 1.41 & 0.00 & 13.75 & 11.43 & 1839.14 & 2574.76 \\
safety+math+code+medical & diagonal\_sst\_main & 0.40 & 43.0 & 0.00 & 0.78 & 0.00 & 13.75 & 11.82 & 3875.43 & 2238.06 \\
safety+math+code+medical & diagonal\_sst\_main & 0.40 & 44.0 & 0.00 & 0.78 & 0.00 & 13.75 & 11.90 & 2915.80 & 1819.17 \\
safety+math+code+medical & diagonal\_sst\_main & 0.60 & 42.0 & 0.00 & 0.62 & 0.00 & 40.00 & 10.41 & 392.09 & 168.34 \\
safety+math+code+medical & diagonal\_sst\_main & 0.60 & 43.0 & 0.00 & 0.31 & 0.00 & 75.31 & 10.48 & 252.79 & 112.69 \\
safety+math+code+medical & diagonal\_sst\_main & 0.60 & 44.0 & 0.00 & 0.31 & 0.00 & 56.56 & 9.59 & 314.23 & 137.59 \\
safety+math+code+medical & diagonal\_sst\_main & 0.80 & 42.0 & 0.08 & 0.47 & 0.00 & 85.00 & 10.55 & 167.99 & 100.74 \\
safety+math+code+medical & diagonal\_sst\_main & 0.80 & 43.0 & 0.16 & 0.47 & 0.00 & 85.00 & 9.31 & 73.59 & 52.21 \\
safety+math+code+medical & diagonal\_sst\_main & 0.80 & 44.0 & 0.16 & 0.16 & 0.00 & 78.44 & 8.76 & 116.08 & 72.05 \\
safety+math+code+medical & diagonal\_sst\_main & 1.00 & 42.0 & 0.16 & 0.47 & 0.00 & 85.94 & 12.28 & 119.75 & 108.05 \\
safety+math+code+medical & diagonal\_sst\_main & 1.00 & 43.0 & 0.16 & 0.31 & 0.00 & 85.62 & 11.31 & 40.18 & 43.44 \\
safety+math+code+medical & diagonal\_sst\_main & 1.00 & 44.0 & 0.16 & 0.00 & 0.00 & 70.94 & 9.71 & 75.56 & 83.97 \\
\bottomrule
\end{tabularx}
\end{table}


% --- Alpha スイープ 集計 - Method: diagonal_sst_main, Pattern: safety+math ---
\begin{table}[H]
\centering
\caption{Alpha スイープ 集計 - Method: diagonal_sst_main, Pattern: safety+math}
\label{tab:sweep_diagonalsstmain_safety_math}
\small
\begin{tabular}{\textwidth}{lllRRRRRRR}
\toprule
\textbf{Pattern} & \textbf{Method} & \textbf{Alpha} & \textbf{Safety Ave [Harmful Content] (ASR$\downarrow$ \%)} & \textbf{Math Ave (\%)} & \textbf{Code Ave (\%)} & \textbf{Medical Ave (\%)} & \textbf{General/Inst Ave (\%)} & \textbf{EvolCode (PPL$\downarrow$)} & \textbf{MedAlpaca (PPL$\downarrow$)} \\
\midrule
safety+math & diagonal\_sst\_main & 0.00 & 34.32 $\pm$ 0.20 & 6.04 $\pm$ 0.18 & 2.19 $\pm$ 0.00 & 63.44 $\pm$ 28.69 & 24.87 $\pm$ 9.33 & 1.99 $\pm$ 0.00 & 2.71 $\pm$ 0.00 \\
safety+math & diagonal\_sst\_main & 0.20 & 31.84 $\pm$ 1.33 & 6.46 $\pm$ 0.63 & 4.53 $\pm$ 0.83 & 62.71 $\pm$ 28.61 & 21.02 $\pm$ 10.84 & 1.91 $\pm$ 0.00 & 2.66 $\pm$ 0.01 \\
safety+math & diagonal\_sst\_main & 0.40 & 32.86 $\pm$ 2.20 & 6.20 $\pm$ 0.39 & 4.84 $\pm$ 0.68 & 71.61 $\pm$ 15.47 & 27.05 $\pm$ 11.87 & 1.86 $\pm$ 0.01 & 2.63 $\pm$ 0.02 \\
safety+math & diagonal\_sst\_main & 0.60 & 17.09 $\pm$ 5.01 & 6.82 $\pm$ 0.86 & 6.72 $\pm$ 0.95 & 72.55 $\pm$ 15.91 & 22.48 $\pm$ 8.92 & 1.84 $\pm$ 0.01 & 2.75 $\pm$ 0.05 \\
safety+math & diagonal\_sst\_main & 0.80 & 2.76 $\pm$ 1.73 & 8.44 $\pm$ 1.02 & 5.89 $\pm$ 2.13 & 73.02 $\pm$ 17.23 & 17.72 $\pm$ 6.02 & 1.88 $\pm$ 0.04 & 3.15 $\pm$ 0.15 \\
safety+math & diagonal\_sst\_main & 1.00 & 0.00 $\pm$ 0.00 & 9.84 $\pm$ 0.47 & 7.76 $\pm$ 3.63 & 65.42 $\pm$ 31.57 & 20.65 $\pm$ 6.27 & 1.98 $\pm$ 0.07 & 4.36 $\pm$ 0.49 \\
\bottomrule
\end{tabularx}
\end{table}


% --- Alpha スイープ シード別詳細 - Method: diagonal_sst_main, Pattern: safety+math ---
\begin{table}[H]
\centering
\caption{Alpha スイープ シード別詳細 - Method: diagonal_sst_main, Pattern: safety+math}
\label{tab:sweep_seeds_diagonalsstmain_safety_math}
\small
\begin{tabular}{\textwidth}{llllRRRRRRR}
\toprule
\textbf{Pattern} & \textbf{Method} & \textbf{Alpha} & \textbf{Seed} & \textbf{Safety Ave [Harmful Content] (ASR$\downarrow$ \%)} & \textbf{Math Ave (\%)} & \textbf{Code Ave (\%)} & \textbf{Medical Ave (\%)} & \textbf{General/Inst Ave (\%)} & \textbf{EvolCode (PPL$\downarrow$)} & \textbf{MedAlpaca (PPL$\downarrow$)} \\
\midrule
safety+math & diagonal\_sst\_main & 0.00 & 42.0 & 34.43 & 5.94 & 2.19 & 30.31 & 14.70 & 1.99 & 2.71 \\
safety+math & diagonal\_sst\_main & 0.00 & 43.0 & 34.44 & 5.94 & 2.19 & 80.00 & 33.04 & 1.99 & 2.71 \\
safety+math & diagonal\_sst\_main & 0.00 & 44.0 & 34.10 & 6.25 & 2.19 & 80.00 & 26.88 & 1.99 & 2.71 \\
safety+math & diagonal\_sst\_main & 0.20 & 42.0 & 31.57 & 6.09 & 5.47 & 29.69 & 14.42 & 1.91 & 2.65 \\
safety+math & diagonal\_sst\_main & 0.20 & 43.0 & 33.28 & 6.09 & 4.22 & 78.44 & 15.12 & 1.91 & 2.66 \\
safety+math & diagonal\_sst\_main & 0.20 & 44.0 & 30.66 & 7.19 & 3.91 & 80.00 & 33.53 & 1.92 & 2.65 \\
safety+math & diagonal\_sst\_main & 0.40 & 42.0 & 34.17 & 6.25 & 4.06 & 53.91 & 13.80 & 1.85 & 2.60 \\
safety+math & diagonal\_sst\_main & 0.40 & 43.0 & 30.32 & 5.78 & 5.16 & 78.44 & 30.63 & 1.86 & 2.65 \\
safety+math & diagonal\_sst\_main & 0.40 & 44.0 & 34.11 & 6.56 & 5.31 & 82.50 & 36.72 & 1.86 & 2.64 \\
safety+math & diagonal\_sst\_main & 0.60 & 42.0 & 22.60 & 6.25 & 7.34 & 54.22 & 13.40 & 1.83 & 2.69 \\
safety+math & diagonal\_sst\_main & 0.60 & 43.0 & 15.85 & 6.41 & 7.19 & 80.62 & 31.23 & 1.84 & 2.79 \\
safety+math & diagonal\_sst\_main & 0.60 & 44.0 & 12.81 & 7.81 & 5.62 & 82.81 & 22.82 & 1.85 & 2.78 \\
safety+math & diagonal\_sst\_main & 0.80 & 42.0 & 3.99 & 7.50 & 6.88 & 53.12 & 13.76 & 1.83 & 2.97 \\
safety+math & diagonal\_sst\_main & 0.80 & 43.0 & 0.78 & 8.28 & 7.34 & 82.81 & 24.65 & 1.90 & 3.23 \\
safety+math & diagonal\_sst\_main & 0.80 & 44.0 & 3.52 & 9.53 & 3.44 & 83.12 & 14.76 & 1.90 & 3.24 \\
safety+math & diagonal\_sst\_main & 1.00 & 42.0 & 0.00 & 9.38 & 11.88 & 29.06 & 13.45 & 1.90 & 3.82 \\
safety+math & diagonal\_sst\_main & 1.00 & 43.0 & 0.00 & 10.31 & 6.41 & 81.25 & 24.94 & 2.03 & 4.79 \\
safety+math & diagonal\_sst\_main & 1.00 & 44.0 & 0.00 & 9.84 & 5.00 & 85.94 & 23.56 & 2.00 & 4.46 \\
\bottomrule
\end{tabularx}
\end{table}


% --- Alpha スイープ 集計 - Method: diagonal_sst_main, Pattern: safety+code ---
\begin{table}[H]
\centering
\caption{Alpha スイープ 集計 - Method: diagonal_sst_main, Pattern: safety+code}
\label{tab:sweep_diagonalsstmain_safety_code}
\small
\begin{tabular}{\textwidth}{lllRRRRRRR}
\toprule
\textbf{Pattern} & \textbf{Method} & \textbf{Alpha} & \textbf{Safety Ave [Harmful Content] (ASR$\downarrow$ \%)} & \textbf{Math Ave (\%)} & \textbf{Code Ave (\%)} & \textbf{Medical Ave (\%)} & \textbf{General/Inst Ave (\%)} & \textbf{EvolCode (PPL$\downarrow$)} & \textbf{MedAlpaca (PPL$\downarrow$)} \\
\midrule
safety+code & diagonal\_sst\_main & 0.00 & 0.08 $\pm$ 0.00 & 0.47 $\pm$ 0.00 & 0.00 $\pm$ 0.00 & 85.62 $\pm$ 0.00 & 13.66 $\pm$ 0.00 & 485.34 $\pm$ 0.00 & 92.12 $\pm$ 0.00 \\
safety+code & diagonal\_sst\_main & 0.20 & 0.08 $\pm$ 0.08 & 0.36 $\pm$ 0.18 & 0.00 $\pm$ 0.00 & 86.35 $\pm$ 0.18 & 13.66 $\pm$ 0.15 & 304.22 $\pm$ 54.14 & 58.85 $\pm$ 7.32 \\
safety+code & diagonal\_sst\_main & 0.40 & 0.32 $\pm$ 0.14 & 0.83 $\pm$ 0.63 & 0.00 $\pm$ 0.00 & 86.46 $\pm$ 0.18 & 13.44 $\pm$ 0.46 & 296.86 $\pm$ 96.53 & 42.89 $\pm$ 11.58 \\
safety+code & diagonal\_sst\_main & 0.60 & 0.97 $\pm$ 0.42 & 0.89 $\pm$ 0.24 & 0.00 $\pm$ 0.00 & 86.25 $\pm$ 0.31 & 12.26 $\pm$ 0.55 & 191.50 $\pm$ 74.10 & 31.40 $\pm$ 8.25 \\
safety+code & diagonal\_sst\_main & 0.80 & 0.82 $\pm$ 0.33 & 1.09 $\pm$ 0.41 & 0.00 $\pm$ 0.00 & 80.73 $\pm$ 3.96 & 11.67 $\pm$ 0.24 & 48.92 $\pm$ 16.08 & 26.83 $\pm$ 3.50 \\
safety+code & diagonal\_sst\_main & 1.00 & 0.00 $\pm$ 0.00 & 0.42 $\pm$ 0.24 & 0.00 $\pm$ 0.00 & 15.21 $\pm$ 0.95 & 10.63 $\pm$ 0.27 & 187.93 $\pm$ 30.20 & 249.26 $\pm$ 144.36 \\
\bottomrule
\end{tabularx}
\end{table}


% --- Alpha スイープ シード別詳細 - Method: diagonal_sst_main, Pattern: safety+code ---
\begin{table}[H]
\centering
\caption{Alpha スイープ シード別詳細 - Method: diagonal_sst_main, Pattern: safety+code}
\label{tab:sweep_seeds_diagonalsstmain_safety_code}
\small
\begin{tabular}{\textwidth}{llllRRRRRRR}
\toprule
\textbf{Pattern} & \textbf{Method} & \textbf{Alpha} & \textbf{Seed} & \textbf{Safety Ave [Harmful Content] (ASR$\downarrow$ \%)} & \textbf{Math Ave (\%)} & \textbf{Code Ave (\%)} & \textbf{Medical Ave (\%)} & \textbf{General/Inst Ave (\%)} & \textbf{EvolCode (PPL$\downarrow$)} & \textbf{MedAlpaca (PPL$\downarrow$)} \\
\midrule
safety+code & diagonal\_sst\_main & 0.00 & 42.0 & 0.08 & 0.47 & 0.00 & 85.62 & 13.66 & 485.34 & 92.12 \\
safety+code & diagonal\_sst\_main & 0.00 & 43.0 & 0.08 & 0.47 & 0.00 & 85.62 & 13.66 & 485.34 & 92.12 \\
safety+code & diagonal\_sst\_main & 0.00 & 44.0 & 0.08 & 0.47 & 0.00 & 85.62 & 13.66 & 485.34 & 92.12 \\
safety+code & diagonal\_sst\_main & 0.20 & 42.0 & 0.08 & 0.47 & 0.00 & 86.56 & 13.49 & 259.22 & 51.06 \\
safety+code & diagonal\_sst\_main & 0.20 & 43.0 & 0.00 & 0.47 & 0.00 & 86.25 & 13.79 & 364.30 & 65.57 \\
safety+code & diagonal\_sst\_main & 0.20 & 44.0 & 0.16 & 0.16 & 0.00 & 86.25 & 13.72 & 289.15 & 59.92 \\
safety+code & diagonal\_sst\_main & 0.40 & 42.0 & 0.24 & 0.47 & 0.00 & 86.25 & 12.92 & 223.61 & 29.53 \\
safety+code & diagonal\_sst\_main & 0.40 & 43.0 & 0.24 & 1.56 & 0.00 & 86.56 & 13.65 & 406.24 & 49.45 \\
safety+code & diagonal\_sst\_main & 0.40 & 44.0 & 0.48 & 0.47 & 0.00 & 86.56 & 13.75 & 260.72 & 49.71 \\
safety+code & diagonal\_sst\_main & 0.60 & 42.0 & 1.21 & 1.09 & 0.00 & 86.25 & 12.88 & 135.96 & 21.94 \\
safety+code & diagonal\_sst\_main & 0.60 & 43.0 & 1.20 & 0.94 & 0.00 & 85.94 & 11.85 & 275.64 & 35.11 \\
safety+code & diagonal\_sst\_main & 0.60 & 44.0 & 0.49 & 0.62 & 0.00 & 86.56 & 12.05 & 162.89 & 37.14 \\
safety+code & diagonal\_sst\_main & 0.80 & 42.0 & 0.55 & 0.62 & 0.00 & 77.19 & 11.40 & 39.85 & 24.20 \\
safety+code & diagonal\_sst\_main & 0.80 & 43.0 & 0.71 & 1.41 & 0.00 & 80.00 & 11.86 & 67.49 & 25.49 \\
safety+code & diagonal\_sst\_main & 0.80 & 44.0 & 1.18 & 1.25 & 0.00 & 85.00 & 11.77 & 39.42 & 30.80 \\
safety+code & diagonal\_sst\_main & 1.00 & 42.0 & 0.00 & 0.47 & 0.00 & 16.25 & 10.94 & 221.80 & 409.75 \\
safety+code & diagonal\_sst\_main & 1.00 & 43.0 & 0.00 & 0.62 & 0.00 & 14.37 & 10.47 & 178.18 & 130.00 \\
safety+code & diagonal\_sst\_main & 1.00 & 44.0 & 0.00 & 0.16 & 0.00 & 15.00 & 10.49 & 163.80 & 208.02 \\
\bottomrule
\end{tabularx}
\end{table}


% --- Alpha スイープ 集計 - Method: diagonal_sst_main, Pattern: safety+medical ---
\begin{table}[H]
\centering
\caption{Alpha スイープ 集計 - Method: diagonal_sst_main, Pattern: safety+medical}
\label{tab:sweep_diagonalsstmain_safety_medical}
\small
\begin{tabular}{\textwidth}{lllRRRRRRR}
\toprule
\textbf{Pattern} & \textbf{Method} & \textbf{Alpha} & \textbf{Safety Ave [Harmful Content] (ASR$\downarrow$ \%)} & \textbf{Math Ave (\%)} & \textbf{Code Ave (\%)} & \textbf{Medical Ave (\%)} & \textbf{General/Inst Ave (\%)} & \textbf{EvolCode (PPL$\downarrow$)} & \textbf{MedAlpaca (PPL$\downarrow$)} \\
\midrule
safety+medical & diagonal\_sst\_main & 0.00 & 0.00 $\pm$ 0.00 & 0.00 $\pm$ 0.00 & 0.00 $\pm$ 0.00 & 33.12 $\pm$ 0.00 & 12.06 $\pm$ 0.10 & 180460.88 $\pm$ 0.00 & 169745.34 $\pm$ 0.00 \\
safety+medical & diagonal\_sst\_main & 0.20 & 0.13 $\pm$ 0.05 & 0.00 $\pm$ 0.00 & 0.00 $\pm$ 0.00 & 13.75 $\pm$ 0.31 & 12.12 $\pm$ 0.08 & 94281.86 $\pm$ 26088.99 & 81185.07 $\pm$ 29847.41 \\
safety+medical & diagonal\_sst\_main & 0.40 & 0.19 $\pm$ 0.20 & 0.00 $\pm$ 0.00 & 0.00 $\pm$ 0.00 & 85.73 $\pm$ 0.65 & 12.13 $\pm$ 0.32 & 131598.61 $\pm$ 37516.58 & 83673.35 $\pm$ 23943.66 \\
safety+medical & diagonal\_sst\_main & 0.60 & 0.05 $\pm$ 0.05 & 0.00 $\pm$ 0.00 & 0.00 $\pm$ 0.00 & 9.69 $\pm$ 8.46 & 12.06 $\pm$ 0.13 & 75380.55 $\pm$ 8006.67 & 78424.43 $\pm$ 9712.01 \\
safety+medical & diagonal\_sst\_main & 0.80 & 0.05 $\pm$ 0.05 & 0.00 $\pm$ 0.00 & 0.00 $\pm$ 0.00 & 4.06 $\pm$ 7.04 & 11.83 $\pm$ 0.45 & 135732.82 $\pm$ 163601.13 & 175313.17 $\pm$ 210610.19 \\
safety+medical & diagonal\_sst\_main & 1.00 & 0.05 $\pm$ 0.05 & 0.00 $\pm$ 0.00 & 0.00 $\pm$ 0.00 & 26.25 $\pm$ 22.89 & 11.87 $\pm$ 0.98 & 138717.43 $\pm$ 145875.51 & 95092.19 $\pm$ 65045.99 \\
\bottomrule
\end{tabularx}
\end{table}


% --- Alpha スイープ シード別詳細 - Method: diagonal_sst_main, Pattern: safety+medical ---
\begin{table}[H]
\centering
\caption{Alpha スイープ シード別詳細 - Method: diagonal_sst_main, Pattern: safety+medical}
\label{tab:sweep_seeds_diagonalsstmain_safety_medical}
\small
\begin{tabular}{\textwidth}{llllRRRRRRR}
\toprule
\textbf{Pattern} & \textbf{Method} & \textbf{Alpha} & \textbf{Seed} & \textbf{Safety Ave [Harmful Content] (ASR$\downarrow$ \%)} & \textbf{Math Ave (\%)} & \textbf{Code Ave (\%)} & \textbf{Medical Ave (\%)} & \textbf{General/Inst Ave (\%)} & \textbf{EvolCode (PPL$\downarrow$)} & \textbf{MedAlpaca (PPL$\downarrow$)} \\
\midrule
safety+medical & diagonal\_sst\_main & 0.00 & 42.0 & 0.00 & 0.00 & 0.00 & 33.12 & 11.96 & 180460.88 & 169745.34 \\
safety+medical & diagonal\_sst\_main & 0.00 & 43.0 & 0.00 & 0.00 & 0.00 & 33.12 & 12.06 & 180460.88 & 169745.34 \\
safety+medical & diagonal\_sst\_main & 0.00 & 44.0 & 0.00 & 0.00 & 0.00 & 33.12 & 12.17 & 180460.88 & 169745.34 \\
safety+medical & diagonal\_sst\_main & 0.20 & 42.0 & 0.16 & 0.00 & 0.00 & 13.75 & 12.21 & 122142.18 & 114130.75 \\
safety+medical & diagonal\_sst\_main & 0.20 & 43.0 & 0.08 & 0.00 & 0.00 & 13.44 & 12.05 & 90275.76 & 73476.06 \\
safety+medical & diagonal\_sst\_main & 0.20 & 44.0 & 0.16 & 0.00 & 0.00 & 14.06 & 12.10 & 70427.63 & 55948.41 \\
safety+medical & diagonal\_sst\_main & 0.40 & 42.0 & 0.40 & 0.00 & 0.00 & 85.00 & 11.87 & 167508.43 & 99865.73 \\
safety+medical & diagonal\_sst\_main & 0.40 & 43.0 & 0.16 & 0.00 & 0.00 & 85.94 & 12.49 & 92659.00 & 56169.58 \\
safety+medical & diagonal\_sst\_main & 0.40 & 44.0 & 0.00 & 0.00 & 0.00 & 86.25 & 12.03 & 134628.39 & 94984.74 \\
safety+medical & diagonal\_sst\_main & 0.60 & 42.0 & 0.08 & 0.00 & 0.00 & 13.44 & 11.91 & 72783.78 & 67362.99 \\
safety+medical & diagonal\_sst\_main & 0.60 & 43.0 & 0.00 & 0.00 & 0.00 & 0.00 & 12.16 & 84363.30 & 85554.10 \\
safety+medical & diagonal\_sst\_main & 0.60 & 44.0 & 0.08 & 0.00 & 0.00 & 15.62 & 12.11 & 68994.57 & 82356.19 \\
safety+medical & diagonal\_sst\_main & 0.80 & 42.0 & 0.00 & 0.00 & 0.00 & 12.19 & 12.33 & 8763.53 & 20450.35 \\
safety+medical & diagonal\_sst\_main & 0.80 & 43.0 & 0.08 & 0.00 & 0.00 & 0.00 & 11.47 & 320355.53 & 415132.33 \\
safety+medical & diagonal\_sst\_main & 0.80 & 44.0 & 0.08 & 0.00 & 0.00 & 0.00 & 11.68 & 78079.41 & 90356.82 \\
safety+medical & diagonal\_sst\_main & 1.00 & 42.0 & 0.00 & 0.00 & 0.00 & 7.81 & 11.46 & 11982.52 & 20544.71 \\
safety+medical & diagonal\_sst\_main & 1.00 & 43.0 & 0.08 & 0.00 & 0.00 & 51.88 & 11.16 & 298175.26 & 140302.29 \\
safety+medical & diagonal\_sst\_main & 1.00 & 44.0 & 0.08 & 0.00 & 0.00 & 19.06 & 13.00 & 105994.50 & 124429.57 \\
\bottomrule
\end{tabularx}
\end{table}


% --- Alpha スイープ 集計 - Method: dare, Pattern: safety+math+code+medical ---
\begin{table}[H]
\centering
\caption{Alpha スイープ 集計 - Method: dare, Pattern: safety+math+code+medical}
\label{tab:sweep_dare_safety_math_code_medical}
\small
\begin{tabular}{\textwidth}{lllRRRRRRR}
\toprule
\textbf{Pattern} & \textbf{Method} & \textbf{Alpha} & \textbf{Safety Ave [Harmful Content] (ASR$\downarrow$ \%)} & \textbf{Math Ave (\%)} & \textbf{Code Ave (\%)} & \textbf{Medical Ave (\%)} & \textbf{General/Inst Ave (\%)} & \textbf{EvolCode (PPL$\downarrow$)} & \textbf{MedAlpaca (PPL$\downarrow$)} \\
\midrule
safety+math+code+medical & dare & 0.00 & 0.11 $\pm$ 0.05 & 0.00 $\pm$ 0.00 & 0.00 $\pm$ 0.00 & 26.15 $\pm$ 45.29 & 11.63 $\pm$ 0.34 & 6818688.18 $\pm$ 7065619.14 & 7378650.47 $\pm$ 4295362.07 \\
safety+math+code+medical & dare & 0.20 & 0.13 $\pm$ 0.05 & 0.00 $\pm$ 0.00 & 0.00 $\pm$ 0.00 & 11.98 $\pm$ 7.81 & 12.06 $\pm$ 0.50 & 7173662.71 $\pm$ 2942125.13 & 6687130.41 $\pm$ 3055033.24 \\
safety+math+code+medical & dare & 0.40 & 0.03 $\pm$ 0.05 & 0.00 $\pm$ 0.00 & 0.00 $\pm$ 0.00 & 13.75 $\pm$ 0.00 & 12.07 $\pm$ 0.90 & 6979776.99 $\pm$ 6033524.78 & 2986909.55 $\pm$ 2662560.22 \\
safety+math+code+medical & dare & 0.60 & 0.08 $\pm$ 0.08 & 0.00 $\pm$ 0.00 & 0.00 $\pm$ 0.00 & 38.02 $\pm$ 41.77 & 11.53 $\pm$ 0.23 & 5368837.19 $\pm$ 6493560.14 & 6550135.45 $\pm$ 3035748.48 \\
safety+math+code+medical & dare & 0.80 & 0.05 $\pm$ 0.05 & 0.00 $\pm$ 0.00 & 0.00 $\pm$ 0.00 & 29.58 $\pm$ 48.81 & 12.30 $\pm$ 0.47 & 953765.09 $\pm$ 912212.52 & 828217.06 $\pm$ 375799.63 \\
safety+math+code+medical & dare & 1.00 & 0.00 $\pm$ 0.00 & 0.68 $\pm$ 0.70 & 5.36 $\pm$ 2.13 & 55.83 $\pm$ 25.27 & 17.94 $\pm$ 12.11 & 2.35 $\pm$ 0.03 & 7.15 $\pm$ 0.10 \\
\bottomrule
\end{tabularx}
\end{table}


% --- Alpha スイープ シード別詳細 - Method: dare, Pattern: safety+math+code+medical ---
\begin{table}[H]
\centering
\caption{Alpha スイープ シード別詳細 - Method: dare, Pattern: safety+math+code+medical}
\label{tab:sweep_seeds_dare_safety_math_code_medical}
\small
\begin{tabular}{\textwidth}{llllRRRRRRR}
\toprule
\textbf{Pattern} & \textbf{Method} & \textbf{Alpha} & \textbf{Seed} & \textbf{Safety Ave [Harmful Content] (ASR$\downarrow$ \%)} & \textbf{Math Ave (\%)} & \textbf{Code Ave (\%)} & \textbf{Medical Ave (\%)} & \textbf{General/Inst Ave (\%)} & \textbf{EvolCode (PPL$\downarrow$)} & \textbf{MedAlpaca (PPL$\downarrow$)} \\
\midrule
safety+math+code+medical & dare & 0.00 & 42.0 & 0.16 & 0.00 & 0.00 & 0.00 & 12.02 & 14730567.06 & 12336031.02 \\
safety+math+code+medical & dare & 0.00 & 43.0 & 0.08 & 0.00 & 0.00 & 78.44 & 11.39 & 1138044.95 & 4764244.25 \\
safety+math+code+medical & dare & 0.00 & 44.0 & 0.08 & 0.00 & 0.00 & 0.00 & 11.47 & 4587452.52 & 5035676.13 \\
safety+math+code+medical & dare & 0.20 & 42.0 & 0.16 & 0.00 & 0.00 & 13.75 & 12.58 & 4799372.95 & 3193735.46 \\
safety+math+code+medical & dare & 0.20 & 43.0 & 0.08 & 0.00 & 0.00 & 3.44 & 12.02 & 6256487.91 & 8009125.74 \\
safety+math+code+medical & dare & 0.20 & 44.0 & 0.16 & 0.00 & 0.00 & 18.75 & 11.58 & 10465127.28 & 8858530.02 \\
safety+math+code+medical & dare & 0.40 & 42.0 & 0.00 & 0.00 & 0.00 & 13.75 & 11.19 & 12369089.71 & 5828590.40 \\
safety+math+code+medical & dare & 0.40 & 43.0 & 0.00 & 0.00 & 0.00 & 13.75 & 12.99 & 8108718.76 & 2582368.48 \\
safety+math+code+medical & dare & 0.40 & 44.0 & 0.08 & 0.00 & 0.00 & 13.75 & 12.03 & 461522.51 & 549769.77 \\
safety+math+code+medical & dare & 0.60 & 42.0 & 0.16 & 0.00 & 0.00 & 86.25 & 11.42 & 12857479.52 & 10037855.54 \\
safety+math+code+medical & dare & 0.60 & 43.0 & 0.08 & 0.00 & 0.00 & 14.06 & 11.38 & 1950859.88 & 4501931.51 \\
safety+math+code+medical & dare & 0.60 & 44.0 & 0.00 & 0.00 & 0.00 & 13.75 & 11.79 & 1298172.17 & 5110619.29 \\
safety+math+code+medical & dare & 0.80 & 42.0 & 0.00 & 0.00 & 0.00 & 2.19 & 12.80 & 2005704.45 & 1261537.06 \\
safety+math+code+medical & dare & 0.80 & 43.0 & 0.08 & 0.00 & 0.00 & 85.94 & 12.21 & 380898.02 & 591539.57 \\
safety+math+code+medical & dare & 0.80 & 44.0 & 0.08 & 0.00 & 0.00 & 0.62 & 11.89 & 474692.80 & 631574.56 \\
safety+math+code+medical & dare & 1.00 & 42.0 & 0.00 & 1.41 & 7.81 & 84.06 & 31.92 & 2.31 & 7.26 \\
safety+math+code+medical & dare & 1.00 & 43.0 & 0.00 & 0.62 & 3.91 & 48.12 & 10.60 & 2.36 & 7.08 \\
safety+math+code+medical & dare & 1.00 & 44.0 & 0.00 & 0.00 & 4.38 & 35.31 & 11.30 & 2.37 & 7.10 \\
\bottomrule
\end{tabularx}
\end{table}


% --- Alpha スイープ 集計 - Method: dare, Pattern: safety+math ---
\begin{table}[H]
\centering
\caption{Alpha スイープ 集計 - Method: dare, Pattern: safety+math}
\label{tab:sweep_dare_safety_math}
\small
\begin{tabular}{\textwidth}{lllRRRRRRR}
\toprule
\textbf{Pattern} & \textbf{Method} & \textbf{Alpha} & \textbf{Safety Ave [Harmful Content] (ASR$\downarrow$ \%)} & \textbf{Math Ave (\%)} & \textbf{Code Ave (\%)} & \textbf{Medical Ave (\%)} & \textbf{General/Inst Ave (\%)} & \textbf{EvolCode (PPL$\downarrow$)} & \textbf{MedAlpaca (PPL$\downarrow$)} \\
\midrule
safety+math & dare & 0.00 & 33.14 $\pm$ 0.78 & 5.47 $\pm$ 0.47 & 3.33 $\pm$ 0.59 & 79.79 $\pm$ 0.79 & 30.41 $\pm$ 4.79 & 2.02 $\pm$ 0.00 & 2.77 $\pm$ 0.00 \\
safety+math & dare & 0.20 & 10.27 $\pm$ 1.33 & 4.79 $\pm$ 1.13 & 4.64 $\pm$ 1.26 & 84.27 $\pm$ 0.79 & 28.15 $\pm$ 5.63 & 2.36 $\pm$ 0.04 & 3.61 $\pm$ 0.05 \\
safety+math & dare & 0.40 & 9.09 $\pm$ 1.71 & 5.36 $\pm$ 0.65 & 4.95 $\pm$ 0.65 & 85.00 $\pm$ 0.54 & 23.52 $\pm$ 2.16 & 2.38 $\pm$ 0.05 & 3.85 $\pm$ 0.08 \\
safety+math & dare & 0.60 & 6.14 $\pm$ 5.14 & 5.42 $\pm$ 0.79 & 5.00 $\pm$ 1.02 & 84.27 $\pm$ 2.35 & 22.85 $\pm$ 8.57 & 2.43 $\pm$ 0.07 & 4.44 $\pm$ 0.28 \\
safety+math & dare & 0.80 & 3.59 $\pm$ 4.78 & 6.35 $\pm$ 0.70 & 5.68 $\pm$ 1.04 & 85.52 $\pm$ 1.00 & 12.79 $\pm$ 0.68 & 2.47 $\pm$ 0.07 & 5.66 $\pm$ 0.26 \\
safety+math & dare & 1.00 & 0.00 $\pm$ 0.00 & 0.57 $\pm$ 0.63 & 5.68 $\pm$ 3.61 & 55.94 $\pm$ 25.88 & 11.22 $\pm$ 0.57 & 2.35 $\pm$ 0.03 & 7.18 $\pm$ 0.22 \\
\bottomrule
\end{tabularx}
\end{table}


% --- Alpha スイープ シード別詳細 - Method: dare, Pattern: safety+math ---
\begin{table}[H]
\centering
\caption{Alpha スイープ シード別詳細 - Method: dare, Pattern: safety+math}
\label{tab:sweep_seeds_dare_safety_math}
\small
\begin{tabular}{\textwidth}{llllRRRRRRR}
\toprule
\textbf{Pattern} & \textbf{Method} & \textbf{Alpha} & \textbf{Seed} & \textbf{Safety Ave [Harmful Content] (ASR$\downarrow$ \%)} & \textbf{Math Ave (\%)} & \textbf{Code Ave (\%)} & \textbf{Medical Ave (\%)} & \textbf{General/Inst Ave (\%)} & \textbf{EvolCode (PPL$\downarrow$)} & \textbf{MedAlpaca (PPL$\downarrow$)} \\
\midrule
safety+math & dare & 0.00 & 42.0 & 32.72 & 5.94 & 2.66 & 80.62 & 34.49 & 2.02 & 2.77 \\
safety+math & dare & 0.00 & 43.0 & 32.67 & 5.00 & 3.75 & 79.69 & 25.14 & 2.03 & 2.77 \\
safety+math & dare & 0.00 & 44.0 & 34.05 & 5.47 & 3.59 & 79.06 & 31.59 & 2.02 & 2.76 \\
safety+math & dare & 0.20 & 42.0 & 11.64 & 4.06 & 4.84 & 84.38 & 34.65 & 2.32 & 3.57 \\
safety+math & dare & 0.20 & 43.0 & 10.17 & 4.22 & 5.78 & 83.44 & 24.94 & 2.37 & 3.60 \\
safety+math & dare & 0.20 & 44.0 & 8.99 & 6.09 & 3.28 & 85.00 & 24.86 & 2.40 & 3.66 \\
safety+math & dare & 0.40 & 42.0 & 11.06 & 5.16 & 5.16 & 85.31 & 21.61 & 2.32 & 3.78 \\
safety+math & dare & 0.40 & 43.0 & 8.32 & 4.84 & 5.47 & 84.38 & 25.86 & 2.39 & 3.84 \\
safety+math & dare & 0.40 & 44.0 & 7.90 & 6.09 & 4.22 & 85.31 & 23.09 & 2.42 & 3.94 \\
safety+math & dare & 0.60 & 42.0 & 0.32 & 6.25 & 5.16 & 84.38 & 15.04 & 2.35 & 4.18 \\
safety+math & dare & 0.60 & 43.0 & 8.06 & 4.69 & 5.94 & 81.88 & 21.51 & 2.46 & 4.40 \\
safety+math & dare & 0.60 & 44.0 & 10.04 & 5.31 & 3.91 & 86.56 & 32.02 & 2.48 & 4.73 \\
safety+math & dare & 0.80 & 42.0 & 0.00 & 7.03 & 5.94 & 84.38 & 13.56 & 2.40 & 5.37 \\
safety+math & dare & 0.80 & 43.0 & 9.02 & 6.41 & 6.56 & 85.94 & 12.26 & 2.51 & 5.86 \\
safety+math & dare & 0.80 & 44.0 & 1.74 & 5.62 & 4.53 & 86.25 & 12.55 & 2.52 & 5.75 \\
safety+math & dare & 1.00 & 42.0 & 0.00 & 1.25 & 9.06 & 85.62 & 11.49 & 2.32 & 7.37 \\
safety+math & dare & 1.00 & 43.0 & 0.00 & 0.47 & 1.88 & 44.06 & 10.56 & 2.37 & 7.23 \\
safety+math & dare & 1.00 & 44.0 & 0.00 & 0.00 & 6.09 & 38.12 & 11.61 & 2.36 & 6.93 \\
\bottomrule
\end{tabularx}
\end{table}


% --- Alpha スイープ 集計 - Method: dare, Pattern: safety+code ---
\begin{table}[H]
\centering
\caption{Alpha スイープ 集計 - Method: dare, Pattern: safety+code}
\label{tab:sweep_dare_safety_code}
\small
\begin{tabular}{\textwidth}{lllRRRRRRR}
\toprule
\textbf{Pattern} & \textbf{Method} & \textbf{Alpha} & \textbf{Safety Ave [Harmful Content] (ASR$\downarrow$ \%)} & \textbf{Math Ave (\%)} & \textbf{Code Ave (\%)} & \textbf{Medical Ave (\%)} & \textbf{General/Inst Ave (\%)} & \textbf{EvolCode (PPL$\downarrow$)} & \textbf{MedAlpaca (PPL$\downarrow$)} \\
\midrule
safety+code & dare & 0.00 & 0.05 $\pm$ 0.09 & 0.10 $\pm$ 0.09 & 0.00 $\pm$ 0.00 & 26.77 $\pm$ 13.01 & 11.60 $\pm$ 0.87 & 879426.80 $\pm$ 305130.82 & 609083.49 $\pm$ 167433.25 \\
safety+code & dare & 0.20 & 0.03 $\pm$ 0.05 & 0.00 $\pm$ 0.00 & 0.00 $\pm$ 0.00 & 25.21 $\pm$ 1.60 & 12.05 $\pm$ 0.25 & 421186.26 $\pm$ 138822.67 & 559487.26 $\pm$ 208077.06 \\
safety+code & dare & 0.40 & 0.03 $\pm$ 0.05 & 0.05 $\pm$ 0.09 & 0.00 $\pm$ 0.00 & 34.69 $\pm$ 36.57 & 12.05 $\pm$ 0.65 & 274915.81 $\pm$ 142634.75 & 546472.43 $\pm$ 179293.80 \\
safety+code & dare & 0.60 & 0.00 $\pm$ 0.00 & 0.05 $\pm$ 0.09 & 0.00 $\pm$ 0.00 & 29.58 $\pm$ 12.27 & 11.73 $\pm$ 0.70 & 202107.17 $\pm$ 205144.66 & 343746.27 $\pm$ 163477.65 \\
safety+code & dare & 0.80 & 0.00 $\pm$ 0.00 & 0.16 $\pm$ 0.16 & 0.00 $\pm$ 0.00 & 25.31 $\pm$ 16.83 & 11.58 $\pm$ 0.96 & 183440.18 $\pm$ 68554.99 & 437650.47 $\pm$ 213227.25 \\
safety+code & dare & 1.00 & 0.00 $\pm$ 0.00 & 0.73 $\pm$ 0.50 & 5.21 $\pm$ 2.67 & 57.81 $\pm$ 22.81 & 21.54 $\pm$ 10.74 & 2.35 $\pm$ 0.04 & 7.07 $\pm$ 0.07 \\
\bottomrule
\end{tabularx}
\end{table}


% --- Alpha スイープ シード別詳細 - Method: dare, Pattern: safety+code ---
\begin{table}[H]
\centering
\caption{Alpha スイープ シード別詳細 - Method: dare, Pattern: safety+code}
\label{tab:sweep_seeds_dare_safety_code}
\small
\begin{tabular}{\textwidth}{llllRRRRRRR}
\toprule
\textbf{Pattern} & \textbf{Method} & \textbf{Alpha} & \textbf{Seed} & \textbf{Safety Ave [Harmful Content] (ASR$\downarrow$ \%)} & \textbf{Math Ave (\%)} & \textbf{Code Ave (\%)} & \textbf{Medical Ave (\%)} & \textbf{General/Inst Ave (\%)} & \textbf{EvolCode (PPL$\downarrow$)} & \textbf{MedAlpaca (PPL$\downarrow$)} \\
\midrule
safety+code & dare & 0.00 & 42.0 & 0.00 & 0.16 & 0.00 & 40.31 & 12.14 & 1187511.19 & 491920.23 \\
safety+code & dare & 0.00 & 43.0 & 0.00 & 0.16 & 0.00 & 25.62 & 12.07 & 577337.91 & 800850.96 \\
safety+code & dare & 0.00 & 44.0 & 0.16 & 0.00 & 0.00 & 14.37 & 10.60 & 873431.29 & 534479.27 \\
safety+code & dare & 0.20 & 42.0 & 0.00 & 0.00 & 0.00 & 25.62 & 12.32 & 386759.01 & 457123.94 \\
safety+code & dare & 0.20 & 43.0 & 0.08 & 0.00 & 0.00 & 23.44 & 11.83 & 573983.09 & 798917.10 \\
safety+code & dare & 0.20 & 44.0 & 0.00 & 0.00 & 0.00 & 26.56 & 11.99 & 302816.67 & 422420.74 \\
safety+code & dare & 0.40 & 42.0 & 0.08 & 0.00 & 0.00 & 9.06 & 11.80 & 403013.26 & 566389.97 \\
safety+code & dare & 0.40 & 43.0 & 0.00 & 0.00 & 0.00 & 76.56 & 12.80 & 300521.74 & 714975.79 \\
safety+code & dare & 0.40 & 44.0 & 0.00 & 0.16 & 0.00 & 18.44 & 11.57 & 121212.45 & 358051.52 \\
safety+code & dare & 0.60 & 42.0 & 0.00 & 0.00 & 0.00 & 31.25 & 11.10 & 84728.55 & 334537.08 \\
safety+code & dare & 0.60 & 43.0 & 0.00 & 0.16 & 0.00 & 16.56 & 11.60 & 82608.30 & 185067.88 \\
safety+code & dare & 0.60 & 44.0 & 0.00 & 0.00 & 0.00 & 40.94 & 12.48 & 438984.66 & 511633.87 \\
safety+code & dare & 0.80 & 42.0 & 0.00 & 0.16 & 0.00 & 14.37 & 10.78 & 139942.12 & 559428.50 \\
safety+code & dare & 0.80 & 43.0 & 0.00 & 0.00 & 0.00 & 44.69 & 12.64 & 262466.83 & 562081.29 \\
safety+code & dare & 0.80 & 44.0 & 0.00 & 0.31 & 0.00 & 16.88 & 11.32 & 147911.59 & 191441.61 \\
safety+code & dare & 1.00 & 42.0 & 0.00 & 0.94 & 8.28 & 84.06 & 32.23 & 2.30 & 7.08 \\
safety+code & dare & 1.00 & 43.0 & 0.00 & 1.09 & 3.44 & 42.81 & 10.74 & 2.37 & 7.13 \\
safety+code & dare & 1.00 & 44.0 & 0.00 & 0.16 & 3.91 & 46.56 & 21.65 & 2.37 & 6.99 \\
\bottomrule
\end{tabularx}
\end{table}


% --- Alpha スイープ 集計 - Method: dare, Pattern: safety+medical ---
\begin{table}[H]
\centering
\caption{Alpha スイープ 集計 - Method: dare, Pattern: safety+medical}
\label{tab:sweep_dare_safety_medical}
\small
\begin{tabular}{\textwidth}{lllRRRRRRR}
\toprule
\textbf{Pattern} & \textbf{Method} & \textbf{Alpha} & \textbf{Safety Ave [Harmful Content] (ASR$\downarrow$ \%)} & \textbf{Math Ave (\%)} & \textbf{Code Ave (\%)} & \textbf{Medical Ave (\%)} & \textbf{General/Inst Ave (\%)} & \textbf{EvolCode (PPL$\downarrow$)} & \textbf{MedAlpaca (PPL$\downarrow$)} \\
\midrule
safety+medical & dare & 0.00 & 0.03 $\pm$ 0.05 & 0.00 $\pm$ 0.00 & 0.00 $\pm$ 0.00 & 51.98 $\pm$ 35.56 & 12.48 $\pm$ 0.41 & 25893511.35 $\pm$ 8569831.17 & 21157726.46 $\pm$ 5807072.10 \\
safety+medical & dare & 0.20 & 0.05 $\pm$ 0.05 & 0.00 $\pm$ 0.00 & 0.00 $\pm$ 0.00 & 4.69 $\pm$ 6.53 & 12.69 $\pm$ 1.18 & 30069390.20 $\pm$ 19622450.50 & 62005630.52 $\pm$ 57668015.31 \\
safety+medical & dare & 0.40 & 0.05 $\pm$ 0.05 & 0.00 $\pm$ 0.00 & 0.00 $\pm$ 0.00 & 12.19 $\pm$ 11.49 & 11.56 $\pm$ 0.21 & 30373008.34 $\pm$ 42093120.51 & 137500579.89 $\pm$ 224029054.92 \\
safety+medical & dare & 0.60 & 0.03 $\pm$ 0.05 & 0.00 $\pm$ 0.00 & 0.00 $\pm$ 0.00 & 22.08 $\pm$ 10.81 & 12.39 $\pm$ 0.44 & 12863372.97 $\pm$ 9716067.34 & 13423154.39 $\pm$ 11060518.87 \\
safety+medical & dare & 0.80 & 0.08 $\pm$ 0.08 & 0.00 $\pm$ 0.00 & 0.00 $\pm$ 0.00 & 16.56 $\pm$ 16.59 & 11.94 $\pm$ 0.18 & 79051784.95 $\pm$ 108121693.22 & 73992085.11 $\pm$ 102576019.23 \\
safety+medical & dare & 1.00 & 0.00 $\pm$ 0.00 & 0.47 $\pm$ 0.56 & 4.38 $\pm$ 3.53 & 62.40 $\pm$ 23.88 & 23.44 $\pm$ 10.58 & 2.34 $\pm$ 0.04 & 7.13 $\pm$ 0.18 \\
\bottomrule
\end{tabularx}
\end{table}


% --- Alpha スイープ シード別詳細 - Method: dare, Pattern: safety+medical ---
\begin{table}[H]
\centering
\caption{Alpha スイープ シード別詳細 - Method: dare, Pattern: safety+medical}
\label{tab:sweep_seeds_dare_safety_medical}
\small
\begin{tabular}{\textwidth}{llllRRRRRRR}
\toprule
\textbf{Pattern} & \textbf{Method} & \textbf{Alpha} & \textbf{Seed} & \textbf{Safety Ave [Harmful Content] (ASR$\downarrow$ \%)} & \textbf{Math Ave (\%)} & \textbf{Code Ave (\%)} & \textbf{Medical Ave (\%)} & \textbf{General/Inst Ave (\%)} & \textbf{EvolCode (PPL$\downarrow$)} & \textbf{MedAlpaca (PPL$\downarrow$)} \\
\midrule
safety+medical & dare & 0.00 & 42.0 & 0.08 & 0.00 & 0.00 & 58.13 & 12.29 & 31609289.83 & 20636305.98 \\
safety+medical & dare & 0.00 & 43.0 & 0.00 & 0.00 & 0.00 & 84.06 & 12.20 & 30031293.17 & 15628948.22 \\
safety+medical & dare & 0.00 & 44.0 & 0.00 & 0.00 & 0.00 & 13.75 & 12.94 & 16039951.05 & 27207925.18 \\
safety+medical & dare & 0.20 & 42.0 & 0.08 & 0.00 & 0.00 & 12.19 & 11.93 & 7413500.08 & 13387384.99 \\
safety+medical & dare & 0.20 & 43.0 & 0.08 & 0.00 & 0.00 & 1.56 & 14.06 & 41668528.86 & 125720332.10 \\
safety+medical & dare & 0.20 & 44.0 & 0.00 & 0.00 & 0.00 & 0.31 & 12.10 & 41126141.66 & 46909174.49 \\
safety+medical & dare & 0.40 & 42.0 & 0.08 & 0.00 & 0.00 & 0.00 & 11.35 & 4591460.66 & 6233091.27 \\
safety+medical & dare & 0.40 & 43.0 & 0.00 & 0.00 & 0.00 & 13.75 & 11.57 & 7580247.55 & 10091188.42 \\
safety+medical & dare & 0.40 & 44.0 & 0.08 & 0.00 & 0.00 & 22.81 & 11.77 & 78947316.81 & 396177459.98 \\
safety+medical & dare & 0.60 & 42.0 & 0.00 & 0.00 & 0.00 & 17.81 & 12.81 & 23273464.09 & 25783754.01 \\
safety+medical & dare & 0.60 & 43.0 & 0.08 & 0.00 & 0.00 & 14.06 & 11.93 & 4035556.82 & 10026154.84 \\
safety+medical & dare & 0.60 & 44.0 & 0.00 & 0.00 & 0.00 & 34.38 & 12.42 & 11281098.00 & 4459554.31 \\
safety+medical & dare & 0.80 & 42.0 & 0.00 & 0.00 & 0.00 & 1.56 & 12.13 & 28606774.28 & 15417965.64 \\
safety+medical & dare & 0.80 & 43.0 & 0.08 & 0.00 & 0.00 & 13.75 & 11.91 & 5371427.76 & 14123976.08 \\
safety+medical & dare & 0.80 & 44.0 & 0.16 & 0.00 & 0.00 & 34.38 & 11.78 & 203177152.83 & 192434313.61 \\
safety+medical & dare & 1.00 & 42.0 & 0.00 & 1.09 & 8.44 & 84.69 & 30.19 & 2.30 & 7.33 \\
safety+medical & dare & 1.00 & 43.0 & 0.00 & 0.31 & 2.66 & 65.31 & 28.88 & 2.36 & 7.10 \\
safety+medical & dare & 1.00 & 44.0 & 0.00 & 0.00 & 2.03 & 37.19 & 11.24 & 2.37 & 6.97 \\
\bottomrule
\end{tabularx}
\end{table}


% --- Alpha スイープ 集計 - Method: della, Pattern: safety+math+code+medical ---
\begin{table}[H]
\centering
\caption{Alpha スイープ 集計 - Method: della, Pattern: safety+math+code+medical}
\label{tab:sweep_della_safety_math_code_medical}
\small
\begin{tabular}{\textwidth}{lllRRRRRRR}
\toprule
\textbf{Pattern} & \textbf{Method} & \textbf{Alpha} & \textbf{Safety Ave [Harmful Content] (ASR$\downarrow$ \%)} & \textbf{Math Ave (\%)} & \textbf{Code Ave (\%)} & \textbf{Medical Ave (\%)} & \textbf{General/Inst Ave (\%)} & \textbf{EvolCode (PPL$\downarrow$)} & \textbf{MedAlpaca (PPL$\downarrow$)} \\
\midrule
safety+math+code+medical & della & 0.00 & 0.11 $\pm$ 0.05 & 0.00 $\pm$ 0.00 & 0.00 $\pm$ 0.00 & 46.25 $\pm$ 43.46 & 12.13 $\pm$ 0.12 & 4511054.09 $\pm$ 4014528.21 & 2709873.15 $\pm$ 1621650.35 \\
safety+math+code+medical & della & 0.20 & 0.11 $\pm$ 0.05 & 0.00 $\pm$ 0.00 & 0.00 $\pm$ 0.00 & 44.79 $\pm$ 35.50 & 11.56 $\pm$ 0.35 & 5166971.66 $\pm$ 3651027.91 & 9317429.43 $\pm$ 6260181.15 \\
safety+math+code+medical & della & 0.40 & 0.13 $\pm$ 0.05 & 0.00 $\pm$ 0.00 & 0.00 $\pm$ 0.00 & 41.46 $\pm$ 38.91 & 12.38 $\pm$ 0.49 & 3796940.54 $\pm$ 1902635.09 & 2718893.17 $\pm$ 2050580.37 \\
safety+math+code+medical & della & 0.60 & 0.08 $\pm$ 0.00 & 0.00 $\pm$ 0.00 & 0.00 $\pm$ 0.00 & 15.42 $\pm$ 16.18 & 12.53 $\pm$ 0.78 & 1556853.48 $\pm$ 450599.15 & 2115706.72 $\pm$ 117978.22 \\
safety+math+code+medical & della & 0.80 & 0.05 $\pm$ 0.09 & 0.00 $\pm$ 0.00 & 0.00 $\pm$ 0.00 & 65.52 $\pm$ 18.69 & 11.74 $\pm$ 0.15 & 4357578.11 $\pm$ 3708292.54 & 4314491.46 $\pm$ 4478659.77 \\
safety+math+code+medical & della & 1.00 & 0.00 $\pm$ 0.00 & 0.62 $\pm$ 0.31 & 4.58 $\pm$ 3.10 & 57.60 $\pm$ 23.98 & 18.90 $\pm$ 12.88 & 2.39 $\pm$ 0.04 & 7.38 $\pm$ 0.06 \\
\bottomrule
\end{tabularx}
\end{table}


% --- Alpha スイープ シード別詳細 - Method: della, Pattern: safety+math+code+medical ---
\begin{table}[H]
\centering
\caption{Alpha スイープ シード別詳細 - Method: della, Pattern: safety+math+code+medical}
\label{tab:sweep_seeds_della_safety_math_code_medical}
\small
\begin{tabular}{\textwidth}{llllRRRRRRR}
\toprule
\textbf{Pattern} & \textbf{Method} & \textbf{Alpha} & \textbf{Seed} & \textbf{Safety Ave [Harmful Content] (ASR$\downarrow$ \%)} & \textbf{Math Ave (\%)} & \textbf{Code Ave (\%)} & \textbf{Medical Ave (\%)} & \textbf{General/Inst Ave (\%)} & \textbf{EvolCode (PPL$\downarrow$)} & \textbf{MedAlpaca (PPL$\downarrow$)} \\
\midrule
safety+math+code+medical & della & 0.00 & 42.0 & 0.08 & 0.00 & 0.00 & 0.00 & 11.99 & 1076322.47 & 1937178.27 \\
safety+math+code+medical & della & 0.00 & 43.0 & 0.16 & 0.00 & 0.00 & 52.50 & 12.22 & 3532422.73 & 1619075.77 \\
safety+math+code+medical & della & 0.00 & 44.0 & 0.08 & 0.00 & 0.00 & 86.25 & 12.17 & 8924417.07 & 4573365.42 \\
safety+math+code+medical & della & 0.20 & 42.0 & 0.08 & 0.00 & 0.00 & 85.62 & 11.95 & 3367617.78 & 9340200.62 \\
safety+math+code+medical & della & 0.20 & 43.0 & 0.16 & 0.00 & 0.00 & 27.50 & 11.25 & 2764868.94 & 3045893.75 \\
safety+math+code+medical & della & 0.20 & 44.0 & 0.08 & 0.00 & 0.00 & 21.25 & 11.48 & 9368428.27 & 15566193.92 \\
safety+math+code+medical & della & 0.40 & 42.0 & 0.16 & 0.00 & 0.00 & 13.75 & 11.83 & 5819408.74 & 5081745.67 \\
safety+math+code+medical & della & 0.40 & 43.0 & 0.16 & 0.00 & 0.00 & 85.94 & 12.52 & 2042579.89 & 1404892.78 \\
safety+math+code+medical & della & 0.40 & 44.0 & 0.08 & 0.00 & 0.00 & 24.69 & 12.78 & 3528832.98 & 1670041.07 \\
safety+math+code+medical & della & 0.60 & 42.0 & 0.08 & 0.00 & 0.00 & 13.44 & 11.64 & 1547222.80 & 1979585.04 \\
safety+math+code+medical & della & 0.60 & 43.0 & 0.08 & 0.00 & 0.00 & 0.31 & 12.97 & 1111146.86 & 2188460.85 \\
safety+math+code+medical & della & 0.60 & 44.0 & 0.08 & 0.00 & 0.00 & 32.50 & 13.00 & 2012190.77 & 2179074.27 \\
safety+math+code+medical & della & 0.80 & 42.0 & 0.16 & 0.00 & 0.00 & 81.56 & 11.58 & 5216210.59 & 9270274.79 \\
safety+math+code+medical & della & 0.80 & 43.0 & 0.00 & 0.00 & 0.00 & 70.00 & 11.75 & 7561235.35 & 3116659.19 \\
safety+math+code+medical & della & 0.80 & 44.0 & 0.00 & 0.00 & 0.00 & 45.00 & 11.89 & 295288.40 & 556540.40 \\
safety+math+code+medical & della & 1.00 & 42.0 & 0.00 & 0.94 & 8.12 & 84.69 & 33.77 & 2.34 & 7.44 \\
safety+math+code+medical & della & 1.00 & 43.0 & 0.00 & 0.62 & 2.34 & 49.06 & 11.21 & 2.41 & 7.34 \\
safety+math+code+medical & della & 1.00 & 44.0 & 0.00 & 0.31 & 3.28 & 39.06 & 11.72 & 2.42 & 7.35 \\
\bottomrule
\end{tabularx}
\end{table}


% --- Alpha スイープ 集計 - Method: della, Pattern: safety+math ---
\begin{table}[H]
\centering
\caption{Alpha スイープ 集計 - Method: della, Pattern: safety+math}
\label{tab:sweep_della_safety_math}
\small
\begin{tabular}{\textwidth}{lllRRRRRRR}
\toprule
\textbf{Pattern} & \textbf{Method} & \textbf{Alpha} & \textbf{Safety Ave [Harmful Content] (ASR$\downarrow$ \%)} & \textbf{Math Ave (\%)} & \textbf{Code Ave (\%)} & \textbf{Medical Ave (\%)} & \textbf{General/Inst Ave (\%)} & \textbf{EvolCode (PPL$\downarrow$)} & \textbf{MedAlpaca (PPL$\downarrow$)} \\
\midrule
safety+math & della & 0.00 & 33.26 $\pm$ 0.87 & 6.56 $\pm$ 0.16 & 3.33 $\pm$ 0.59 & 83.44 $\pm$ 1.74 & 34.20 $\pm$ 1.62 & 2.04 $\pm$ 0.00 & 2.78 $\pm$ 0.02 \\
safety+math & della & 0.20 & 11.63 $\pm$ 0.83 & 5.57 $\pm$ 0.89 & 4.64 $\pm$ 0.24 & 83.23 $\pm$ 1.18 & 32.96 $\pm$ 2.11 & 2.38 $\pm$ 0.04 & 3.70 $\pm$ 0.07 \\
safety+math & della & 0.40 & 7.61 $\pm$ 2.70 & 5.52 $\pm$ 0.09 & 5.31 $\pm$ 0.68 & 84.58 $\pm$ 2.03 & 32.47 $\pm$ 2.15 & 2.41 $\pm$ 0.05 & 3.93 $\pm$ 0.10 \\
safety+math & della & 0.60 & 6.66 $\pm$ 5.52 & 5.52 $\pm$ 0.63 & 5.31 $\pm$ 1.36 & 85.62 $\pm$ 0.54 & 31.03 $\pm$ 4.06 & 2.46 $\pm$ 0.06 & 4.50 $\pm$ 0.30 \\
safety+math & della & 0.80 & 2.71 $\pm$ 3.71 & 5.78 $\pm$ 1.08 & 4.90 $\pm$ 0.39 & 85.52 $\pm$ 0.79 & 19.17 $\pm$ 5.48 & 2.51 $\pm$ 0.09 & 5.81 $\pm$ 0.42 \\
safety+math & della & 1.00 & 0.00 $\pm$ 0.00 & 0.62 $\pm$ 0.41 & 4.69 $\pm$ 3.29 & 58.54 $\pm$ 23.27 & 17.78 $\pm$ 11.72 & 2.38 $\pm$ 0.03 & 7.33 $\pm$ 0.14 \\
\bottomrule
\end{tabularx}
\end{table}


% --- Alpha スイープ シード別詳細 - Method: della, Pattern: safety+math ---
\begin{table}[H]
\centering
\caption{Alpha スイープ シード別詳細 - Method: della, Pattern: safety+math}
\label{tab:sweep_seeds_della_safety_math}
\small
\begin{tabular}{\textwidth}{llllRRRRRRR}
\toprule
\textbf{Pattern} & \textbf{Method} & \textbf{Alpha} & \textbf{Seed} & \textbf{Safety Ave [Harmful Content] (ASR$\downarrow$ \%)} & \textbf{Math Ave (\%)} & \textbf{Code Ave (\%)} & \textbf{Medical Ave (\%)} & \textbf{General/Inst Ave (\%)} & \textbf{EvolCode (PPL$\downarrow$)} & \textbf{MedAlpaca (PPL$\downarrow$)} \\
\midrule
safety+math & della & 0.00 & 42.0 & 34.22 & 6.72 & 3.59 & 81.88 & 34.10 & 2.05 & 2.80 \\
safety+math & della & 0.00 & 43.0 & 32.52 & 6.56 & 3.75 & 83.12 & 35.87 & 2.05 & 2.76 \\
safety+math & della & 0.00 & 44.0 & 33.04 & 6.41 & 2.66 & 85.31 & 32.63 & 2.04 & 2.78 \\
safety+math & della & 0.20 & 42.0 & 10.84 & 4.84 & 4.84 & 84.06 & 31.09 & 2.34 & 3.66 \\
safety+math & della & 0.20 & 43.0 & 12.50 & 5.31 & 4.38 & 81.88 & 35.25 & 2.39 & 3.67 \\
safety+math & della & 0.20 & 44.0 & 11.53 & 6.56 & 4.69 & 83.75 & 32.54 & 2.41 & 3.78 \\
safety+math & della & 0.40 & 42.0 & 10.66 & 5.47 & 6.09 & 84.69 & 31.22 & 2.35 & 3.83 \\
safety+math & della & 0.40 & 43.0 & 5.56 & 5.62 & 5.00 & 82.50 & 34.96 & 2.41 & 3.91 \\
safety+math & della & 0.40 & 44.0 & 6.61 & 5.47 & 4.84 & 86.56 & 31.25 & 2.46 & 4.04 \\
safety+math & della & 0.60 & 42.0 & 0.63 & 4.84 & 6.25 & 85.31 & 29.00 & 2.39 & 4.26 \\
safety+math & della & 0.60 & 43.0 & 11.45 & 5.62 & 5.94 & 85.31 & 28.39 & 2.47 & 4.41 \\
safety+math & della & 0.60 & 44.0 & 7.92 & 6.09 & 3.75 & 86.25 & 35.70 & 2.51 & 4.84 \\
safety+math & della & 0.80 & 42.0 & 0.00 & 7.03 & 4.84 & 84.69 & 23.57 & 2.42 & 5.34 \\
safety+math & della & 0.80 & 43.0 & 6.94 & 5.16 & 5.31 & 85.62 & 20.90 & 2.54 & 5.98 \\
safety+math & della & 0.80 & 44.0 & 1.19 & 5.16 & 4.53 & 86.25 & 13.03 & 2.58 & 6.12 \\
safety+math & della & 1.00 & 42.0 & 0.00 & 0.78 & 7.81 & 85.00 & 31.29 & 2.35 & 7.45 \\
safety+math & della & 1.00 & 43.0 & 0.00 & 0.94 & 1.25 & 49.38 & 10.38 & 2.41 & 7.37 \\
safety+math & della & 1.00 & 44.0 & 0.00 & 0.16 & 5.00 & 41.25 & 11.68 & 2.40 & 7.17 \\
\bottomrule
\end{tabularx}
\end{table}


% --- Alpha スイープ 集計 - Method: della, Pattern: safety+code ---
\begin{table}[H]
\centering
\caption{Alpha スイープ 集計 - Method: della, Pattern: safety+code}
\label{tab:sweep_della_safety_code}
\small
\begin{tabular}{\textwidth}{lllRRRRRRR}
\toprule
\textbf{Pattern} & \textbf{Method} & \textbf{Alpha} & \textbf{Safety Ave [Harmful Content] (ASR$\downarrow$ \%)} & \textbf{Math Ave (\%)} & \textbf{Code Ave (\%)} & \textbf{Medical Ave (\%)} & \textbf{General/Inst Ave (\%)} & \textbf{EvolCode (PPL$\downarrow$)} & \textbf{MedAlpaca (PPL$\downarrow$)} \\
\midrule
safety+code & della & 0.00 & 0.00 $\pm$ 0.00 & 0.05 $\pm$ 0.09 & 0.00 $\pm$ 0.00 & 41.04 $\pm$ 10.74 & 11.94 $\pm$ 0.59 & 569125.61 $\pm$ 394534.46 & 1151826.80 $\pm$ 842498.02 \\
safety+code & della & 0.20 & 0.00 $\pm$ 0.00 & 0.52 $\pm$ 0.33 & 0.00 $\pm$ 0.00 & 22.08 $\pm$ 3.31 & 11.43 $\pm$ 0.95 & 212820.11 $\pm$ 119967.94 & 320482.73 $\pm$ 211524.60 \\
safety+code & della & 0.40 & 0.05 $\pm$ 0.05 & 0.05 $\pm$ 0.09 & 0.00 $\pm$ 0.00 & 21.15 $\pm$ 10.37 & 11.39 $\pm$ 0.22 & 199062.51 $\pm$ 165620.73 & 394587.54 $\pm$ 221751.31 \\
safety+code & della & 0.60 & 0.00 $\pm$ 0.00 & 0.62 $\pm$ 0.41 & 0.00 $\pm$ 0.00 & 43.54 $\pm$ 14.08 & 12.43 $\pm$ 0.46 & 107944.55 $\pm$ 59575.64 & 211408.90 $\pm$ 76715.67 \\
safety+code & della & 0.80 & 0.05 $\pm$ 0.09 & 0.62 $\pm$ 0.83 & 0.00 $\pm$ 0.00 & 35.10 $\pm$ 27.13 & 12.00 $\pm$ 0.78 & 63830.50 $\pm$ 29959.41 & 217419.43 $\pm$ 83074.00 \\
safety+code & della & 1.00 & 0.00 $\pm$ 0.00 & 0.62 $\pm$ 0.72 & 4.43 $\pm$ 3.60 & 53.54 $\pm$ 27.30 & 11.04 $\pm$ 0.62 & 2.38 $\pm$ 0.04 & 7.38 $\pm$ 0.07 \\
\bottomrule
\end{tabularx}
\end{table}


% --- Alpha スイープ シード別詳細 - Method: della, Pattern: safety+code ---
\begin{table}[H]
\centering
\caption{Alpha スイープ シード別詳細 - Method: della, Pattern: safety+code}
\label{tab:sweep_seeds_della_safety_code}
\small
\begin{tabular}{\textwidth}{llllRRRRRRR}
\toprule
\textbf{Pattern} & \textbf{Method} & \textbf{Alpha} & \textbf{Seed} & \textbf{Safety Ave [Harmful Content] (ASR$\downarrow$ \%)} & \textbf{Math Ave (\%)} & \textbf{Code Ave (\%)} & \textbf{Medical Ave (\%)} & \textbf{General/Inst Ave (\%)} & \textbf{EvolCode (PPL$\downarrow$)} & \textbf{MedAlpaca (PPL$\downarrow$)} \\
\midrule
safety+code & della & 0.00 & 42.0 & 0.00 & 0.00 & 0.00 & 35.00 & 11.85 & 318866.94 & 1074126.52 \\
safety+code & della & 0.00 & 43.0 & 0.00 & 0.00 & 0.00 & 53.44 & 11.40 & 364580.27 & 350870.46 \\
safety+code & della & 0.00 & 44.0 & 0.00 & 0.16 & 0.00 & 34.69 & 12.56 & 1023929.61 & 2030483.41 \\
safety+code & della & 0.20 & 42.0 & 0.00 & 0.62 & 0.00 & 25.62 & 12.29 & 284396.34 & 501382.55 \\
safety+code & della & 0.20 & 43.0 & 0.00 & 0.78 & 0.00 & 19.06 & 10.42 & 74319.10 & 87909.29 \\
safety+code & della & 0.20 & 44.0 & 0.00 & 0.16 & 0.00 & 21.56 & 11.58 & 279744.90 & 372156.36 \\
safety+code & della & 0.40 & 42.0 & 0.08 & 0.00 & 0.00 & 12.19 & 11.65 & 386282.25 & 378643.94 \\
safety+code & della & 0.40 & 43.0 & 0.00 & 0.00 & 0.00 & 18.75 & 11.26 & 139243.22 & 623880.35 \\
safety+code & della & 0.40 & 44.0 & 0.08 & 0.16 & 0.00 & 32.50 & 11.28 & 71662.06 & 181238.32 \\
safety+code & della & 0.60 & 42.0 & 0.00 & 0.94 & 0.00 & 57.19 & 12.96 & 173731.57 & 126277.55 \\
safety+code & della & 0.60 & 43.0 & 0.00 & 0.16 & 0.00 & 29.06 & 12.10 & 92466.80 & 275182.68 \\
safety+code & della & 0.60 & 44.0 & 0.00 & 0.78 & 0.00 & 44.38 & 12.23 & 57635.29 & 232766.47 \\
safety+code & della & 0.80 & 42.0 & 0.16 & 1.56 & 0.00 & 13.75 & 12.58 & 52116.52 & 262634.90 \\
safety+code & della & 0.80 & 43.0 & 0.00 & 0.00 & 0.00 & 65.62 & 12.30 & 41497.89 & 121545.33 \\
safety+code & della & 0.80 & 44.0 & 0.00 & 0.31 & 0.00 & 25.94 & 11.12 & 97877.08 & 268078.07 \\
safety+code & della & 1.00 & 42.0 & 0.00 & 1.41 & 8.12 & 84.69 & 11.16 & 2.34 & 7.41 \\
safety+code & della & 1.00 & 43.0 & 0.00 & 0.47 & 0.94 & 42.19 & 10.37 & 2.40 & 7.43 \\
safety+code & della & 1.00 & 44.0 & 0.00 & 0.00 & 4.22 & 33.75 & 11.59 & 2.40 & 7.29 \\
\bottomrule
\end{tabularx}
\end{table}


% --- Alpha スイープ 集計 - Method: della, Pattern: safety+medical ---
\begin{table}[H]
\centering
\caption{Alpha スイープ 集計 - Method: della, Pattern: safety+medical}
\label{tab:sweep_della_safety_medical}
\small
\begin{tabular}{\textwidth}{lllRRRRRRR}
\toprule
\textbf{Pattern} & \textbf{Method} & \textbf{Alpha} & \textbf{Safety Ave [Harmful Content] (ASR$\downarrow$ \%)} & \textbf{Math Ave (\%)} & \textbf{Code Ave (\%)} & \textbf{Medical Ave (\%)} & \textbf{General/Inst Ave (\%)} & \textbf{EvolCode (PPL$\downarrow$)} & \textbf{MedAlpaca (PPL$\downarrow$)} \\
\midrule
safety+medical & della & 0.00 & 0.08 $\pm$ 0.08 & 0.00 $\pm$ 0.00 & 0.00 $\pm$ 0.00 & 74.90 $\pm$ 18.85 & 12.24 $\pm$ 0.24 & 26812751.49 $\pm$ 26729427.27 & 43518202.81 $\pm$ 44548820.85 \\
safety+medical & della & 0.20 & 0.03 $\pm$ 0.05 & 0.00 $\pm$ 0.00 & 0.00 $\pm$ 0.00 & 35.94 $\pm$ 37.89 & 11.96 $\pm$ 0.91 & 7535389.77 $\pm$ 3979386.81 & 9534444.12 $\pm$ 6845498.98 \\
safety+medical & della & 0.40 & 0.05 $\pm$ 0.09 & 0.00 $\pm$ 0.00 & 0.00 $\pm$ 0.00 & 18.23 $\pm$ 16.24 & 12.48 $\pm$ 0.36 & 24905029231.27 $\pm$ 43111038373.55 & 2796226504.59 $\pm$ 4644885775.60 \\
safety+medical & della & 0.60 & 0.03 $\pm$ 0.05 & 0.00 $\pm$ 0.00 & 0.00 $\pm$ 0.00 & 27.29 $\pm$ 23.45 & 12.10 $\pm$ 1.01 & 10754022.73 $\pm$ 9470200.24 & 8631379.06 $\pm$ 4139779.72 \\
safety+medical & della & 0.80 & 0.05 $\pm$ 0.05 & 0.00 $\pm$ 0.00 & 0.00 $\pm$ 0.00 & 17.08 $\pm$ 16.20 & 12.07 $\pm$ 0.80 & 40599833.68 $\pm$ 51603645.67 & 33461048.41 $\pm$ 46722626.90 \\
safety+medical & della & 1.00 & 0.00 $\pm$ 0.00 & 0.62 $\pm$ 0.68 & 4.01 $\pm$ 3.43 & 56.56 $\pm$ 22.05 & 17.08 $\pm$ 10.75 & 2.38 $\pm$ 0.04 & 7.38 $\pm$ 0.14 \\
\bottomrule
\end{tabularx}
\end{table}


% --- Alpha スイープ シード別詳細 - Method: della, Pattern: safety+medical ---
\begin{table}[H]
\centering
\caption{Alpha スイープ シード別詳細 - Method: della, Pattern: safety+medical}
\label{tab:sweep_seeds_della_safety_medical}
\small
\begin{tabular}{\textwidth}{llllRRRRRRR}
\toprule
\textbf{Pattern} & \textbf{Method} & \textbf{Alpha} & \textbf{Seed} & \textbf{Safety Ave [Harmful Content] (ASR$\downarrow$ \%)} & \textbf{Math Ave (\%)} & \textbf{Code Ave (\%)} & \textbf{Medical Ave (\%)} & \textbf{General/Inst Ave (\%)} & \textbf{EvolCode (PPL$\downarrow$)} & \textbf{MedAlpaca (PPL$\downarrow$)} \\
\midrule
safety+medical & della & 0.00 & 42.0 & 0.00 & 0.00 & 0.00 & 85.62 & 12.22 & 10500490.26 & 26900946.82 \\
safety+medical & della & 0.00 & 43.0 & 0.08 & 0.00 & 0.00 & 53.12 & 12.02 & 12277586.79 & 9666451.99 \\
safety+medical & della & 0.00 & 44.0 & 0.16 & 0.00 & 0.00 & 85.94 & 12.50 & 57660177.40 & 93987209.62 \\
safety+medical & della & 0.20 & 42.0 & 0.00 & 0.00 & 0.00 & 13.75 & 11.12 & 7593960.41 & 16608425.88 \\
safety+medical & della & 0.20 & 43.0 & 0.08 & 0.00 & 0.00 & 14.37 & 11.83 & 3527040.92 & 2942957.72 \\
safety+medical & della & 0.20 & 44.0 & 0.00 & 0.00 & 0.00 & 79.69 & 12.94 & 11485167.97 & 9051948.75 \\
safety+medical & della & 0.40 & 42.0 & 0.16 & 0.00 & 0.00 & 5.62 & 12.29 & 15375847.04 & 223089326.53 \\
safety+medical & della & 0.40 & 43.0 & 0.00 & 0.00 & 0.00 & 12.50 & 12.89 & 74685368447.73 & 8158232207.05 \\
safety+medical & della & 0.40 & 44.0 & 0.00 & 0.00 & 0.00 & 36.56 & 12.25 & 14343399.05 & 7357980.19 \\
safety+medical & della & 0.60 & 42.0 & 0.08 & 0.00 & 0.00 & 13.75 & 12.84 & 21562404.57 & 13042229.15 \\
safety+medical & della & 0.60 & 43.0 & 0.00 & 0.00 & 0.00 & 54.37 & 12.52 & 6788182.90 & 4830310.81 \\
safety+medical & della & 0.60 & 44.0 & 0.00 & 0.00 & 0.00 & 13.75 & 10.95 & 3911480.70 & 8021597.21 \\
safety+medical & della & 0.80 & 42.0 & 0.00 & 0.00 & 0.00 & 34.69 & 12.60 & 6850838.62 & 6993248.65 \\
safety+medical & della & 0.80 & 43.0 & 0.08 & 0.00 & 0.00 & 2.81 & 12.45 & 100003030.30 & 87408527.68 \\
safety+medical & della & 0.80 & 44.0 & 0.08 & 0.00 & 0.00 & 13.75 & 11.15 & 14945632.14 & 5981368.91 \\
safety+medical & della & 1.00 & 42.0 & 0.00 & 1.41 & 7.97 & 81.88 & 29.49 & 2.33 & 7.51 \\
safety+medical & della & 1.00 & 43.0 & 0.00 & 0.31 & 2.19 & 46.25 & 10.77 & 2.39 & 7.41 \\
safety+medical & della & 1.00 & 44.0 & 0.00 & 0.16 & 1.88 & 41.56 & 10.97 & 2.41 & 7.23 \\
\bottomrule
\end{tabularx}
\end{table}


% --- Alpha スイープ 集計 - Method: data_free_sst_main, Pattern: safety+math+code+medical ---
\begin{table}[H]
\centering
\caption{Alpha スイープ 集計 - Method: data_free_sst_main, Pattern: safety+math+code+medical}
\label{tab:sweep_datafreesstmain_safety_math_code_medical}
\small
\begin{tabular}{\textwidth}{lllRRRRRRR}
\toprule
\textbf{Pattern} & \textbf{Method} & \textbf{Alpha} & \textbf{Safety Ave [Harmful Content] (ASR$\downarrow$ \%)} & \textbf{Math Ave (\%)} & \textbf{Code Ave (\%)} & \textbf{Medical Ave (\%)} & \textbf{General/Inst Ave (\%)} & \textbf{EvolCode (PPL$\downarrow$)} & \textbf{MedAlpaca (PPL$\downarrow$)} \\
\midrule
safety+math+code+medical & data\_free\_sst\_main & 0.00 & 0.16 $\pm$ 0.00 & 0.00 $\pm$ 0.00 & 0.00 $\pm$ 0.00 & 12.19 $\pm$ 0.00 & 10.65 $\pm$ 0.12 & 48514.43 $\pm$ 0.00 & 60540.72 $\pm$ 0.00 \\
safety+math+code+medical & data\_free\_sst\_main & 0.20 & 0.13 $\pm$ 0.09 & 0.00 $\pm$ 0.00 & 0.00 $\pm$ 0.00 & 0.62 $\pm$ 0.83 & 11.32 $\pm$ 0.24 & 23973.29 $\pm$ 1724.29 & 42451.04 $\pm$ 921.18 \\
safety+math+code+medical & data\_free\_sst\_main & 0.40 & 0.00 $\pm$ 0.00 & 0.00 $\pm$ 0.00 & 0.00 $\pm$ 0.00 & 6.25 $\pm$ 6.25 & 12.00 $\pm$ 0.20 & 12564.36 $\pm$ 2806.77 & 25747.98 $\pm$ 3655.02 \\
safety+math+code+medical & data\_free\_sst\_main & 0.60 & 0.08 $\pm$ 0.00 & 0.00 $\pm$ 0.00 & 0.00 $\pm$ 0.00 & 26.67 $\pm$ 24.56 & 11.85 $\pm$ 0.59 & 8840.75 $\pm$ 1530.06 & 17110.70 $\pm$ 5533.52 \\
safety+math+code+medical & data\_free\_sst\_main & 0.80 & 0.03 $\pm$ 0.05 & 0.00 $\pm$ 0.00 & 0.00 $\pm$ 0.00 & 57.71 $\pm$ 27.97 & 12.37 $\pm$ 1.27 & 6526.30 $\pm$ 426.62 & 7716.72 $\pm$ 2112.88 \\
safety+math+code+medical & data\_free\_sst\_main & 1.00 & 0.03 $\pm$ 0.05 & 0.00 $\pm$ 0.00 & 0.00 $\pm$ 0.00 & 62.40 $\pm$ 27.79 & 11.48 $\pm$ 0.20 & 6377.05 $\pm$ 1134.31 & 5459.66 $\pm$ 490.26 \\
\bottomrule
\end{tabularx}
\end{table}


% --- Alpha スイープ シード別詳細 - Method: data_free_sst_main, Pattern: safety+math+code+medical ---
\begin{table}[H]
\centering
\caption{Alpha スイープ シード別詳細 - Method: data_free_sst_main, Pattern: safety+math+code+medical}
\label{tab:sweep_seeds_datafreesstmain_safety_math_code_medical}
\small
\begin{tabular}{\textwidth}{llllRRRRRRR}
\toprule
\textbf{Pattern} & \textbf{Method} & \textbf{Alpha} & \textbf{Seed} & \textbf{Safety Ave [Harmful Content] (ASR$\downarrow$ \%)} & \textbf{Math Ave (\%)} & \textbf{Code Ave (\%)} & \textbf{Medical Ave (\%)} & \textbf{General/Inst Ave (\%)} & \textbf{EvolCode (PPL$\downarrow$)} & \textbf{MedAlpaca (PPL$\downarrow$)} \\
\midrule
safety+math+code+medical & data\_free\_sst\_main & 0.00 & 42.0 & 0.16 & 0.00 & 0.00 & 12.19 & 10.79 & 48514.43 & 60540.72 \\
safety+math+code+medical & data\_free\_sst\_main & 0.00 & 43.0 & 0.16 & 0.00 & 0.00 & 12.19 & 10.58 & 48514.43 & 60540.72 \\
safety+math+code+medical & data\_free\_sst\_main & 0.00 & 44.0 & 0.16 & 0.00 & 0.00 & 12.19 & 10.58 & 48514.43 & 60540.72 \\
safety+math+code+medical & data\_free\_sst\_main & 0.20 & 42.0 & 0.08 & 0.00 & 0.00 & 1.56 & 11.07 & 24088.01 & 43085.80 \\
safety+math+code+medical & data\_free\_sst\_main & 0.20 & 43.0 & 0.08 & 0.00 & 0.00 & 0.00 & 11.55 & 22194.51 & 41394.49 \\
safety+math+code+medical & data\_free\_sst\_main & 0.20 & 44.0 & 0.24 & 0.00 & 0.00 & 0.31 & 11.35 & 25637.36 & 42872.84 \\
safety+math+code+medical & data\_free\_sst\_main & 0.40 & 42.0 & 0.00 & 0.00 & 0.00 & 6.25 & 12.00 & 12708.72 & 29794.72 \\
safety+math+code+medical & data\_free\_sst\_main & 0.40 & 43.0 & 0.00 & 0.00 & 0.00 & 0.00 & 11.80 & 9688.20 & 22686.80 \\
safety+math+code+medical & data\_free\_sst\_main & 0.40 & 44.0 & 0.00 & 0.00 & 0.00 & 12.50 & 12.20 & 15296.16 & 24762.41 \\
safety+math+code+medical & data\_free\_sst\_main & 0.60 & 42.0 & 0.08 & 0.00 & 0.00 & 13.44 & 12.53 & 8906.95 & 19263.61 \\
safety+math+code+medical & data\_free\_sst\_main & 0.60 & 43.0 & 0.08 & 0.00 & 0.00 & 11.56 & 11.50 & 7278.66 & 10824.29 \\
safety+math+code+medical & data\_free\_sst\_main & 0.60 & 44.0 & 0.08 & 0.00 & 0.00 & 55.00 & 11.51 & 10336.63 & 21244.18 \\
safety+math+code+medical & data\_free\_sst\_main & 0.80 & 42.0 & 0.00 & 0.00 & 0.00 & 30.00 & 13.79 & 6338.87 & 8330.75 \\
safety+math+code+medical & data\_free\_sst\_main & 0.80 & 43.0 & 0.08 & 0.00 & 0.00 & 57.19 & 11.99 & 7014.55 & 5364.83 \\
safety+math+code+medical & data\_free\_sst\_main & 0.80 & 44.0 & 0.00 & 0.00 & 0.00 & 85.94 & 11.34 & 6225.47 & 9454.56 \\
safety+math+code+medical & data\_free\_sst\_main & 1.00 & 42.0 & 0.08 & 0.00 & 0.00 & 31.87 & 11.30 & 7255.67 & 4971.49 \\
safety+math+code+medical & data\_free\_sst\_main & 1.00 & 43.0 & 0.00 & 0.00 & 0.00 & 69.06 & 11.69 & 6778.98 & 5455.51 \\
safety+math+code+medical & data\_free\_sst\_main & 1.00 & 44.0 & 0.00 & 0.00 & 0.00 & 86.25 & 11.44 & 5096.50 & 5951.99 \\
\bottomrule
\end{tabularx}
\end{table}


% --- Alpha スイープ 集計 - Method: data_free_sst_main, Pattern: safety+math ---
\begin{table}[H]
\centering
\caption{Alpha スイープ 集計 - Method: data_free_sst_main, Pattern: safety+math}
\label{tab:sweep_datafreesstmain_safety_math}
\small
\begin{tabular}{\textwidth}{lllRRRRRRR}
\toprule
\textbf{Pattern} & \textbf{Method} & \textbf{Alpha} & \textbf{Safety Ave [Harmful Content] (ASR$\downarrow$ \%)} & \textbf{Math Ave (\%)} & \textbf{Code Ave (\%)} & \textbf{Medical Ave (\%)} & \textbf{General/Inst Ave (\%)} & \textbf{EvolCode (PPL$\downarrow$)} & \textbf{MedAlpaca (PPL$\downarrow$)} \\
\midrule
safety+math & data\_free\_sst\_main & 0.00 & 34.32 $\pm$ 0.20 & 5.99 $\pm$ 0.09 & 2.92 $\pm$ 1.26 & 63.44 $\pm$ 28.69 & 30.93 $\pm$ 7.17 & 1.99 $\pm$ 0.00 & 2.71 $\pm$ 0.00 \\
safety+math & data\_free\_sst\_main & 0.20 & 32.36 $\pm$ 0.85 & 3.96 $\pm$ 2.07 & 4.84 $\pm$ 2.30 & 70.57 $\pm$ 14.22 & 27.02 $\pm$ 6.95 & 1.99 $\pm$ 0.01 & 2.72 $\pm$ 0.01 \\
safety+math & data\_free\_sst\_main & 0.40 & 35.94 $\pm$ 4.77 & 5.26 $\pm$ 0.18 & 6.04 $\pm$ 3.16 & 62.60 $\pm$ 28.36 & 28.39 $\pm$ 3.75 & 1.99 $\pm$ 0.01 & 2.80 $\pm$ 0.02 \\
safety+math & data\_free\_sst\_main & 0.60 & 14.73 $\pm$ 3.41 & 5.57 $\pm$ 0.24 & 7.19 $\pm$ 4.33 & 69.38 $\pm$ 13.52 & 24.69 $\pm$ 0.38 & 2.02 $\pm$ 0.02 & 3.02 $\pm$ 0.02 \\
safety+math & data\_free\_sst\_main & 0.80 & 15.59 $\pm$ 6.01 & 4.90 $\pm$ 0.86 & 5.62 $\pm$ 0.83 & 66.72 $\pm$ 13.56 & 26.05 $\pm$ 4.68 & 2.08 $\pm$ 0.03 & 3.55 $\pm$ 0.03 \\
safety+math & data\_free\_sst\_main & 1.00 & 11.49 $\pm$ 9.25 & 5.89 $\pm$ 0.79 & 5.57 $\pm$ 0.24 & 56.77 $\pm$ 26.43 & 21.14 $\pm$ 7.66 & 2.19 $\pm$ 0.04 & 4.64 $\pm$ 0.10 \\
\bottomrule
\end{tabularx}
\end{table}


% --- Alpha スイープ シード別詳細 - Method: data_free_sst_main, Pattern: safety+math ---
\begin{table}[H]
\centering
\caption{Alpha スイープ シード別詳細 - Method: data_free_sst_main, Pattern: safety+math}
\label{tab:sweep_seeds_datafreesstmain_safety_math}
\small
\begin{tabular}{\textwidth}{llllRRRRRRR}
\toprule
\textbf{Pattern} & \textbf{Method} & \textbf{Alpha} & \textbf{Seed} & \textbf{Safety Ave [Harmful Content] (ASR$\downarrow$ \%)} & \textbf{Math Ave (\%)} & \textbf{Code Ave (\%)} & \textbf{Medical Ave (\%)} & \textbf{General/Inst Ave (\%)} & \textbf{EvolCode (PPL$\downarrow$)} & \textbf{MedAlpaca (PPL$\downarrow$)} \\
\midrule
safety+math & data\_free\_sst\_main & 0.00 & 42.0 & 34.43 & 6.09 & 4.38 & 30.31 & 39.21 & 1.99 & 2.71 \\
safety+math & data\_free\_sst\_main & 0.00 & 43.0 & 34.10 & 5.94 & 2.19 & 80.00 & 26.72 & 1.99 & 2.71 \\
safety+math & data\_free\_sst\_main & 0.00 & 44.0 & 34.44 & 5.94 & 2.19 & 80.00 & 26.88 & 1.99 & 2.71 \\
safety+math & data\_free\_sst\_main & 0.20 & 42.0 & 33.03 & 1.56 & 7.50 & 54.22 & 26.52 & 1.98 & 2.71 \\
safety+math & data\_free\_sst\_main & 0.20 & 43.0 & 32.65 & 5.16 & 3.44 & 77.50 & 34.21 & 1.99 & 2.73 \\
safety+math & data\_free\_sst\_main & 0.20 & 44.0 & 31.41 & 5.16 & 3.59 & 80.00 & 20.33 & 1.99 & 2.73 \\
safety+math & data\_free\_sst\_main & 0.40 & 42.0 & 41.42 & 5.16 & 9.69 & 30.00 & 26.62 & 1.98 & 2.78 \\
safety+math & data\_free\_sst\_main & 0.40 & 43.0 & 32.71 & 5.16 & 4.06 & 76.25 & 25.86 & 2.00 & 2.79 \\
safety+math & data\_free\_sst\_main & 0.40 & 44.0 & 33.70 & 5.47 & 4.38 & 81.56 & 32.70 & 2.00 & 2.82 \\
safety+math & data\_free\_sst\_main & 0.60 & 42.0 & 10.89 & 5.31 & 12.19 & 54.06 & 24.30 & 1.99 & 3.03 \\
safety+math & data\_free\_sst\_main & 0.60 & 43.0 & 15.92 & 5.78 & 4.84 & 74.38 & 25.07 & 2.02 & 3.04 \\
safety+math & data\_free\_sst\_main & 0.60 & 44.0 & 17.39 & 5.62 & 4.53 & 79.69 & 24.70 & 2.04 & 3.00 \\
safety+math & data\_free\_sst\_main & 0.80 & 42.0 & 9.86 & 4.06 & 6.56 & 51.41 & 31.46 & 2.05 & 3.55 \\
safety+math & data\_free\_sst\_main & 0.80 & 43.0 & 15.06 & 5.78 & 5.00 & 71.56 & 23.39 & 2.10 & 3.51 \\
safety+math & data\_free\_sst\_main & 0.80 & 44.0 & 21.85 & 4.84 & 5.31 & 77.19 & 23.31 & 2.10 & 3.58 \\
safety+math & data\_free\_sst\_main & 1.00 & 42.0 & 0.95 & 5.78 & 5.31 & 28.75 & 13.26 & 2.14 & 4.52 \\
safety+math & data\_free\_sst\_main & 1.00 & 43.0 & 18.26 & 6.72 & 5.78 & 60.31 & 28.55 & 2.23 & 4.71 \\
safety+math & data\_free\_sst\_main & 1.00 & 44.0 & 15.26 & 5.16 & 5.62 & 81.25 & 21.63 & 2.21 & 4.69 \\
\bottomrule
\end{tabularx}
\end{table}


% --- Alpha スイープ 集計 - Method: data_free_sst_main, Pattern: safety+code ---
\begin{table}[H]
\centering
\caption{Alpha スイープ 集計 - Method: data_free_sst_main, Pattern: safety+code}
\label{tab:sweep_datafreesstmain_safety_code}
\small
\begin{tabular}{\textwidth}{lllRRRRRRR}
\toprule
\textbf{Pattern} & \textbf{Method} & \textbf{Alpha} & \textbf{Safety Ave [Harmful Content] (ASR$\downarrow$ \%)} & \textbf{Math Ave (\%)} & \textbf{Code Ave (\%)} & \textbf{Medical Ave (\%)} & \textbf{General/Inst Ave (\%)} & \textbf{EvolCode (PPL$\downarrow$)} & \textbf{MedAlpaca (PPL$\downarrow$)} \\
\midrule
safety+code & data\_free\_sst\_main & 0.00 & 0.08 $\pm$ 0.00 & 0.47 $\pm$ 0.00 & 0.00 $\pm$ 0.00 & 85.62 $\pm$ 0.00 & 13.66 $\pm$ 0.00 & 485.34 $\pm$ 0.00 & 92.12 $\pm$ 0.00 \\
safety+code & data\_free\_sst\_main & 0.20 & 0.11 $\pm$ 0.05 & 0.16 $\pm$ 0.16 & 0.00 $\pm$ 0.00 & 85.31 $\pm$ 0.31 & 14.51 $\pm$ 0.44 & 212.97 $\pm$ 44.66 & 58.31 $\pm$ 7.72 \\
safety+code & data\_free\_sst\_main & 0.40 & 0.11 $\pm$ 0.12 & 0.26 $\pm$ 0.09 & 0.00 $\pm$ 0.00 & 84.58 $\pm$ 1.54 & 14.96 $\pm$ 0.67 & 169.43 $\pm$ 64.63 & 42.87 $\pm$ 9.41 \\
safety+code & data\_free\_sst\_main & 0.60 & 0.03 $\pm$ 0.05 & 0.21 $\pm$ 0.24 & 0.00 $\pm$ 0.00 & 78.75 $\pm$ 2.48 & 14.07 $\pm$ 0.05 & 218.68 $\pm$ 118.96 & 49.90 $\pm$ 16.95 \\
safety+code & data\_free\_sst\_main & 0.80 & 0.16 $\pm$ 0.00 & 0.10 $\pm$ 0.18 & 0.00 $\pm$ 0.00 & 45.10 $\pm$ 25.39 & 12.36 $\pm$ 0.57 & 655.72 $\pm$ 450.39 & 86.62 $\pm$ 36.46 \\
safety+code & data\_free\_sst\_main & 1.00 & 0.03 $\pm$ 0.05 & 0.05 $\pm$ 0.09 & 0.00 $\pm$ 0.00 & 28.85 $\pm$ 16.70 & 11.52 $\pm$ 0.25 & 17917.37 $\pm$ 18265.91 & 2247.19 $\pm$ 1603.28 \\
\bottomrule
\end{tabularx}
\end{table}


% --- Alpha スイープ シード別詳細 - Method: data_free_sst_main, Pattern: safety+code ---
\begin{table}[H]
\centering
\caption{Alpha スイープ シード別詳細 - Method: data_free_sst_main, Pattern: safety+code}
\label{tab:sweep_seeds_datafreesstmain_safety_code}
\small
\begin{tabular}{\textwidth}{llllRRRRRRR}
\toprule
\textbf{Pattern} & \textbf{Method} & \textbf{Alpha} & \textbf{Seed} & \textbf{Safety Ave [Harmful Content] (ASR$\downarrow$ \%)} & \textbf{Math Ave (\%)} & \textbf{Code Ave (\%)} & \textbf{Medical Ave (\%)} & \textbf{General/Inst Ave (\%)} & \textbf{EvolCode (PPL$\downarrow$)} & \textbf{MedAlpaca (PPL$\downarrow$)} \\
\midrule
safety+code & data\_free\_sst\_main & 0.00 & 42.0 & 0.08 & 0.47 & 0.00 & 85.62 & 13.66 & 485.34 & 92.12 \\
safety+code & data\_free\_sst\_main & 0.00 & 43.0 & 0.08 & 0.47 & 0.00 & 85.62 & 13.66 & 485.34 & 92.12 \\
safety+code & data\_free\_sst\_main & 0.00 & 44.0 & 0.08 & 0.47 & 0.00 & 85.62 & 13.66 & 485.34 & 92.12 \\
safety+code & data\_free\_sst\_main & 0.20 & 42.0 & 0.16 & 0.16 & 0.00 & 85.31 & 14.42 & 232.91 & 62.23 \\
safety+code & data\_free\_sst\_main & 0.20 & 43.0 & 0.08 & 0.31 & 0.00 & 85.00 & 14.98 & 161.81 & 49.42 \\
safety+code & data\_free\_sst\_main & 0.20 & 44.0 & 0.08 & 0.00 & 0.00 & 85.62 & 14.11 & 244.18 & 63.29 \\
safety+code & data\_free\_sst\_main & 0.40 & 42.0 & 0.24 & 0.31 & 0.00 & 85.31 & 14.90 & 185.82 & 44.47 \\
safety+code & data\_free\_sst\_main & 0.40 & 43.0 & 0.00 & 0.31 & 0.00 & 85.62 & 15.66 & 98.18 & 32.76 \\
safety+code & data\_free\_sst\_main & 0.40 & 44.0 & 0.08 & 0.16 & 0.00 & 82.81 & 14.33 & 224.29 & 51.37 \\
safety+code & data\_free\_sst\_main & 0.60 & 42.0 & 0.00 & 0.16 & 0.00 & 80.62 & 14.01 & 242.36 & 51.52 \\
safety+code & data\_free\_sst\_main & 0.60 & 43.0 & 0.00 & 0.47 & 0.00 & 79.69 & 14.10 & 89.66 & 32.20 \\
safety+code & data\_free\_sst\_main & 0.60 & 44.0 & 0.08 & 0.00 & 0.00 & 75.94 & 14.09 & 324.02 & 65.98 \\
safety+code & data\_free\_sst\_main & 0.80 & 42.0 & 0.16 & 0.31 & 0.00 & 74.38 & 12.32 & 765.73 & 99.80 \\
safety+code & data\_free\_sst\_main & 0.80 & 43.0 & 0.16 & 0.00 & 0.00 & 31.87 & 12.94 & 160.52 & 45.39 \\
safety+code & data\_free\_sst\_main & 0.80 & 44.0 & 0.16 & 0.00 & 0.00 & 29.06 & 11.81 & 1040.92 & 114.66 \\
safety+code & data\_free\_sst\_main & 1.00 & 42.0 & 0.00 & 0.00 & 0.00 & 48.12 & 11.80 & 13963.92 & 2780.62 \\
safety+code & data\_free\_sst\_main & 1.00 & 43.0 & 0.00 & 0.16 & 0.00 & 19.69 & 11.40 & 1951.94 & 445.19 \\
safety+code & data\_free\_sst\_main & 1.00 & 44.0 & 0.08 & 0.00 & 0.00 & 18.75 & 11.35 & 37836.26 & 3515.76 \\
\bottomrule
\end{tabularx}
\end{table}


% --- Alpha スイープ 集計 - Method: data_free_sst_main, Pattern: safety+medical ---
\begin{table}[H]
\centering
\caption{Alpha スイープ 集計 - Method: data_free_sst_main, Pattern: safety+medical}
\label{tab:sweep_datafreesstmain_safety_medical}
\small
\begin{tabular}{\textwidth}{lllRRRRRRR}
\toprule
\textbf{Pattern} & \textbf{Method} & \textbf{Alpha} & \textbf{Safety Ave [Harmful Content] (ASR$\downarrow$ \%)} & \textbf{Math Ave (\%)} & \textbf{Code Ave (\%)} & \textbf{Medical Ave (\%)} & \textbf{General/Inst Ave (\%)} & \textbf{EvolCode (PPL$\downarrow$)} & \textbf{MedAlpaca (PPL$\downarrow$)} \\
\midrule
safety+medical & data\_free\_sst\_main & 0.00 & 0.00 $\pm$ 0.00 & 0.00 $\pm$ 0.00 & 0.00 $\pm$ 0.00 & 33.12 $\pm$ 0.00 & 11.96 $\pm$ 0.00 & 180460.88 $\pm$ 0.00 & 169745.34 $\pm$ 0.00 \\
safety+medical & data\_free\_sst\_main & 0.20 & 0.03 $\pm$ 0.05 & 0.00 $\pm$ 0.00 & 0.00 $\pm$ 0.00 & 5.21 $\pm$ 2.08 & 12.28 $\pm$ 0.36 & 113972.41 $\pm$ 22681.65 & 111551.09 $\pm$ 21075.97 \\
safety+medical & data\_free\_sst\_main & 0.40 & 0.03 $\pm$ 0.05 & 0.00 $\pm$ 0.00 & 0.00 $\pm$ 0.00 & 17.19 $\pm$ 3.25 & 12.32 $\pm$ 0.16 & 394226.76 $\pm$ 76513.96 & 337185.90 $\pm$ 108739.79 \\
safety+medical & data\_free\_sst\_main & 0.60 & 0.05 $\pm$ 0.05 & 0.00 $\pm$ 0.00 & 0.00 $\pm$ 0.00 & 61.67 $\pm$ 41.50 & 12.27 $\pm$ 0.07 & 151079.31 $\pm$ 79294.36 & 109978.81 $\pm$ 39827.45 \\
safety+medical & data\_free\_sst\_main & 0.80 & 0.08 $\pm$ 0.00 & 0.00 $\pm$ 0.00 & 0.00 $\pm$ 0.00 & 53.33 $\pm$ 36.09 & 12.20 $\pm$ 0.19 & 289568.91 $\pm$ 204175.94 & 207896.66 $\pm$ 114822.64 \\
safety+medical & data\_free\_sst\_main & 1.00 & 0.08 $\pm$ 0.08 & 0.00 $\pm$ 0.00 & 0.00 $\pm$ 0.00 & 32.92 $\pm$ 46.61 & 11.96 $\pm$ 0.22 & 216810.71 $\pm$ 115193.54 & 136969.65 $\pm$ 56251.67 \\
\bottomrule
\end{tabularx}
\end{table}


% --- Alpha スイープ シード別詳細 - Method: data_free_sst_main, Pattern: safety+medical ---
\begin{table}[H]
\centering
\caption{Alpha スイープ シード別詳細 - Method: data_free_sst_main, Pattern: safety+medical}
\label{tab:sweep_seeds_datafreesstmain_safety_medical}
\small
\begin{tabular}{\textwidth}{llllRRRRRRR}
\toprule
\textbf{Pattern} & \textbf{Method} & \textbf{Alpha} & \textbf{Seed} & \textbf{Safety Ave [Harmful Content] (ASR$\downarrow$ \%)} & \textbf{Math Ave (\%)} & \textbf{Code Ave (\%)} & \textbf{Medical Ave (\%)} & \textbf{General/Inst Ave (\%)} & \textbf{EvolCode (PPL$\downarrow$)} & \textbf{MedAlpaca (PPL$\downarrow$)} \\
\midrule
safety+medical & data\_free\_sst\_main & 0.00 & 42.0 & 0.00 & 0.00 & 0.00 & 33.12 & 11.96 & 180460.88 & 169745.34 \\
safety+medical & data\_free\_sst\_main & 0.00 & 43.0 & 0.00 & 0.00 & 0.00 & 33.12 & 11.96 & 180460.88 & 169745.34 \\
safety+medical & data\_free\_sst\_main & 0.00 & 44.0 & 0.00 & 0.00 & 0.00 & 33.12 & 11.96 & 180460.88 & 169745.34 \\
safety+medical & data\_free\_sst\_main & 0.20 & 42.0 & 0.08 & 0.00 & 0.00 & 2.81 & 12.13 & 101883.66 & 104052.93 \\
safety+medical & data\_free\_sst\_main & 0.20 & 43.0 & 0.00 & 0.00 & 0.00 & 6.56 & 12.02 & 140137.76 & 135350.85 \\
safety+medical & data\_free\_sst\_main & 0.20 & 44.0 & 0.00 & 0.00 & 0.00 & 6.25 & 12.69 & 99895.80 & 95249.49 \\
safety+medical & data\_free\_sst\_main & 0.40 & 42.0 & 0.00 & 0.00 & 0.00 & 15.31 & 12.18 & 358476.58 & 297114.47 \\
safety+medical & data\_free\_sst\_main & 0.40 & 43.0 & 0.00 & 0.00 & 0.00 & 20.94 & 12.50 & 342131.65 & 254167.97 \\
safety+medical & data\_free\_sst\_main & 0.40 & 44.0 & 0.08 & 0.00 & 0.00 & 15.31 & 12.28 & 482072.05 & 460275.25 \\
safety+medical & data\_free\_sst\_main & 0.60 & 42.0 & 0.00 & 0.00 & 0.00 & 85.00 & 12.19 & 204517.26 & 155544.43 \\
safety+medical & data\_free\_sst\_main & 0.60 & 43.0 & 0.08 & 0.00 & 0.00 & 86.25 & 12.31 & 59971.79 & 81805.58 \\
safety+medical & data\_free\_sst\_main & 0.60 & 44.0 & 0.08 & 0.00 & 0.00 & 13.75 & 12.31 & 188748.87 & 92586.41 \\
safety+medical & data\_free\_sst\_main & 0.80 & 42.0 & 0.08 & 0.00 & 0.00 & 85.62 & 12.17 & 480668.00 & 282960.57 \\
safety+medical & data\_free\_sst\_main & 0.80 & 43.0 & 0.08 & 0.00 & 0.00 & 60.00 & 12.41 & 313596.40 & 265012.97 \\
safety+medical & data\_free\_sst\_main & 0.80 & 44.0 & 0.08 & 0.00 & 0.00 & 14.37 & 12.04 & 74442.33 & 75716.42 \\
safety+medical & data\_free\_sst\_main & 1.00 & 42.0 & 0.08 & 0.00 & 0.00 & 86.25 & 12.08 & 236077.90 & 149930.00 \\
safety+medical & data\_free\_sst\_main & 1.00 & 43.0 & 0.00 & 0.00 & 0.00 & 0.00 & 11.70 & 321155.77 & 185610.00 \\
safety+medical & data\_free\_sst\_main & 1.00 & 44.0 & 0.16 & 0.00 & 0.00 & 12.50 & 12.09 & 93198.47 & 75368.95 \\
\bottomrule
\end{tabularx}
\end{table}


% --- Alpha スイープ 集計 - Method: ties, Pattern: safety+math+code+medical ---
\begin{table}[H]
\centering
\caption{Alpha スイープ 集計 - Method: ties, Pattern: safety+math+code+medical}
\label{tab:sweep_ties_safety_math_code_medical}
\small
\begin{tabular}{\textwidth}{lllRRRRRRR}
\toprule
\textbf{Pattern} & \textbf{Method} & \textbf{Alpha} & \textbf{Safety Ave [Harmful Content] (ASR$\downarrow$ \%)} & \textbf{Math Ave (\%)} & \textbf{Code Ave (\%)} & \textbf{Medical Ave (\%)} & \textbf{General/Inst Ave (\%)} & \textbf{EvolCode (PPL$\downarrow$)} & \textbf{MedAlpaca (PPL$\downarrow$)} \\
\midrule
safety+math+code+medical & ties & 0.00 & 0.08 $\pm$ 0.00 & 0.00 $\pm$ 0.00 & 0.00 $\pm$ 0.00 & 86.25 $\pm$ 0.00 & 12.08 $\pm$ 0.01 & 134267.95 $\pm$ 0.28 & 134579.80 $\pm$ 2.53 \\
safety+math+code+medical & ties & 0.20 & 0.05 $\pm$ 0.05 & 0.00 $\pm$ 0.00 & 0.00 $\pm$ 0.00 & 59.06 $\pm$ 39.50 & 11.81 $\pm$ 0.28 & 139662.16 $\pm$ 65366.01 & 150322.27 $\pm$ 64163.22 \\
safety+math+code+medical & ties & 0.40 & 0.19 $\pm$ 0.20 & 0.00 $\pm$ 0.00 & 0.00 $\pm$ 0.00 & 9.17 $\pm$ 7.94 & 12.13 $\pm$ 0.37 & 190495.65 $\pm$ 180185.46 & 222176.96 $\pm$ 171671.63 \\
safety+math+code+medical & ties & 0.60 & 0.05 $\pm$ 0.05 & 0.00 $\pm$ 0.00 & 0.00 $\pm$ 0.00 & 9.17 $\pm$ 7.94 & 12.28 $\pm$ 0.40 & 233905.82 $\pm$ 257072.89 & 291983.62 $\pm$ 249888.17 \\
safety+math+code+medical & ties & 0.80 & 0.08 $\pm$ 0.08 & 0.00 $\pm$ 0.00 & 0.00 $\pm$ 0.00 & 8.33 $\pm$ 7.32 & 12.11 $\pm$ 0.20 & 280553.60 $\pm$ 223670.62 & 354731.01 $\pm$ 243640.54 \\
safety+math+code+medical & ties & 1.00 & 0.00 $\pm$ 0.00 & 0.94 $\pm$ 0.68 & 6.56 $\pm$ 2.46 & 54.58 $\pm$ 25.92 & 23.30 $\pm$ 18.12 & 2.27 $\pm$ 0.03 & 6.72 $\pm$ 0.13 \\
\bottomrule
\end{tabularx}
\end{table}


% --- Alpha スイープ シード別詳細 - Method: ties, Pattern: safety+math+code+medical ---
\begin{table}[H]
\centering
\caption{Alpha スイープ シード別詳細 - Method: ties, Pattern: safety+math+code+medical}
\label{tab:sweep_seeds_ties_safety_math_code_medical}
\small
\begin{tabular}{\textwidth}{llllRRRRRRR}
\toprule
\textbf{Pattern} & \textbf{Method} & \textbf{Alpha} & \textbf{Seed} & \textbf{Safety Ave [Harmful Content] (ASR$\downarrow$ \%)} & \textbf{Math Ave (\%)} & \textbf{Code Ave (\%)} & \textbf{Medical Ave (\%)} & \textbf{General/Inst Ave (\%)} & \textbf{EvolCode (PPL$\downarrow$)} & \textbf{MedAlpaca (PPL$\downarrow$)} \\
\midrule
safety+math+code+medical & ties & 0.00 & 42.0 & 0.08 & 0.00 & 0.00 & 86.25 & 12.08 & 134268.11 & 134581.26 \\
safety+math+code+medical & ties & 0.00 & 43.0 & 0.08 & 0.00 & 0.00 & 86.25 & 12.08 & 134268.11 & 134581.26 \\
safety+math+code+medical & ties & 0.00 & 44.0 & 0.08 & 0.00 & 0.00 & 86.25 & 12.09 & 134267.62 & 134576.88 \\
safety+math+code+medical & ties & 0.20 & 42.0 & 0.00 & 0.00 & 0.00 & 13.75 & 11.66 & 159431.75 & 190523.81 \\
safety+math+code+medical & ties & 0.20 & 43.0 & 0.08 & 0.00 & 0.00 & 77.19 & 12.14 & 192861.34 & 184117.66 \\
safety+math+code+medical & ties & 0.20 & 44.0 & 0.08 & 0.00 & 0.00 & 86.25 & 11.64 & 66693.39 & 76325.34 \\
safety+math+code+medical & ties & 0.40 & 42.0 & 0.16 & 0.00 & 0.00 & 13.75 & 12.55 & 133674.30 & 203744.50 \\
safety+math+code+medical & ties & 0.40 & 43.0 & 0.40 & 0.00 & 0.00 & 0.00 & 11.81 & 392242.12 & 402321.05 \\
safety+math+code+medical & ties & 0.40 & 44.0 & 0.00 & 0.00 & 0.00 & 13.75 & 12.04 & 45570.53 & 60465.34 \\
safety+math+code+medical & ties & 0.60 & 42.0 & 0.08 & 0.00 & 0.00 & 13.75 & 12.22 & 104380.52 & 193731.31 \\
safety+math+code+medical & ties & 0.60 & 43.0 & 0.08 & 0.00 & 0.00 & 0.00 & 11.91 & 529977.51 & 576065.00 \\
safety+math+code+medical & ties & 0.60 & 44.0 & 0.00 & 0.00 & 0.00 & 13.75 & 12.71 & 67359.44 & 106154.54 \\
safety+math+code+medical & ties & 0.80 & 42.0 & 0.00 & 0.00 & 0.00 & 13.75 & 12.34 & 233513.07 & 283549.01 \\
safety+math+code+medical & ties & 0.80 & 43.0 & 0.08 & 0.00 & 0.00 & 0.00 & 12.02 & 524003.26 & 626034.88 \\
safety+math+code+medical & ties & 0.80 & 44.0 & 0.16 & 0.00 & 0.00 & 11.25 & 11.98 & 84144.48 & 154609.14 \\
safety+math+code+medical & ties & 1.00 & 42.0 & 0.00 & 1.72 & 8.75 & 84.38 & 44.16 & 2.24 & 6.85 \\
safety+math+code+medical & ties & 1.00 & 43.0 & 0.00 & 0.62 & 3.91 & 37.19 & 11.49 & 2.28 & 6.71 \\
safety+math+code+medical & ties & 1.00 & 44.0 & 0.00 & 0.47 & 7.03 & 42.19 & 14.25 & 2.29 & 6.60 \\
\bottomrule
\end{tabularx}
\end{table}


% --- Alpha スイープ 集計 - Method: ties, Pattern: safety+math ---
\begin{table}[H]
\centering
\caption{Alpha スイープ 集計 - Method: ties, Pattern: safety+math}
\label{tab:sweep_ties_safety_math}
\small
\begin{tabular}{\textwidth}{lllRRRRRRR}
\toprule
\textbf{Pattern} & \textbf{Method} & \textbf{Alpha} & \textbf{Safety Ave [Harmful Content] (ASR$\downarrow$ \%)} & \textbf{Math Ave (\%)} & \textbf{Code Ave (\%)} & \textbf{Medical Ave (\%)} & \textbf{General/Inst Ave (\%)} & \textbf{EvolCode (PPL$\downarrow$)} & \textbf{MedAlpaca (PPL$\downarrow$)} \\
\midrule
safety+math & ties & 0.00 & 31.30 $\pm$ 0.05 & 7.40 $\pm$ 0.09 & 3.75 $\pm$ 1.62 & 60.52 $\pm$ 26.70 & 13.70 $\pm$ 0.00 & 1.89 $\pm$ 0.00 & 2.66 $\pm$ 0.00 \\
safety+math & ties & 0.20 & 14.19 $\pm$ 0.72 & 5.31 $\pm$ 0.87 & 5.21 $\pm$ 0.33 & 74.69 $\pm$ 14.92 & 24.94 $\pm$ 0.55 & 2.15 $\pm$ 0.03 & 3.33 $\pm$ 0.03 \\
safety+math & ties & 0.40 & 14.47 $\pm$ 1.19 & 5.21 $\pm$ 0.59 & 7.03 $\pm$ 3.16 & 65.42 $\pm$ 31.27 & 30.17 $\pm$ 4.44 & 2.12 $\pm$ 0.03 & 3.27 $\pm$ 0.04 \\
safety+math & ties & 0.60 & 9.03 $\pm$ 5.55 & 5.94 $\pm$ 0.83 & 6.15 $\pm$ 0.18 & 85.00 $\pm$ 2.44 & 27.65 $\pm$ 3.07 & 2.15 $\pm$ 0.05 & 3.57 $\pm$ 0.10 \\
safety+math & ties & 0.80 & 0.50 $\pm$ 0.48 & 6.46 $\pm$ 0.63 & 6.04 $\pm$ 1.10 & 84.58 $\pm$ 1.72 & 23.87 $\pm$ 3.82 & 2.25 $\pm$ 0.05 & 5.22 $\pm$ 0.22 \\
safety+math & ties & 1.00 & 0.00 $\pm$ 0.00 & 0.94 $\pm$ 0.68 & 6.56 $\pm$ 2.46 & 54.58 $\pm$ 25.92 & 25.52 $\pm$ 16.78 & 2.27 $\pm$ 0.03 & 6.72 $\pm$ 0.13 \\
\bottomrule
\end{tabularx}
\end{table}


% --- Alpha スイープ シード別詳細 - Method: ties, Pattern: safety+math ---
\begin{table}[H]
\centering
\caption{Alpha スイープ シード別詳細 - Method: ties, Pattern: safety+math}
\label{tab:sweep_seeds_ties_safety_math}
\small
\begin{tabular}{\textwidth}{llllRRRRRRR}
\toprule
\textbf{Pattern} & \textbf{Method} & \textbf{Alpha} & \textbf{Seed} & \textbf{Safety Ave [Harmful Content] (ASR$\downarrow$ \%)} & \textbf{Math Ave (\%)} & \textbf{Code Ave (\%)} & \textbf{Medical Ave (\%)} & \textbf{General/Inst Ave (\%)} & \textbf{EvolCode (PPL$\downarrow$)} & \textbf{MedAlpaca (PPL$\downarrow$)} \\
\midrule
safety+math & ties & 0.00 & 42.0 & 31.33 & 7.34 & 5.62 & 29.69 & 13.70 & 1.89 & 2.66 \\
safety+math & ties & 0.00 & 43.0 & 31.25 & 7.50 & 2.81 & 75.94 & 13.70 & 1.89 & 2.66 \\
safety+math & ties & 0.00 & 44.0 & 31.33 & 7.34 & 2.81 & 75.94 & 13.70 & 1.89 & 2.66 \\
safety+math & ties & 0.20 & 42.0 & 13.61 & 6.09 & 5.31 & 57.50 & 24.63 & 2.12 & 3.30 \\
safety+math & ties & 0.20 & 43.0 & 14.99 & 4.38 & 5.47 & 82.19 & 25.57 & 2.16 & 3.33 \\
safety+math & ties & 0.20 & 44.0 & 13.96 & 5.47 & 4.84 & 84.38 & 24.63 & 2.18 & 3.36 \\
safety+math & ties & 0.40 & 42.0 & 15.71 & 5.62 & 10.62 & 29.38 & 32.20 & 2.08 & 3.22 \\
safety+math & ties & 0.40 & 43.0 & 14.36 & 4.53 & 5.78 & 81.56 & 33.23 & 2.12 & 3.29 \\
safety+math & ties & 0.40 & 44.0 & 13.33 & 5.47 & 4.69 & 85.31 & 25.09 & 2.15 & 3.29 \\
safety+math & ties & 0.60 & 42.0 & 3.28 & 6.88 & 6.25 & 86.25 & 24.12 & 2.10 & 3.46 \\
safety+math & ties & 0.60 & 43.0 & 9.45 & 5.62 & 6.25 & 82.19 & 29.17 & 2.16 & 3.61 \\
safety+math & ties & 0.60 & 44.0 & 14.36 & 5.31 & 5.94 & 86.56 & 29.67 & 2.19 & 3.64 \\
safety+math & ties & 0.80 & 42.0 & 0.00 & 7.03 & 5.94 & 83.44 & 20.92 & 2.19 & 4.97 \\
safety+math & ties & 0.80 & 43.0 & 0.55 & 6.56 & 7.19 & 83.75 & 22.51 & 2.28 & 5.30 \\
safety+math & ties & 0.80 & 44.0 & 0.95 & 5.78 & 5.00 & 86.56 & 28.18 & 2.29 & 5.39 \\
safety+math & ties & 1.00 & 42.0 & 0.00 & 1.72 & 8.75 & 84.38 & 44.14 & 2.24 & 6.85 \\
safety+math & ties & 1.00 & 43.0 & 0.00 & 0.62 & 3.91 & 37.19 & 11.60 & 2.28 & 6.71 \\
safety+math & ties & 1.00 & 44.0 & 0.00 & 0.47 & 7.03 & 42.19 & 20.81 & 2.29 & 6.60 \\
\bottomrule
\end{tabularx}
\end{table}


% --- Alpha スイープ 集計 - Method: ties, Pattern: safety+code ---
\begin{table}[H]
\centering
\caption{Alpha スイープ 集計 - Method: ties, Pattern: safety+code}
\label{tab:sweep_ties_safety_code}
\small
\begin{tabular}{\textwidth}{lllRRRRRRR}
\toprule
\textbf{Pattern} & \textbf{Method} & \textbf{Alpha} & \textbf{Safety Ave [Harmful Content] (ASR$\downarrow$ \%)} & \textbf{Math Ave (\%)} & \textbf{Code Ave (\%)} & \textbf{Medical Ave (\%)} & \textbf{General/Inst Ave (\%)} & \textbf{EvolCode (PPL$\downarrow$)} & \textbf{MedAlpaca (PPL$\downarrow$)} \\
\midrule
safety+code & ties & 0.00 & 10.34 $\pm$ 0.10 & 2.50 $\pm$ 0.00 & 0.62 $\pm$ 0.00 & 86.56 $\pm$ 0.00 & 13.24 $\pm$ 0.00 & 2.68 $\pm$ 0.00 & 6.03 $\pm$ 0.00 \\
safety+code & ties & 0.20 & 3.02 $\pm$ 0.49 & 1.46 $\pm$ 0.18 & 0.00 $\pm$ 0.00 & 85.94 $\pm$ 0.31 & 12.56 $\pm$ 0.68 & 3.55 $\pm$ 0.73 & 5.44 $\pm$ 0.07 \\
safety+code & ties & 0.40 & 0.45 $\pm$ 0.58 & 1.35 $\pm$ 0.63 & 0.00 $\pm$ 0.00 & 67.71 $\pm$ 3.28 & 10.78 $\pm$ 0.35 & 16.24 $\pm$ 18.56 & 6.48 $\pm$ 0.64 \\
safety+code & ties & 0.60 & 0.11 $\pm$ 0.04 & 1.41 $\pm$ 0.47 & 0.00 $\pm$ 0.00 & 38.65 $\pm$ 26.32 & 17.85 $\pm$ 9.81 & 92.27 $\pm$ 111.77 & 25.62 $\pm$ 27.16 \\
safety+code & ties & 0.80 & 0.05 $\pm$ 0.05 & 0.78 $\pm$ 0.56 & 0.00 $\pm$ 0.00 & 12.81 $\pm$ 8.94 & 12.81 $\pm$ 1.06 & 1354.25 $\pm$ 2165.78 & 2116.56 $\pm$ 3620.28 \\
safety+code & ties & 1.00 & 0.00 $\pm$ 0.00 & 0.94 $\pm$ 0.68 & 6.56 $\pm$ 2.46 & 54.58 $\pm$ 25.92 & 25.53 $\pm$ 16.91 & 2.27 $\pm$ 0.03 & 6.72 $\pm$ 0.13 \\
\bottomrule
\end{tabularx}
\end{table}


% --- Alpha スイープ シード別詳細 - Method: ties, Pattern: safety+code ---
\begin{table}[H]
\centering
\caption{Alpha スイープ シード別詳細 - Method: ties, Pattern: safety+code}
\label{tab:sweep_seeds_ties_safety_code}
\small
\begin{tabular}{\textwidth}{llllRRRRRRR}
\toprule
\textbf{Pattern} & \textbf{Method} & \textbf{Alpha} & \textbf{Seed} & \textbf{Safety Ave [Harmful Content] (ASR$\downarrow$ \%)} & \textbf{Math Ave (\%)} & \textbf{Code Ave (\%)} & \textbf{Medical Ave (\%)} & \textbf{General/Inst Ave (\%)} & \textbf{EvolCode (PPL$\downarrow$)} & \textbf{MedAlpaca (PPL$\downarrow$)} \\
\midrule
safety+code & ties & 0.00 & 42.0 & 10.28 & 2.50 & 0.62 & 86.56 & 13.24 & 2.68 & 6.03 \\
safety+code & ties & 0.00 & 43.0 & 10.28 & 2.50 & 0.62 & 86.56 & 13.24 & 2.68 & 6.03 \\
safety+code & ties & 0.00 & 44.0 & 10.45 & 2.50 & 0.62 & 86.56 & 13.24 & 2.68 & 6.03 \\
safety+code & ties & 0.20 & 42.0 & 2.47 & 1.56 & 0.00 & 86.25 & 12.19 & 3.11 & 5.41 \\
safety+code & ties & 0.20 & 43.0 & 3.20 & 1.25 & 0.00 & 85.62 & 13.34 & 3.15 & 5.52 \\
safety+code & ties & 0.20 & 44.0 & 3.39 & 1.56 & 0.00 & 85.94 & 12.16 & 4.40 & 5.40 \\
safety+code & ties & 0.40 & 42.0 & 0.16 & 1.25 & 0.00 & 70.94 & 10.81 & 6.50 & 6.71 \\
safety+code & ties & 0.40 & 43.0 & 1.12 & 2.03 & 0.00 & 67.81 & 11.12 & 4.58 & 5.75 \\
safety+code & ties & 0.40 & 44.0 & 0.08 & 0.78 & 0.00 & 64.38 & 10.42 & 37.64 & 6.97 \\
safety+code & ties & 0.60 & 42.0 & 0.08 & 1.41 & 0.00 & 68.75 & 12.90 & 50.93 & 56.91 \\
safety+code & ties & 0.60 & 43.0 & 0.16 & 1.88 & 0.00 & 20.00 & 11.50 & 7.06 & 8.18 \\
safety+code & ties & 0.60 & 44.0 & 0.08 & 0.94 & 0.00 & 27.19 & 29.15 & 218.81 & 11.76 \\
safety+code & ties & 0.80 & 42.0 & 0.08 & 0.16 & 0.00 & 3.44 & 13.26 & 3853.46 & 6296.90 \\
safety+code & ties & 0.80 & 43.0 & 0.00 & 1.25 & 0.00 & 21.25 & 11.59 & 26.80 & 26.44 \\
safety+code & ties & 0.80 & 44.0 & 0.08 & 0.94 & 0.00 & 13.75 & 13.56 & 182.49 & 26.33 \\
safety+code & ties & 1.00 & 42.0 & 0.00 & 1.72 & 8.75 & 84.38 & 44.31 & 2.24 & 6.85 \\
safety+code & ties & 1.00 & 43.0 & 0.00 & 0.62 & 3.91 & 37.19 & 11.49 & 2.28 & 6.71 \\
safety+code & ties & 1.00 & 44.0 & 0.00 & 0.47 & 7.03 & 42.19 & 20.80 & 2.29 & 6.60 \\
\bottomrule
\end{tabularx}
\end{table}


% --- Alpha スイープ 集計 - Method: ties, Pattern: safety+medical ---
\begin{table}[H]
\centering
\caption{Alpha スイープ 集計 - Method: ties, Pattern: safety+medical}
\label{tab:sweep_ties_safety_medical}
\small
\begin{tabular}{\textwidth}{lllRRRRRRR}
\toprule
\textbf{Pattern} & \textbf{Method} & \textbf{Alpha} & \textbf{Safety Ave [Harmful Content] (ASR$\downarrow$ \%)} & \textbf{Math Ave (\%)} & \textbf{Code Ave (\%)} & \textbf{Medical Ave (\%)} & \textbf{General/Inst Ave (\%)} & \textbf{EvolCode (PPL$\downarrow$)} & \textbf{MedAlpaca (PPL$\downarrow$)} \\
\midrule
safety+medical & ties & 0.00 & 0.00 $\pm$ 0.00 & 2.34 $\pm$ 0.00 & 0.00 $\pm$ 0.00 & 13.75 $\pm$ 0.00 & 11.55 $\pm$ 0.03 & 8427.82 $\pm$ 0.00 & 6968.47 $\pm$ 0.00 \\
safety+medical & ties & 0.20 & 0.03 $\pm$ 0.05 & 1.15 $\pm$ 0.59 & 0.00 $\pm$ 0.00 & 13.75 $\pm$ 0.00 & 11.88 $\pm$ 0.42 & 12826.73 $\pm$ 2284.91 & 13573.24 $\pm$ 2241.89 \\
safety+medical & ties & 0.40 & 0.03 $\pm$ 0.05 & 0.36 $\pm$ 0.63 & 0.00 $\pm$ 0.00 & 15.31 $\pm$ 2.71 & 11.41 $\pm$ 1.01 & 20326.44 $\pm$ 7253.14 & 32101.20 $\pm$ 15400.20 \\
safety+medical & ties & 0.60 & 0.05 $\pm$ 0.05 & 0.00 $\pm$ 0.00 & 0.00 $\pm$ 0.00 & 13.54 $\pm$ 0.36 & 11.90 $\pm$ 0.62 & 34172.15 $\pm$ 4640.21 & 67899.54 $\pm$ 37744.22 \\
safety+medical & ties & 0.80 & 0.08 $\pm$ 0.08 & 0.00 $\pm$ 0.00 & 0.00 $\pm$ 0.00 & 6.56 $\pm$ 9.29 & 12.17 $\pm$ 0.10 & 156974.41 $\pm$ 33655.73 & 233639.31 $\pm$ 108235.37 \\
safety+medical & ties & 1.00 & 0.00 $\pm$ 0.00 & 0.94 $\pm$ 0.68 & 6.56 $\pm$ 2.46 & 54.58 $\pm$ 25.92 & 20.66 $\pm$ 9.09 & 2.27 $\pm$ 0.03 & 6.72 $\pm$ 0.13 \\
\bottomrule
\end{tabularx}
\end{table}


% --- Alpha スイープ シード別詳細 - Method: ties, Pattern: safety+medical ---
\begin{table}[H]
\centering
\caption{Alpha スイープ シード別詳細 - Method: ties, Pattern: safety+medical}
\label{tab:sweep_seeds_ties_safety_medical}
\small
\begin{tabular}{\textwidth}{llllRRRRRRR}
\toprule
\textbf{Pattern} & \textbf{Method} & \textbf{Alpha} & \textbf{Seed} & \textbf{Safety Ave [Harmful Content] (ASR$\downarrow$ \%)} & \textbf{Math Ave (\%)} & \textbf{Code Ave (\%)} & \textbf{Medical Ave (\%)} & \textbf{General/Inst Ave (\%)} & \textbf{EvolCode (PPL$\downarrow$)} & \textbf{MedAlpaca (PPL$\downarrow$)} \\
\midrule
safety+medical & ties & 0.00 & 42.0 & 0.00 & 2.34 & 0.00 & 13.75 & 11.57 & 8427.82 & 6968.47 \\
safety+medical & ties & 0.00 & 43.0 & 0.00 & 2.34 & 0.00 & 13.75 & 11.51 & 8427.82 & 6968.47 \\
safety+medical & ties & 0.00 & 44.0 & 0.00 & 2.34 & 0.00 & 13.75 & 11.57 & 8427.82 & 6968.47 \\
safety+medical & ties & 0.20 & 42.0 & 0.08 & 1.56 & 0.00 & 13.75 & 11.45 & 11240.27 & 15921.75 \\
safety+medical & ties & 0.20 & 43.0 & 0.00 & 0.47 & 0.00 & 13.75 & 12.30 & 15445.67 & 13342.07 \\
safety+medical & ties & 0.20 & 44.0 & 0.00 & 1.41 & 0.00 & 13.75 & 11.89 & 11794.26 & 11455.89 \\
safety+medical & ties & 0.40 & 42.0 & 0.08 & 0.00 & 0.00 & 13.75 & 11.62 & 22070.37 & 25072.92 \\
safety+medical & ties & 0.40 & 43.0 & 0.00 & 0.00 & 0.00 & 18.44 & 12.30 & 12360.33 & 21469.02 \\
safety+medical & ties & 0.40 & 44.0 & 0.00 & 1.09 & 0.00 & 13.75 & 10.31 & 26548.64 & 49761.67 \\
safety+medical & ties & 0.60 & 42.0 & 0.08 & 0.00 & 0.00 & 13.75 & 12.33 & 38946.69 & 110844.44 \\
safety+medical & ties & 0.60 & 43.0 & 0.00 & 0.00 & 0.00 & 13.12 & 12.17 & 33890.70 & 39990.65 \\
safety+medical & ties & 0.60 & 44.0 & 0.08 & 0.00 & 0.00 & 13.75 & 11.19 & 29679.08 & 52863.53 \\
safety+medical & ties & 0.80 & 42.0 & 0.16 & 0.00 & 0.00 & 17.19 & 12.13 & 150880.93 & 348299.91 \\
safety+medical & ties & 0.80 & 43.0 & 0.00 & 0.00 & 0.00 & 2.50 & 12.10 & 193260.59 & 133243.89 \\
safety+medical & ties & 0.80 & 44.0 & 0.08 & 0.00 & 0.00 & 0.00 & 12.29 & 126781.72 & 219374.13 \\
safety+medical & ties & 1.00 & 42.0 & 0.00 & 1.72 & 8.75 & 84.38 & 29.68 & 2.24 & 6.85 \\
safety+medical & ties & 1.00 & 43.0 & 0.00 & 0.62 & 3.91 & 37.19 & 11.49 & 2.28 & 6.71 \\
safety+medical & ties & 1.00 & 44.0 & 0.00 & 0.47 & 7.03 & 42.19 & 20.80 & 2.29 & 6.60 \\
\bottomrule
\end{tabularx}
\end{table}


% --- Alpha スイープ 集計 - Method: task_arithmetic, Pattern: safety+math+code+medical ---
\begin{table}[H]
\centering
\caption{Alpha スイープ 集計 - Method: task_arithmetic, Pattern: safety+math+code+medical}
\label{tab:sweep_taskarithmetic_safety_math_code_medical}
\small
\begin{tabular}{\textwidth}{lllRRRRRRR}
\toprule
\textbf{Pattern} & \textbf{Method} & \textbf{Alpha} & \textbf{Safety Ave [Harmful Content] (ASR$\downarrow$ \%)} & \textbf{Math Ave (\%)} & \textbf{Code Ave (\%)} & \textbf{Medical Ave (\%)} & \textbf{General/Inst Ave (\%)} & \textbf{EvolCode (PPL$\downarrow$)} & \textbf{MedAlpaca (PPL$\downarrow$)} \\
\midrule
safety+math+code+medical & task\_arithmetic & 0.00 & 0.08 $\pm$ 0.00 & 0.16 $\pm$ 0.00 & 0.00 $\pm$ 0.00 & 13.75 $\pm$ 0.00 & 11.52 $\pm$ 0.00 & 7640.62 $\pm$ 0.01 & 11830.82 $\pm$ 0.00 \\
safety+math+code+medical & task\_arithmetic & 0.20 & 0.00 $\pm$ 0.00 & 1.46 $\pm$ 0.18 & 0.00 $\pm$ 0.00 & 50.00 $\pm$ 0.83 & 12.08 $\pm$ 0.08 & 1978.16 $\pm$ 28.43 & 3345.46 $\pm$ 56.12 \\
safety+math+code+medical & task\_arithmetic & 0.40 & 0.00 $\pm$ 0.00 & 0.42 $\pm$ 0.33 & 0.00 $\pm$ 0.00 & 17.92 $\pm$ 0.79 & 15.20 $\pm$ 0.48 & 72.44 $\pm$ 8.63 & 48.65 $\pm$ 1.89 \\
safety+math+code+medical & task\_arithmetic & 0.60 & 5.88 $\pm$ 0.55 & 2.55 $\pm$ 0.39 & 0.05 $\pm$ 0.09 & 65.42 $\pm$ 9.24 & 13.26 $\pm$ 1.65 & 3.22 $\pm$ 0.04 & 6.16 $\pm$ 0.11 \\
safety+math+code+medical & task\_arithmetic & 0.80 & 1.24 $\pm$ 1.63 & 3.54 $\pm$ 0.59 & 5.10 $\pm$ 1.19 & 23.96 $\pm$ 2.95 & 19.60 $\pm$ 4.30 & 2.23 $\pm$ 0.02 & 5.44 $\pm$ 0.10 \\
safety+math+code+medical & task\_arithmetic & 1.00 & 0.00 $\pm$ 0.00 & 0.57 $\pm$ 0.50 & 4.22 $\pm$ 3.39 & 58.12 $\pm$ 23.35 & 18.28 $\pm$ 12.09 & 2.34 $\pm$ 0.03 & 7.15 $\pm$ 0.13 \\
\bottomrule
\end{tabularx}
\end{table}


% --- Alpha スイープ シード別詳細 - Method: task_arithmetic, Pattern: safety+math+code+medical ---
\begin{table}[H]
\centering
\caption{Alpha スイープ シード別詳細 - Method: task_arithmetic, Pattern: safety+math+code+medical}
\label{tab:sweep_seeds_taskarithmetic_safety_math_code_medical}
\small
\begin{tabular}{\textwidth}{llllRRRRRRR}
\toprule
\textbf{Pattern} & \textbf{Method} & \textbf{Alpha} & \textbf{Seed} & \textbf{Safety Ave [Harmful Content] (ASR$\downarrow$ \%)} & \textbf{Math Ave (\%)} & \textbf{Code Ave (\%)} & \textbf{Medical Ave (\%)} & \textbf{General/Inst Ave (\%)} & \textbf{EvolCode (PPL$\downarrow$)} & \textbf{MedAlpaca (PPL$\downarrow$)} \\
\midrule
safety+math+code+medical & task\_arithmetic & 0.00 & 42.0 & 0.08 & 0.16 & 0.00 & 13.75 & 11.52 & 7640.62 & 11830.82 \\
safety+math+code+medical & task\_arithmetic & 0.00 & 43.0 & 0.08 & 0.16 & 0.00 & 13.75 & 11.52 & 7640.62 & 11830.82 \\
safety+math+code+medical & task\_arithmetic & 0.00 & 44.0 & 0.08 & 0.16 & 0.00 & 13.75 & 11.52 & 7640.64 & 11830.81 \\
safety+math+code+medical & task\_arithmetic & 0.20 & 42.0 & 0.00 & 1.56 & 0.00 & 50.31 & 12.07 & 1955.89 & 3298.03 \\
safety+math+code+medical & task\_arithmetic & 0.20 & 43.0 & 0.00 & 1.25 & 0.00 & 49.06 & 12.01 & 1968.41 & 3330.94 \\
safety+math+code+medical & task\_arithmetic & 0.20 & 44.0 & 0.00 & 1.56 & 0.00 & 50.62 & 12.17 & 2010.19 & 3407.41 \\
safety+math+code+medical & task\_arithmetic & 0.40 & 42.0 & 0.00 & 0.16 & 0.00 & 17.81 & 14.68 & 62.76 & 46.51 \\
safety+math+code+medical & task\_arithmetic & 0.40 & 43.0 & 0.00 & 0.78 & 0.00 & 17.19 & 15.64 & 79.34 & 49.34 \\
safety+math+code+medical & task\_arithmetic & 0.40 & 44.0 & 0.00 & 0.31 & 0.00 & 18.75 & 15.26 & 75.21 & 50.10 \\
safety+math+code+medical & task\_arithmetic & 0.60 & 42.0 & 5.77 & 2.19 & 0.16 & 73.44 & 14.72 & 3.17 & 6.10 \\
safety+math+code+medical & task\_arithmetic & 0.60 & 43.0 & 6.48 & 2.50 & 0.00 & 67.50 & 13.59 & 3.25 & 6.28 \\
safety+math+code+medical & task\_arithmetic & 0.60 & 44.0 & 5.40 & 2.97 & 0.00 & 55.31 & 11.48 & 3.24 & 6.10 \\
safety+math+code+medical & task\_arithmetic & 0.80 & 42.0 & 0.64 & 3.28 & 4.84 & 26.25 & 15.41 & 2.21 & 5.56 \\
safety+math+code+medical & task\_arithmetic & 0.80 & 43.0 & 3.09 & 4.22 & 6.41 & 20.62 & 24.01 & 2.23 & 5.37 \\
safety+math+code+medical & task\_arithmetic & 0.80 & 44.0 & 0.00 & 3.12 & 4.06 & 25.00 & 19.39 & 2.24 & 5.38 \\
safety+math+code+medical & task\_arithmetic & 1.00 & 42.0 & 0.00 & 0.94 & 8.12 & 84.38 & 32.24 & 2.31 & 7.26 \\
safety+math+code+medical & task\_arithmetic & 1.00 & 43.0 & 0.00 & 0.78 & 2.03 & 50.31 & 11.23 & 2.36 & 7.18 \\
safety+math+code+medical & task\_arithmetic & 1.00 & 44.0 & 0.00 & 0.00 & 2.50 & 39.69 & 11.37 & 2.36 & 7.00 \\
\bottomrule
\end{tabularx}
\end{table}


% --- Alpha スイープ 集計 - Method: task_arithmetic, Pattern: safety+math ---
\begin{table}[H]
\centering
\caption{Alpha スイープ 集計 - Method: task_arithmetic, Pattern: safety+math}
\label{tab:sweep_taskarithmetic_safety_math}
\small
\begin{tabular}{\textwidth}{lllRRRRRRR}
\toprule
\textbf{Pattern} & \textbf{Method} & \textbf{Alpha} & \textbf{Safety Ave [Harmful Content] (ASR$\downarrow$ \%)} & \textbf{Math Ave (\%)} & \textbf{Code Ave (\%)} & \textbf{Medical Ave (\%)} & \textbf{General/Inst Ave (\%)} & \textbf{EvolCode (PPL$\downarrow$)} & \textbf{MedAlpaca (PPL$\downarrow$)} \\
\midrule
safety+math & task\_arithmetic & 0.00 & 31.01 $\pm$ 0.00 & 6.25 $\pm$ 0.00 & 4.17 $\pm$ 1.80 & 63.33 $\pm$ 28.87 & 33.25 $\pm$ 0.16 & 2.03 $\pm$ 0.00 & 2.77 $\pm$ 0.00 \\
safety+math & task\_arithmetic & 0.20 & 30.94 $\pm$ 0.97 & 5.94 $\pm$ 0.54 & 6.46 $\pm$ 3.37 & 62.81 $\pm$ 28.97 & 31.67 $\pm$ 3.58 & 1.90 $\pm$ 0.01 & 2.67 $\pm$ 0.01 \\
safety+math & task\_arithmetic & 0.40 & 14.84 $\pm$ 1.16 & 5.78 $\pm$ 0.31 & 8.23 $\pm$ 3.46 & 75.52 $\pm$ 15.34 & 28.74 $\pm$ 3.37 & 1.83 $\pm$ 0.01 & 2.75 $\pm$ 0.01 \\
safety+math & task\_arithmetic & 0.60 & 2.48 $\pm$ 2.89 & 6.88 $\pm$ 1.56 & 7.50 $\pm$ 0.68 & 75.21 $\pm$ 15.93 & 27.18 $\pm$ 6.92 & 1.88 $\pm$ 0.02 & 3.47 $\pm$ 0.12 \\
safety+math & task\_arithmetic & 0.80 & 0.00 $\pm$ 0.00 & 9.11 $\pm$ 1.45 & 5.68 $\pm$ 0.24 & 68.70 $\pm$ 11.80 & 24.37 $\pm$ 9.95 & 2.03 $\pm$ 0.02 & 5.17 $\pm$ 0.07 \\
safety+math & task\_arithmetic & 1.00 & 0.00 $\pm$ 0.00 & 0.62 $\pm$ 0.54 & 6.93 $\pm$ 8.08 & 48.80 $\pm$ 8.77 & 18.36 $\pm$ 12.02 & 2.34 $\pm$ 0.03 & 7.15 $\pm$ 0.14 \\
\bottomrule
\end{tabularx}
\end{table}


% --- Alpha スイープ シード別詳細 - Method: task_arithmetic, Pattern: safety+math ---
\begin{table}[H]
\centering
\caption{Alpha スイープ シード別詳細 - Method: task_arithmetic, Pattern: safety+math}
\label{tab:sweep_seeds_taskarithmetic_safety_math}
\small
\begin{tabular}{\textwidth}{llllRRRRRRR}
\toprule
\textbf{Pattern} & \textbf{Method} & \textbf{Alpha} & \textbf{Seed} & \textbf{Safety Ave [Harmful Content] (ASR$\downarrow$ \%)} & \textbf{Math Ave (\%)} & \textbf{Code Ave (\%)} & \textbf{Medical Ave (\%)} & \textbf{General/Inst Ave (\%)} & \textbf{EvolCode (PPL$\downarrow$)} & \textbf{MedAlpaca (PPL$\downarrow$)} \\
\midrule
safety+math & task\_arithmetic & 0.00 & 42.0 & 31.01 & 6.25 & 6.25 & 30.00 & 33.16 & 2.03 & 2.77 \\
safety+math & task\_arithmetic & 0.00 & 43.0 & 31.01 & 6.25 & 3.12 & 80.00 & 33.16 & 2.03 & 2.77 \\
safety+math & task\_arithmetic & 0.00 & 44.0 & 31.01 & 6.25 & 3.12 & 80.00 & 33.43 & 2.03 & 2.77 \\
safety+math & task\_arithmetic & 0.20 & 42.0 & 31.14 & 5.62 & 10.31 & 29.38 & 27.54 & 1.89 & 2.66 \\
safety+math & task\_arithmetic & 0.20 & 43.0 & 31.80 & 6.56 & 5.00 & 78.75 & 33.42 & 1.90 & 2.68 \\
safety+math & task\_arithmetic & 0.20 & 44.0 & 29.88 & 5.62 & 4.06 & 80.31 & 34.04 & 1.90 & 2.68 \\
safety+math & task\_arithmetic & 0.40 & 42.0 & 13.98 & 5.47 & 12.19 & 57.81 & 28.70 & 1.82 & 2.76 \\
safety+math & task\_arithmetic & 0.40 & 43.0 & 14.38 & 6.09 & 6.72 & 84.06 & 25.38 & 1.83 & 2.74 \\
safety+math & task\_arithmetic & 0.40 & 44.0 & 16.15 & 5.78 & 5.78 & 84.69 & 32.13 & 1.85 & 2.76 \\
safety+math & task\_arithmetic & 0.60 & 42.0 & 0.16 & 8.44 & 6.72 & 56.88 & 35.05 & 1.86 & 3.48 \\
safety+math & task\_arithmetic & 0.60 & 43.0 & 1.57 & 5.31 & 7.97 & 83.12 & 22.05 & 1.89 & 3.34 \\
safety+math & task\_arithmetic & 0.60 & 44.0 & 5.72 & 6.88 & 7.81 & 85.62 & 24.45 & 1.89 & 3.57 \\
safety+math & task\_arithmetic & 0.80 & 42.0 & 0.00 & 9.53 & 5.94 & 57.03 & 14.19 & 2.01 & 5.13 \\
safety+math & task\_arithmetic & 0.80 & 43.0 & 0.00 & 10.31 & 5.47 & 68.44 & 34.07 & 2.05 & 5.25 \\
safety+math & task\_arithmetic & 0.80 & 44.0 & 0.00 & 7.50 & 5.62 & 80.62 & 24.85 & 2.04 & 5.13 \\
safety+math & task\_arithmetic & 1.00 & 42.0 & 0.00 & 0.94 & 16.25 & 56.72 & 32.24 & 2.31 & 7.26 \\
safety+math & task\_arithmetic & 1.00 & 43.0 & 0.00 & 0.94 & 2.03 & 50.31 & 11.23 & 2.36 & 7.18 \\
safety+math & task\_arithmetic & 1.00 & 44.0 & 0.00 & 0.00 & 2.50 & 39.38 & 11.62 & 2.36 & 7.00 \\
\bottomrule
\end{tabularx}
\end{table}


% --- Alpha スイープ 集計 - Method: task_arithmetic, Pattern: safety+code ---
\begin{table}[H]
\centering
\caption{Alpha スイープ 集計 - Method: task_arithmetic, Pattern: safety+code}
\label{tab:sweep_taskarithmetic_safety_code}
\small
\begin{tabular}{\textwidth}{lllRRRRRRR}
\toprule
\textbf{Pattern} & \textbf{Method} & \textbf{Alpha} & \textbf{Safety Ave [Harmful Content] (ASR$\downarrow$ \%)} & \textbf{Math Ave (\%)} & \textbf{Code Ave (\%)} & \textbf{Medical Ave (\%)} & \textbf{General/Inst Ave (\%)} & \textbf{EvolCode (PPL$\downarrow$)} & \textbf{MedAlpaca (PPL$\downarrow$)} \\
\midrule
safety+code & task\_arithmetic & 0.00 & 41.06 $\pm$ 0.33 & 5.62 $\pm$ 0.00 & 21.38 $\pm$ 0.00 & 86.25 $\pm$ 0.00 & 21.02 $\pm$ 0.06 & 1.44 $\pm$ 0.00 & 3.84 $\pm$ 0.00 \\
safety+code & task\_arithmetic & 0.20 & 39.04 $\pm$ 0.37 & 3.91 $\pm$ 0.31 & 0.78 $\pm$ 0.00 & 85.94 $\pm$ 0.00 & 34.76 $\pm$ 3.89 & 1.51 $\pm$ 0.00 & 3.66 $\pm$ 0.01 \\
safety+code & task\_arithmetic & 0.40 & 34.56 $\pm$ 1.54 & 2.45 $\pm$ 0.09 & 7.08 $\pm$ 3.58 & 86.88 $\pm$ 0.00 & 22.35 $\pm$ 6.37 & 1.91 $\pm$ 0.00 & 4.26 $\pm$ 0.01 \\
safety+code & task\_arithmetic & 0.60 & 29.17 $\pm$ 3.07 & 2.97 $\pm$ 0.68 & 4.63 $\pm$ 0.09 & 88.85 $\pm$ 0.79 & 17.39 $\pm$ 0.39 & 2.47 $\pm$ 0.01 & 4.84 $\pm$ 0.03 \\
safety+code & task\_arithmetic & 0.80 & 0.05 $\pm$ 0.09 & 2.03 $\pm$ 0.81 & 8.18 $\pm$ 0.86 & 79.79 $\pm$ 6.62 & 20.78 $\pm$ 5.96 & 2.11 $\pm$ 0.02 & 4.33 $\pm$ 0.06 \\
safety+code & task\_arithmetic & 1.00 & 0.00 $\pm$ 0.00 & 0.57 $\pm$ 0.50 & 4.22 $\pm$ 3.39 & 58.02 $\pm$ 23.47 & 18.37 $\pm$ 12.03 & 2.34 $\pm$ 0.03 & 7.15 $\pm$ 0.13 \\
\bottomrule
\end{tabularx}
\end{table}


% --- Alpha スイープ シード別詳細 - Method: task_arithmetic, Pattern: safety+code ---
\begin{table}[H]
\centering
\caption{Alpha スイープ シード別詳細 - Method: task_arithmetic, Pattern: safety+code}
\label{tab:sweep_seeds_taskarithmetic_safety_code}
\small
\begin{tabular}{\textwidth}{llllRRRRRRR}
\toprule
\textbf{Pattern} & \textbf{Method} & \textbf{Alpha} & \textbf{Seed} & \textbf{Safety Ave [Harmful Content] (ASR$\downarrow$ \%)} & \textbf{Math Ave (\%)} & \textbf{Code Ave (\%)} & \textbf{Medical Ave (\%)} & \textbf{General/Inst Ave (\%)} & \textbf{EvolCode (PPL$\downarrow$)} & \textbf{MedAlpaca (PPL$\downarrow$)} \\
\midrule
safety+code & task\_arithmetic & 0.00 & 42.0 & 40.68 & 5.62 & 21.38 & 86.25 & 20.98 & 1.44 & 3.84 \\
safety+code & task\_arithmetic & 0.00 & 43.0 & 41.25 & 5.62 & 21.38 & 86.25 & 20.98 & 1.44 & 3.84 \\
safety+code & task\_arithmetic & 0.00 & 44.0 & 41.25 & 5.62 & 21.38 & 86.25 & 21.08 & 1.44 & 3.84 \\
safety+code & task\_arithmetic & 0.20 & 42.0 & 38.75 & 3.59 & 0.78 & 85.94 & 39.21 & 1.51 & 3.66 \\
safety+code & task\_arithmetic & 0.20 & 43.0 & 39.46 & 3.91 & 0.78 & 85.94 & 32.01 & 1.51 & 3.67 \\
safety+code & task\_arithmetic & 0.20 & 44.0 & 38.91 & 4.22 & 0.78 & 85.94 & 33.06 & 1.51 & 3.65 \\
safety+code & task\_arithmetic & 0.40 & 42.0 & 35.62 & 2.50 & 9.53 & 86.88 & 19.59 & 1.90 & 4.25 \\
safety+code & task\_arithmetic & 0.40 & 43.0 & 35.26 & 2.34 & 2.97 & 86.88 & 29.63 & 1.91 & 4.27 \\
safety+code & task\_arithmetic & 0.40 & 44.0 & 32.80 & 2.50 & 8.75 & 86.88 & 17.82 & 1.91 & 4.26 \\
safety+code & task\_arithmetic & 0.60 & 42.0 & 32.56 & 3.28 & 4.53 & 89.69 & 17.60 & 2.47 & 4.80 \\
safety+code & task\_arithmetic & 0.60 & 43.0 & 26.56 & 3.44 & 4.68 & 88.12 & 17.62 & 2.47 & 4.84 \\
safety+code & task\_arithmetic & 0.60 & 44.0 & 28.39 & 2.19 & 4.68 & 88.75 & 16.93 & 2.48 & 4.87 \\
safety+code & task\_arithmetic & 0.80 & 42.0 & 0.00 & 2.50 & 8.59 & 86.88 & 27.66 & 2.09 & 4.36 \\
safety+code & task\_arithmetic & 0.80 & 43.0 & 0.16 & 2.50 & 7.19 & 78.75 & 17.43 & 2.11 & 4.25 \\
safety+code & task\_arithmetic & 0.80 & 44.0 & 0.00 & 1.09 & 8.75 & 73.75 & 17.24 & 2.13 & 4.36 \\
safety+code & task\_arithmetic & 1.00 & 42.0 & 0.00 & 0.94 & 8.12 & 84.38 & 32.26 & 2.31 & 7.26 \\
safety+code & task\_arithmetic & 1.00 & 43.0 & 0.00 & 0.78 & 2.03 & 50.31 & 11.23 & 2.36 & 7.18 \\
safety+code & task\_arithmetic & 1.00 & 44.0 & 0.00 & 0.00 & 2.50 & 39.38 & 11.62 & 2.36 & 7.00 \\
\bottomrule
\end{tabularx}
\end{table}


% --- Alpha スイープ 集計 - Method: task_arithmetic, Pattern: safety+medical ---
\begin{table}[H]
\centering
\caption{Alpha スイープ 集計 - Method: task_arithmetic, Pattern: safety+medical}
\label{tab:sweep_taskarithmetic_safety_medical}
\small
\begin{tabular}{\textwidth}{lllRRRRRRR}
\toprule
\textbf{Pattern} & \textbf{Method} & \textbf{Alpha} & \textbf{Safety Ave [Harmful Content] (ASR$\downarrow$ \%)} & \textbf{Math Ave (\%)} & \textbf{Code Ave (\%)} & \textbf{Medical Ave (\%)} & \textbf{General/Inst Ave (\%)} & \textbf{EvolCode (PPL$\downarrow$)} & \textbf{MedAlpaca (PPL$\downarrow$)} \\
\midrule
safety+medical & task\_arithmetic & 0.00 & 6.48 $\pm$ 0.00 & 4.06 $\pm$ 0.00 & 1.72 $\pm$ 0.00 & 68.44 $\pm$ 0.00 & 19.09 $\pm$ 0.09 & 1.89 $\pm$ 0.00 & 1.43 $\pm$ 0.00 \\
safety+medical & task\_arithmetic & 0.20 & 0.08 $\pm$ 0.00 & 1.77 $\pm$ 0.09 & 0.00 $\pm$ 0.00 & 16.04 $\pm$ 0.18 & 12.49 $\pm$ 0.11 & 54.04 $\pm$ 0.07 & 25.78 $\pm$ 0.20 \\
safety+medical & task\_arithmetic & 0.40 & 0.00 $\pm$ 0.00 & 0.00 $\pm$ 0.00 & 0.00 $\pm$ 0.00 & 20.00 $\pm$ 7.04 & 12.25 $\pm$ 0.07 & 36901.81 $\pm$ 950.77 & 37492.79 $\pm$ 653.76 \\
safety+medical & task\_arithmetic & 0.60 & 0.16 $\pm$ 0.00 & 0.00 $\pm$ 0.00 & 0.00 $\pm$ 0.00 & 13.85 $\pm$ 0.18 & 12.72 $\pm$ 0.06 & 217753.69 $\pm$ 12313.53 & 198661.81 $\pm$ 12470.98 \\
safety+medical & task\_arithmetic & 0.80 & 0.00 $\pm$ 0.00 & 0.62 $\pm$ 0.41 & 0.00 $\pm$ 0.00 & 38.44 $\pm$ 15.25 & 12.04 $\pm$ 0.26 & 577.39 $\pm$ 106.94 & 257.99 $\pm$ 44.44 \\
safety+medical & task\_arithmetic & 1.00 & 0.00 $\pm$ 0.00 & 0.57 $\pm$ 0.50 & 4.22 $\pm$ 3.39 & 58.02 $\pm$ 23.47 & 18.37 $\pm$ 12.03 & 2.34 $\pm$ 0.03 & 7.15 $\pm$ 0.13 \\
\bottomrule
\end{tabularx}
\end{table}


% --- Alpha スイープ シード別詳細 - Method: task_arithmetic, Pattern: safety+medical ---
\begin{table}[H]
\centering
\caption{Alpha スイープ シード別詳細 - Method: task_arithmetic, Pattern: safety+medical}
\label{tab:sweep_seeds_taskarithmetic_safety_medical}
\small
\begin{tabular}{\textwidth}{llllRRRRRRR}
\toprule
\textbf{Pattern} & \textbf{Method} & \textbf{Alpha} & \textbf{Seed} & \textbf{Safety Ave [Harmful Content] (ASR$\downarrow$ \%)} & \textbf{Math Ave (\%)} & \textbf{Code Ave (\%)} & \textbf{Medical Ave (\%)} & \textbf{General/Inst Ave (\%)} & \textbf{EvolCode (PPL$\downarrow$)} & \textbf{MedAlpaca (PPL$\downarrow$)} \\
\midrule
safety+medical & task\_arithmetic & 0.00 & 42.0 & 6.48 & 4.06 & 1.72 & 68.44 & 19.00 & 1.89 & 1.43 \\
safety+medical & task\_arithmetic & 0.00 & 43.0 & 6.48 & 4.06 & 1.72 & 68.44 & 19.10 & 1.89 & 1.43 \\
safety+medical & task\_arithmetic & 0.00 & 44.0 & 6.48 & 4.06 & 1.72 & 68.44 & 19.17 & 1.89 & 1.43 \\
safety+medical & task\_arithmetic & 0.20 & 42.0 & 0.08 & 1.72 & 0.00 & 16.25 & 12.39 & 54.12 & 25.57 \\
safety+medical & task\_arithmetic & 0.20 & 43.0 & 0.08 & 1.88 & 0.00 & 15.94 & 12.49 & 53.99 & 25.97 \\
safety+medical & task\_arithmetic & 0.20 & 44.0 & 0.08 & 1.72 & 0.00 & 15.94 & 12.60 & 54.02 & 25.79 \\
safety+medical & task\_arithmetic & 0.40 & 42.0 & 0.00 & 0.00 & 0.00 & 28.12 & 12.34 & 35846.12 & 36843.25 \\
safety+medical & task\_arithmetic & 0.40 & 43.0 & 0.00 & 0.00 & 0.00 & 15.62 & 12.20 & 37168.66 & 37484.42 \\
safety+medical & task\_arithmetic & 0.40 & 44.0 & 0.00 & 0.00 & 0.00 & 16.25 & 12.22 & 37690.63 & 38150.69 \\
safety+medical & task\_arithmetic & 0.60 & 42.0 & 0.16 & 0.00 & 0.00 & 13.75 & 12.69 & 231670.29 & 211313.97 \\
safety+medical & task\_arithmetic & 0.60 & 43.0 & 0.16 & 0.00 & 0.00 & 14.06 & 12.69 & 213319.17 & 186380.28 \\
safety+medical & task\_arithmetic & 0.60 & 44.0 & 0.16 & 0.00 & 0.00 & 13.75 & 12.79 & 208271.62 & 198291.16 \\
safety+medical & task\_arithmetic & 0.80 & 42.0 & 0.00 & 0.47 & 0.00 & 51.25 & 11.80 & 520.77 & 209.85 \\
safety+medical & task\_arithmetic & 0.80 & 43.0 & 0.00 & 1.09 & 0.00 & 21.56 & 12.32 & 700.74 & 266.69 \\
safety+medical & task\_arithmetic & 0.80 & 44.0 & 0.00 & 0.31 & 0.00 & 42.50 & 12.01 & 510.67 & 297.44 \\
safety+medical & task\_arithmetic & 1.00 & 42.0 & 0.00 & 0.94 & 8.12 & 84.38 & 32.26 & 2.31 & 7.26 \\
safety+medical & task\_arithmetic & 1.00 & 43.0 & 0.00 & 0.78 & 2.03 & 50.31 & 11.23 & 2.36 & 7.18 \\
safety+medical & task\_arithmetic & 1.00 & 44.0 & 0.00 & 0.00 & 2.50 & 39.38 & 11.62 & 2.36 & 7.00 \\
\bottomrule
\end{tabularx}
\end{table}


% --- Alpha スイープ 集計 - Method: safemerge, Pattern: safety+math ---
\begin{table}[H]
\centering
\caption{Alpha スイープ 集計 - Method: safemerge, Pattern: safety+math}
\label{tab:sweep_safemerge_safety_math}
\small
\begin{tabular}{\textwidth}{lllRRRRRRR}
\toprule
\textbf{Pattern} & \textbf{Method} & \textbf{Alpha} & \textbf{Safety Ave [Harmful Content] (ASR$\downarrow$ \%)} & \textbf{Math Ave (\%)} & \textbf{Code Ave (\%)} & \textbf{Medical Ave (\%)} & \textbf{General/Inst Ave (\%)} & \textbf{EvolCode (PPL$\downarrow$)} & \textbf{MedAlpaca (PPL$\downarrow$)} \\
\midrule
safety+math & safemerge & N/A & 4.61 $\pm$ 7.99 & 1.30 $\pm$ 0.86 & 6.67 $\pm$ 4.86 & 82.50 $\pm$ 1.74 & 13.09 $\pm$ 1.77 & 2.22 $\pm$ 0.48 & 6.41 $\pm$ 3.54 \\
\bottomrule
\end{tabularx}
\end{table}


% --- Alpha スイープ シード別詳細 - Method: safemerge, Pattern: safety+math ---
\begin{table}[H]
\centering
\caption{Alpha スイープ シード別詳細 - Method: safemerge, Pattern: safety+math}
\label{tab:sweep_seeds_safemerge_safety_math}
\small
\begin{tabular}{\textwidth}{llllRRRRRRR}
\toprule
\textbf{Pattern} & \textbf{Method} & \textbf{Alpha} & \textbf{Seed} & \textbf{Safety Ave [Harmful Content] (ASR$\downarrow$ \%)} & \textbf{Math Ave (\%)} & \textbf{Code Ave (\%)} & \textbf{Medical Ave (\%)} & \textbf{General/Inst Ave (\%)} & \textbf{EvolCode (PPL$\downarrow$)} & \textbf{MedAlpaca (PPL$\downarrow$)} \\
\midrule
safety+math & safemerge & N/A & 42.0 & 13.84 & 2.19 & 8.91 & 80.62 & 15.06 & 1.71 & 2.51 \\
safety+math & safemerge & N/A & 43.0 & 0.00 & 0.47 & 1.09 & 84.06 & 11.63 & 2.66 & 9.43 \\
safety+math & safemerge & N/A & 44.0 & 0.00 & 1.25 & 10.00 & 82.81 & 12.60 & 2.30 & 7.29 \\
\bottomrule
\end{tabularx}
\end{table}


% --- Alpha スイープ 集計 - Method: safemerge, Pattern: safety+code ---
\begin{table}[H]
\centering
\caption{Alpha スイープ 集計 - Method: safemerge, Pattern: safety+code}
\label{tab:sweep_safemerge_safety_code}
\small
\begin{tabular}{\textwidth}{lllRRRRRRR}
\toprule
\textbf{Pattern} & \textbf{Method} & \textbf{Alpha} & \textbf{Safety Ave [Harmful Content] (ASR$\downarrow$ \%)} & \textbf{Math Ave (\%)} & \textbf{Code Ave (\%)} & \textbf{Medical Ave (\%)} & \textbf{General/Inst Ave (\%)} & \textbf{EvolCode (PPL$\downarrow$)} & \textbf{MedAlpaca (PPL$\downarrow$)} \\
\midrule
safety+code & safemerge & N/A & 0.10 $\pm$ 0.18 & 1.51 $\pm$ 1.33 & 5.99 $\pm$ 5.39 & 73.75 $\pm$ 2.98 & 19.84 $\pm$ 8.02 & 2.35 $\pm$ 0.78 & 7.03 $\pm$ 4.96 \\
\bottomrule
\end{tabularx}
\end{table}


% --- Alpha スイープ シード別詳細 - Method: safemerge, Pattern: safety+code ---
\begin{table}[H]
\centering
\caption{Alpha スイープ シード別詳細 - Method: safemerge, Pattern: safety+code}
\label{tab:sweep_seeds_safemerge_safety_code}
\small
\begin{tabular}{\textwidth}{llllRRRRRRR}
\toprule
\textbf{Pattern} & \textbf{Method} & \textbf{Alpha} & \textbf{Seed} & \textbf{Safety Ave [Harmful Content] (ASR$\downarrow$ \%)} & \textbf{Math Ave (\%)} & \textbf{Code Ave (\%)} & \textbf{Medical Ave (\%)} & \textbf{General/Inst Ave (\%)} & \textbf{EvolCode (PPL$\downarrow$)} & \textbf{MedAlpaca (PPL$\downarrow$)} \\
\midrule
safety+code & safemerge & N/A & 42.0 & 0.31 & 2.50 & 7.03 & 75.62 & 27.76 & 1.75 & 2.87 \\
safety+code & safemerge & N/A & 43.0 & 0.00 & 2.03 & 10.78 & 70.31 & 20.03 & 2.06 & 5.72 \\
safety+code & safemerge & N/A & 44.0 & 0.00 & 0.00 & 0.16 & 75.31 & 11.72 & 3.23 & 12.51 \\
\bottomrule
\end{tabularx}
\end{table}


% --- Alpha スイープ 集計 - Method: safemerge, Pattern: safety+medical ---
\begin{table}[H]
\centering
\caption{Alpha スイープ 集計 - Method: safemerge, Pattern: safety+medical}
\label{tab:sweep_safemerge_safety_medical}
\small
\begin{tabular}{\textwidth}{lllRRRRRRR}
\toprule
\textbf{Pattern} & \textbf{Method} & \textbf{Alpha} & \textbf{Safety Ave [Harmful Content] (ASR$\downarrow$ \%)} & \textbf{Math Ave (\%)} & \textbf{Code Ave (\%)} & \textbf{Medical Ave (\%)} & \textbf{General/Inst Ave (\%)} & \textbf{EvolCode (PPL$\downarrow$)} & \textbf{MedAlpaca (PPL$\downarrow$)} \\
\midrule
safety+medical & safemerge & N/A & 0.03 $\pm$ 0.05 & 1.09 $\pm$ 1.89 & 3.91 $\pm$ 6.77 & 53.44 $\pm$ 10.84 & 12.81 $\pm$ 1.63 & 3.84 $\pm$ 1.84 & 15.66 $\pm$ 10.55 \\
\bottomrule
\end{tabularx}
\end{table}


% --- Alpha スイープ シード別詳細 - Method: safemerge, Pattern: safety+medical ---
\begin{table}[H]
\centering
\caption{Alpha スイープ シード別詳細 - Method: safemerge, Pattern: safety+medical}
\label{tab:sweep_seeds_safemerge_safety_medical}
\small
\begin{tabular}{\textwidth}{llllRRRRRRR}
\toprule
\textbf{Pattern} & \textbf{Method} & \textbf{Alpha} & \textbf{Seed} & \textbf{Safety Ave [Harmful Content] (ASR$\downarrow$ \%)} & \textbf{Math Ave (\%)} & \textbf{Code Ave (\%)} & \textbf{Medical Ave (\%)} & \textbf{General/Inst Ave (\%)} & \textbf{EvolCode (PPL$\downarrow$)} & \textbf{MedAlpaca (PPL$\downarrow$)} \\
\midrule
safety+medical & safemerge & N/A & 42.0 & 0.00 & 3.28 & 11.72 & 60.31 & 14.69 & 1.87 & 4.37 \\
safety+medical & safemerge & N/A & 43.0 & 0.08 & 0.00 & 0.00 & 40.94 & 11.81 & 5.53 & 25.25 \\
safety+medical & safemerge & N/A & 44.0 & 0.00 & 0.00 & 0.00 & 59.06 & 11.93 & 4.11 & 17.37 \\
\bottomrule
\end{tabularx}
\end{table}


% --- Alpha スイープ 集計 - Method: led_merging, Pattern: safety+math ---
\begin{table}[H]
\centering
\caption{Alpha スイープ 集計 - Method: led_merging, Pattern: safety+math}
\label{tab:sweep_ledmerging_safety_math}
\small
\begin{tabular}{\textwidth}{lllRRRRRRR}
\toprule
\textbf{Pattern} & \textbf{Method} & \textbf{Alpha} & \textbf{Safety Ave [Harmful Content] (ASR$\downarrow$ \%)} & \textbf{Math Ave (\%)} & \textbf{Code Ave (\%)} & \textbf{Medical Ave (\%)} & \textbf{General/Inst Ave (\%)} & \textbf{EvolCode (PPL$\downarrow$)} & \textbf{MedAlpaca (PPL$\downarrow$)} \\
\midrule
safety+math & led\_merging & N/A & 2.80 $\pm$ 0.00 & 1.41 $\pm$ 0.00 & 5.62 $\pm$ 0.00 & 83.75 $\pm$ 0.00 & 22.51 $\pm$ 0.02 & 1.75 $\pm$ 0.00 & 2.72 $\pm$ 0.00 \\
\bottomrule
\end{tabularx}
\end{table}


% --- Alpha スイープ シード別詳細 - Method: led_merging, Pattern: safety+math ---
\begin{table}[H]
\centering
\caption{Alpha スイープ シード別詳細 - Method: led_merging, Pattern: safety+math}
\label{tab:sweep_seeds_ledmerging_safety_math}
\small
\begin{tabular}{\textwidth}{llllRRRRRRR}
\toprule
\textbf{Pattern} & \textbf{Method} & \textbf{Alpha} & \textbf{Seed} & \textbf{Safety Ave [Harmful Content] (ASR$\downarrow$ \%)} & \textbf{Math Ave (\%)} & \textbf{Code Ave (\%)} & \textbf{Medical Ave (\%)} & \textbf{General/Inst Ave (\%)} & \textbf{EvolCode (PPL$\downarrow$)} & \textbf{MedAlpaca (PPL$\downarrow$)} \\
\midrule
safety+math & led\_merging & N/A & 42.0 & 2.80 & 1.41 & 5.62 & 83.75 & 22.50 & 1.75 & 2.72 \\
safety+math & led\_merging & N/A & 43.0 & 2.80 & 1.41 & 5.62 & 83.75 & 22.52 & 1.75 & 2.72 \\
safety+math & led\_merging & N/A & 44.0 & 2.80 & 1.41 & 5.62 & 83.75 & 22.49 & 1.75 & 2.72 \\
\bottomrule
\end{tabularx}
\end{table}


% --- Alpha スイープ 集計 - Method: led_merging, Pattern: safety+code ---
\begin{table}[H]
\centering
\caption{Alpha スイープ 集計 - Method: led_merging, Pattern: safety+code}
\label{tab:sweep_ledmerging_safety_code}
\small
\begin{tabular}{\textwidth}{lllRRRRRRR}
\toprule
\textbf{Pattern} & \textbf{Method} & \textbf{Alpha} & \textbf{Safety Ave [Harmful Content] (ASR$\downarrow$ \%)} & \textbf{Math Ave (\%)} & \textbf{Code Ave (\%)} & \textbf{Medical Ave (\%)} & \textbf{General/Inst Ave (\%)} & \textbf{EvolCode (PPL$\downarrow$)} & \textbf{MedAlpaca (PPL$\downarrow$)} \\
\midrule
safety+code & led\_merging & N/A & 2.88 $\pm$ 0.14 & 1.41 $\pm$ 0.00 & 5.62 $\pm$ 0.00 & 83.75 $\pm$ 0.00 & 22.49 $\pm$ 0.04 & 1.75 $\pm$ 0.00 & 2.72 $\pm$ 0.00 \\
\bottomrule
\end{tabularx}
\end{table}


% --- Alpha スイープ シード別詳細 - Method: led_merging, Pattern: safety+code ---
\begin{table}[H]
\centering
\caption{Alpha スイープ シード別詳細 - Method: led_merging, Pattern: safety+code}
\label{tab:sweep_seeds_ledmerging_safety_code}
\small
\begin{tabular}{\textwidth}{llllRRRRRRR}
\toprule
\textbf{Pattern} & \textbf{Method} & \textbf{Alpha} & \textbf{Seed} & \textbf{Safety Ave [Harmful Content] (ASR$\downarrow$ \%)} & \textbf{Math Ave (\%)} & \textbf{Code Ave (\%)} & \textbf{Medical Ave (\%)} & \textbf{General/Inst Ave (\%)} & \textbf{EvolCode (PPL$\downarrow$)} & \textbf{MedAlpaca (PPL$\downarrow$)} \\
\midrule
safety+code & led\_merging & N/A & 42.0 & 3.04 & 1.41 & 5.62 & 83.75 & 22.49 & 1.75 & 2.72 \\
safety+code & led\_merging & N/A & 43.0 & 2.80 & 1.41 & 5.62 & 83.75 & 22.52 & 1.75 & 2.72 \\
safety+code & led\_merging & N/A & 44.0 & 2.80 & 1.41 & 5.62 & 83.75 & 22.45 & 1.75 & 2.72 \\
\bottomrule
\end{tabularx}
\end{table}


% --- Alpha スイープ 集計 - Method: mergealign, Pattern: safety+math ---
\begin{table}[H]
\centering
\caption{Alpha スイープ 集計 - Method: mergealign, Pattern: safety+math}
\label{tab:sweep_mergealign_safety_math}
\small
\begin{tabular}{\textwidth}{lllRRRRRRR}
\toprule
\textbf{Pattern} & \textbf{Method} & \textbf{Alpha} & \textbf{Safety Ave [Harmful Content] (ASR$\downarrow$ \%)} & \textbf{Math Ave (\%)} & \textbf{Code Ave (\%)} & \textbf{Medical Ave (\%)} & \textbf{General/Inst Ave (\%)} & \textbf{EvolCode (PPL$\downarrow$)} & \textbf{MedAlpaca (PPL$\downarrow$)} \\
\midrule
safety+math & mergealign & N/A & 0.00 $\pm$ 0.00 & 0.00 $\pm$ 0.00 & 0.00 $\pm$ 0.00 & 86.25 $\pm$ 0.00 & 15.73 $\pm$ 3.41 & N/A & N/A \\
\bottomrule
\end{tabularx}
\end{table}


% --- Alpha スイープ シード別詳細 - Method: mergealign, Pattern: safety+math ---
\begin{table}[H]
\centering
\caption{Alpha スイープ シード別詳細 - Method: mergealign, Pattern: safety+math}
\label{tab:sweep_seeds_mergealign_safety_math}
\small
\begin{tabular}{\textwidth}{llllRRRRRRR}
\toprule
\textbf{Pattern} & \textbf{Method} & \textbf{Alpha} & \textbf{Seed} & \textbf{Safety Ave [Harmful Content] (ASR$\downarrow$ \%)} & \textbf{Math Ave (\%)} & \textbf{Code Ave (\%)} & \textbf{Medical Ave (\%)} & \textbf{General/Inst Ave (\%)} & \textbf{EvolCode (PPL$\downarrow$)} & \textbf{MedAlpaca (PPL$\downarrow$)} \\
\midrule
safety+math & mergealign & N/A & 42.0 & 0.00 & 0.00 & 0.00 & 86.25 & 17.71 & N/A & N/A \\
safety+math & mergealign & N/A & 43.0 & 0.00 & 0.00 & 0.00 & 86.25 & 11.79 & N/A & N/A \\
safety+math & mergealign & N/A & 44.0 & 0.00 & 0.00 & 0.00 & 86.25 & 17.71 & N/A & N/A \\
\bottomrule
\end{tabularx}
\end{table}


% --- Alpha スイープ 集計 - Method: mergealign, Pattern: safety+code ---
\begin{table}[H]
\centering
\caption{Alpha スイープ 集計 - Method: mergealign, Pattern: safety+code}
\label{tab:sweep_mergealign_safety_code}
\small
\begin{tabular}{\textwidth}{lllRRRRRRR}
\toprule
\textbf{Pattern} & \textbf{Method} & \textbf{Alpha} & \textbf{Safety Ave [Harmful Content] (ASR$\downarrow$ \%)} & \textbf{Math Ave (\%)} & \textbf{Code Ave (\%)} & \textbf{Medical Ave (\%)} & \textbf{General/Inst Ave (\%)} & \textbf{EvolCode (PPL$\downarrow$)} & \textbf{MedAlpaca (PPL$\downarrow$)} \\
\midrule
safety+code & mergealign & N/A & 0.00 $\pm$ 0.00 & 0.00 $\pm$ 0.00 & 0.00 $\pm$ 0.00 & 86.25 $\pm$ 0.00 & 15.73 $\pm$ 3.41 & N/A & N/A \\
\bottomrule
\end{tabularx}
\end{table}


% --- Alpha スイープ シード別詳細 - Method: mergealign, Pattern: safety+code ---
\begin{table}[H]
\centering
\caption{Alpha スイープ シード別詳細 - Method: mergealign, Pattern: safety+code}
\label{tab:sweep_seeds_mergealign_safety_code}
\small
\begin{tabular}{\textwidth}{llllRRRRRRR}
\toprule
\textbf{Pattern} & \textbf{Method} & \textbf{Alpha} & \textbf{Seed} & \textbf{Safety Ave [Harmful Content] (ASR$\downarrow$ \%)} & \textbf{Math Ave (\%)} & \textbf{Code Ave (\%)} & \textbf{Medical Ave (\%)} & \textbf{General/Inst Ave (\%)} & \textbf{EvolCode (PPL$\downarrow$)} & \textbf{MedAlpaca (PPL$\downarrow$)} \\
\midrule
safety+code & mergealign & N/A & 42.0 & 0.00 & 0.00 & 0.00 & 86.25 & 17.71 & N/A & N/A \\
safety+code & mergealign & N/A & 43.0 & 0.00 & 0.00 & 0.00 & 86.25 & 11.79 & N/A & N/A \\
safety+code & mergealign & N/A & 44.0 & 0.00 & 0.00 & 0.00 & 86.25 & 17.71 & N/A & N/A \\
\bottomrule
\end{tabularx}
\end{table}


% --- Alpha スイープ 集計 - Method: mergealign, Pattern: safety+medical ---
\begin{table}[H]
\centering
\caption{Alpha スイープ 集計 - Method: mergealign, Pattern: safety+medical}
\label{tab:sweep_mergealign_safety_medical}
\small
\begin{tabular}{\textwidth}{lllRRRRRRR}
\toprule
\textbf{Pattern} & \textbf{Method} & \textbf{Alpha} & \textbf{Safety Ave [Harmful Content] (ASR$\downarrow$ \%)} & \textbf{Math Ave (\%)} & \textbf{Code Ave (\%)} & \textbf{Medical Ave (\%)} & \textbf{General/Inst Ave (\%)} & \textbf{EvolCode (PPL$\downarrow$)} & \textbf{MedAlpaca (PPL$\downarrow$)} \\
\midrule
safety+medical & mergealign & N/A & 0.00 $\pm$ 0.00 & 0.00 $\pm$ 0.00 & 0.00 $\pm$ 0.00 & 86.25 $\pm$ 0.00 & 13.76 $\pm$ 3.41 & N/A & N/A \\
\bottomrule
\end{tabularx}
\end{table}


% --- Alpha スイープ シード別詳細 - Method: mergealign, Pattern: safety+medical ---
\begin{table}[H]
\centering
\caption{Alpha スイープ シード別詳細 - Method: mergealign, Pattern: safety+medical}
\label{tab:sweep_seeds_mergealign_safety_medical}
\small
\begin{tabular}{\textwidth}{llllRRRRRRR}
\toprule
\textbf{Pattern} & \textbf{Method} & \textbf{Alpha} & \textbf{Seed} & \textbf{Safety Ave [Harmful Content] (ASR$\downarrow$ \%)} & \textbf{Math Ave (\%)} & \textbf{Code Ave (\%)} & \textbf{Medical Ave (\%)} & \textbf{General/Inst Ave (\%)} & \textbf{EvolCode (PPL$\downarrow$)} & \textbf{MedAlpaca (PPL$\downarrow$)} \\
\midrule
safety+medical & mergealign & N/A & 42.0 & 0.00 & 0.00 & 0.00 & 86.25 & 17.71 & N/A & N/A \\
safety+medical & mergealign & N/A & 43.0 & 0.00 & 0.00 & 0.00 & 86.25 & 11.79 & N/A & N/A \\
safety+medical & mergealign & N/A & 44.0 & 0.00 & 0.00 & 0.00 & 86.25 & 11.79 & N/A & N/A \\
\bottomrule
\end{tabularx}
\end{table}

